% Leçon 12 — Vocabulaire et compréhension de texte : paludisme — Chinois Tle A — Cours signé OnBuch+
% Programme officiel MINESEC. Compiler avec : tectonic ZH12-vocabulaire-paludisme.tex
\newcommand{\DOCMATIERE}{Chinois}
\newcommand{\DOCNIVEAU}{Tle A}
\newcommand{\DOCMODULE}{Module 2 : Environnement et santé}
\newcommand{\DOCLECON}{ZH12}
\newcommand{\DOCTITRE}{Leçon 12 — Vocabulaire et compréhension de texte : paludisme}
% OnBuch+ — préambule « cours signés » ENRICHI (compilé avec tectonic / XeLaTeX)
% Système visuel « Pop » d'OnBuch+ : fond crème (jamais blanc), bordures encre
% épaisses, ombres « dures » décalées, titres Archivo Black en capitales.
% Version MATHÉMATIQUES : graphes (pgfplots/tikz), boîtes pédagogiques
% (définition, propriété, méthode, attention, le savais-tu, exercice résolu,
% exercice, TP, bilan), page de garde et signature OnBuch+ ancrée en bas.
%
% Macros attendues dans main.tex AVANT \input{preamble} :
%   \DOCMATIERE  \DOCNIVEAU  \DOCMODULE  \DOCLECON (numéro)  \DOCTITRE
\documentclass[11pt]{article}

\usepackage{fontspec}
\usepackage[french]{babel} % français = langue principale ; le chinois via \zh{..} / textebox (xeCJK)
\usepackage{xeCJK}
\usepackage{pifont} % \ding{51} \ding{55} utilisés par les modèles
\usepackage[a4paper,margin=2cm,top=3.4cm,bottom=2.8cm,headheight=1.4cm,footskip=1.3cm]{geometry}
\usepackage{amsmath,amssymb,amsfonts}
\usepackage{mathtools} % \xmapsto, \coloneqq... souvent utilisés par les modèles
\usepackage{cancel} % simplifications barrées (bilans de forces)
% \abs{z} (module, valeur absolue) et \norm{u} (norme), utilisés par les modèles
\providecommand{\abs}{}\let\abs\relax\DeclarePairedDelimiter{\abs}{\lvert}{\rvert}
\providecommand{\norm}{}\let\norm\relax\DeclarePairedDelimiter{\norm}{\lVert}{\rVert}
\usepackage[table]{xcolor} % [table] : fournit aussi \rowcolors (alternance de lignes)
\usepackage{pagecolor}
\usepackage{tikz}
\usetikzlibrary{arrows.meta,positioning,calc,shapes.geometric,shapes.misc,shapes.arrows,decorations.pathmorphing,decorations.pathreplacing,decorations.markings,patterns,shadows,mindmap,trees,fit,backgrounds,matrix,intersections,angles,quotes}
\usepackage{pgfplots}
\usepackage[version=4]{mhchem} % équations nucléaires \ce{^{235}_{92}U}
\usepackage[european,siunitx]{circuitikz} % schémas électriques
\pgfplotsset{compat=1.18}
\usepgfplotslibrary{fillbetween,groupplots} % aires entre courbes (calcul intégral)
\usepackage{enumitem}
\usepackage{fancyhdr}
% [most] charge toutes les bibliothèques tcolorbox (listings, minted, raster,
% documentation...) et gonfle inutilement la mémoire TeX sur un document long
% (~300 boîtes) : on ne charge que ce qui est réellement utilisé.
\usepackage[skins,breakable]{tcolorbox}
\usepackage{graphicx}
\usepackage{colortbl}
\usepackage{booktabs}
\usepackage{tabularx}
\usepackage{diagbox} % case d'en-tête diagonale des tableaux à double entrée
% Colonnes extensibles centrée / gauche / droite (C, L, R), souvent utilisées par les modèles
\newcolumntype{C}{>{\centering\arraybackslash}X}
\newcolumntype{L}{>{\raggedright\arraybackslash}X}
\newcolumntype{R}{>{\raggedleft\arraybackslash}X}
\usepackage{array}
\usepackage{multicol}
\usepackage{multirow}
\usepackage{siunitx}
% parse-numbers=false : accepte aussi \SI{2,5 \times 10^{-3}}{...}
\sisetup{locale=FR,output-decimal-marker={,},group-minimum-digits=4,parse-numbers=false}
\DeclareSIUnit\ohm{\ensuremath{\symup{\Omega}}} % U+2126 absent de la police
\DeclareSIUnit\division{div} % réglages d'oscilloscope (ms/div, V/div)
\DeclareSIUnit\tr{tr} % tours (vitesse de rotation, ex. \SI{1500}{\tr\per\minute})
\DeclareSIUnit\molar{M} % concentration molaire (ex. \SI{3}{\molar}, \SI{1}{\micro\molar})
\DeclareSIUnit\an{an} % année (le modèle confond parfois avec \year, un compteur TeX) — ex. \SI{5}{\kilo\meter\per\an}
\DeclareSIUnit\FCFA{FCFA} % franc CFA (ex. \SI{500}{\FCFA\per\kilo\gram})
\DeclareSIUnit\degreeBrix{\textdegree Brix} % teneur en sucre d'un sirop/jus (réfractométrie)
\DeclareSIUnit\month{mois} % le modèle utilise parfois \month, un compteur TeX, comme unité de durée
\usepackage{qrcode}
% \degree : commande fréquemment utilisée par les modèles pour un angle ou une
% température en dehors de \SI{}{} (non fournie par les paquets chargés).
\providecommand{\degree}{^{\circ}}
% \qed (fin de démonstration), utilisé par les modèles sans amsthm
\providecommand{\qed}{\hfill$\square$}
% \barre : double barre des valeurs interdites dans un tableau de variations (tkz-tab)
\providecommand{\barre}{\ensuremath{\|}}
% environnement proof (amsthm non chargé) utilisé par les modèles
\ifdefined\proof\else\newenvironment{proof}[1][Démonstration]{\par\noindent\textit{#1.}\ }{\qed\par}\fi
\usepackage{lastpage}
\usepackage{hyperref}
\hypersetup{hidelinks}

% ============================================================
% Palette « Pop » OnBuch+ — mêmes hex que la classe P (plus_ui.dart)
% ============================================================
\definecolor{poporange}{HTML}{F59321}
\definecolor{popdark}{HTML}{E07A0C}
\definecolor{popcream}{HTML}{FFF6E8}
\definecolor{popink}{HTML}{1C1714}
\definecolor{poppurple}{HTML}{7A5AE0}
\definecolor{popgreen}{HTML}{2E9E6A}
\definecolor{poppink}{HTML}{E0457B}
\definecolor{popblue}{HTML}{2F7DE1}
\definecolor{popgold}{HTML}{C98A12}
\definecolor{popmuted}{HTML}{8A7F76}
\definecolor{popline}{HTML}{EADFD2}
\definecolor{popgray}{HTML}{8A7F76}
\colorlet{popgrey}{popgray}
% Alias tolérants (le modèle nomme parfois les couleurs différemment)
\colorlet{popwhite}{white}
\colorlet{popblack}{popink}
\colorlet{popred}{poppink}
\colorlet{popyellow}{popgold}
% Teintes claires (fonds de tableaux/encadrés) — le modèle nomme parfois
% les couleurs avec un suffixe L (ex. popblueL) sans les avoir définies.
\colorlet{poporangeL}{poporange!20}
\colorlet{popdarkL}{popdark!20}
\colorlet{poppurpleL}{poppurple!20}
\colorlet{popgreenL}{popgreen!20}
\colorlet{poppinkL}{poppink!20}
\colorlet{popblueL}{popblue!20}
\colorlet{popgoldL}{popgold!20}
\colorlet{popredL}{popred!20}
\colorlet{popgrayL}{popgray!20}
% Teintes claires pour fonds de tableaux / zones
\colorlet{popblueL}{popblue!10!white}
\colorlet{poporangeL}{poporange!14!white}
\colorlet{popgreenL}{popgreen!12!white}
\colorlet{poppurpleL}{poppurple!12!white}
\colorlet{poppinkL}{poppink!12!white}
\colorlet{popgoldL}{popgold!14!white}

\pagecolor{popcream}

% ============================================================
% Typographie
% ============================================================
\defaultfontfeatures{Ligatures=TeX}
\setmainfont{PlusJakartaSans-Medium.ttf}[Path=../../../fonts/,BoldFont=PlusJakartaSans-Bold.ttf]
% Symboles, API et emojis absents de la police principale : repli sur DejaVu Sans (caractères actifs)
\newfontface\symfont{DejaVuSans.ttf}[Path=../../../fonts/]
\makeatletter
\newcommand\symactive[1]{\catcode#1=13 \begingroup\lccode`\~=#1 \lowercase{\endgroup\def~}{{\symfont\char#1}}}
\newcommand\symblank[1]{\catcode#1=13 \begingroup\lccode`\~=#1 \lowercase{\endgroup\def~}{}}
\newcommand\symhyph[1]{\catcode#1=13 \begingroup\lccode`\~=#1 \lowercase{\endgroup\def~}{\mbox{-}}}
\newcount\symc
\newcommand\symrange[2]{\symc=#1 \@whilenum\symc<#2 \do{\expandafter\symactive\expandafter{\the\symc}\advance\symc by 1 }}
\newcommand\symblankrange[2]{\symc=#1 \@whilenum\symc<#2 \do{\expandafter\symblank\expandafter{\the\symc}\advance\symc by 1 }}
\makeatother
\symrange{"0250}{"02FF}\symactive{"2713}\symactive{"2714}\symactive{"2717}\symactive{"2718}\symactive{"2610}\symactive{"2611}\symactive{"2612}
\symrange{"2600}{"27BF}
\catcode"2705=13 \begingroup\lccode`\~="2705 \lowercase{\endgroup\def~}{{\symfont\char"2713}}\symblank{"26BD}\symblank{"26A1}
\symrange{"2460}{"257F}\symactive{"223C}\symactive{"2248}\symactive{"2260}
\symhyph{"2011}\symhyph{"2E1A}\symblankrange{"1F300}{"1FAFF}\symblank{"FE0F}
% Caractères chinois : WenQuanYi Zen Hei (sous-ensemble GB2312 + pinyin), gras/italique simulés
\setCJKmainfont{WenQuanYiZenHei-subset.ttf}[Path=../../../fonts/,AutoFakeBold=3,AutoFakeSlant=0.2]
\xeCJKsetup{CJKspace=false}
% Pinyin : voyelles à caron (ǎ ǐ ǒ ǔ ǖ ǘ ǚ ǜ) absentes de la police principale -> caractères actifs
\newfontface\pyfont{WenQuanYiZenHei-subset.ttf}[Path=../../../fonts/]
\newcommand\pyactive[1]{\catcode#1=13 \begingroup\lccode`\~=#1 \lowercase{\endgroup\def~}{{\pyfont\char#1}}}
\pyactive{"01CD}\pyactive{"01CE}\pyactive{"01CF}\pyactive{"01D0}\pyactive{"01D1}\pyactive{"01D2}\pyactive{"01D3}\pyactive{"01D4}\pyactive{"01D5}\pyactive{"01D6}\pyactive{"01D7}\pyactive{"01D8}\pyactive{"01D9}\pyactive{"01DA}\pyactive{"01DB}\pyactive{"01DC}
% Maths en sans-serif (Fira Math) avec chiffres et lettres droites de Plus
% Jakarta, pour que les formules se fondent dans le texte.
\usepackage[math-style=ISO,bold-style=ISO]{unicode-math}
\setmathfont{FiraMath-Regular.otf}[Path=../../../fonts/]
\setmathfont{PlusJakartaSans-Medium.ttf}[Path=../../../fonts/,range={up/num,up/{Latin,latin},\mathup}]
\usepackage{icomma} % virgule décimale sans espace en mode math
% Caractères mathématiques Unicode absents des polices de texte (Plus Jakarta,
% Archivo Black) : rendus par leur équivalent mathématique, en texte comme en
% formule — sinon ils s'affichent en carré vide (ex. « Arithmétique dans ℤ »).
\usepackage{newunicodechar}
\newunicodechar{ℝ}{\ensuremath{\mathbb{R}}}
\newunicodechar{ℤ}{\ensuremath{\mathbb{Z}}}
\newunicodechar{ℂ}{\ensuremath{\mathbb{C}}}
\newunicodechar{ℕ}{\ensuremath{\mathbb{N}}}
\newunicodechar{ℚ}{\ensuremath{\mathbb{Q}}}
\newunicodechar{𝔻}{\ensuremath{\mathbb{D}}}
\newunicodechar{α}{\ensuremath{\alpha}}
\newunicodechar{β}{\ensuremath{\beta}}
\newunicodechar{γ}{\ensuremath{\gamma}}
\newunicodechar{δ}{\ensuremath{\delta}}
\newunicodechar{ε}{\ensuremath{\varepsilon}}
\newunicodechar{θ}{\ensuremath{\theta}}
\newunicodechar{λ}{\ensuremath{\lambda}}
\newunicodechar{μ}{\ensuremath{\mu}}
\newunicodechar{σ}{\ensuremath{\sigma}}
\newunicodechar{τ}{\ensuremath{\tau}}
\newunicodechar{φ}{\ensuremath{\varphi}}
\newunicodechar{ψ}{\ensuremath{\psi}}
\newunicodechar{ω}{\ensuremath{\omega}}
\newunicodechar{Σ}{\ensuremath{\Sigma}}
% majuscules produites par \MakeUppercase dans les titres (α-aminés -> Α)
\newunicodechar{Α}{\ensuremath{\alpha}}
\newunicodechar{Β}{\ensuremath{\beta}}
\newunicodechar{Γ}{\ensuremath{\Gamma}}
\newunicodechar{Δ}{\ensuremath{\Delta}}
\newunicodechar{Ε}{\ensuremath{\varepsilon}}
\newunicodechar{Θ}{\ensuremath{\Theta}}
\newunicodechar{Λ}{\ensuremath{\Lambda}}
\newunicodechar{Μ}{\ensuremath{\mu}}
\newunicodechar{Τ}{\ensuremath{\tau}}
\newunicodechar{Φ}{\ensuremath{\Phi}}
\newunicodechar{Ψ}{\ensuremath{\Psi}}
\newunicodechar{Ω}{\ensuremath{\Omega}}
\newunicodechar{↦}{\ensuremath{\mapsto}}
\newunicodechar{∈}{\ensuremath{\in}}
\newunicodechar{∉}{\ensuremath{\notin}}
\newunicodechar{∧}{\ensuremath{\wedge}}
\newunicodechar{⇔}{\ensuremath{\Leftrightarrow}}
\newunicodechar{⟺}{\ensuremath{\Longleftrightarrow}}
\newunicodechar{⇒}{\ensuremath{\Rightarrow}}
\newunicodechar{⟹}{\ensuremath{\Longrightarrow}}
\newunicodechar{∘}{\ensuremath{\circ}}
\newunicodechar{∪}{\ensuremath{\cup}}
\newunicodechar{∩}{\ensuremath{\cap}}
\newunicodechar{≡}{\ensuremath{\equiv}}
\newunicodechar{⊂}{\ensuremath{\subset}}
\newunicodechar{⊥}{\ensuremath{\perp}}
\newunicodechar{∃}{\ensuremath{\exists}}
\newunicodechar{∀}{\ensuremath{\forall}}
\newunicodechar{∛}{\ensuremath{\sqrt[3]{\,}}}
\newunicodechar{∜}{\ensuremath{\sqrt[4]{\,}}}
\newunicodechar{⃗}{\ensuremath{{}^{\to}}}
\newunicodechar{ⁿ}{\ifmmode^{n}\else\textsuperscript{\ensuremath{n}}\fi}
\newunicodechar{ˣ}{\ifmmode^{x}\else\textsuperscript{\ensuremath{x}}\fi}
\newunicodechar{ᵇ}{\ifmmode^{b}\else\textsuperscript{\ensuremath{b}}\fi}
\newunicodechar{ʳ}{\ifmmode^{r}\else\textsuperscript{\ensuremath{r}}\fi}
\newunicodechar{ᵘ}{\ifmmode^{u}\else\textsuperscript{\ensuremath{u}}\fi}
\newunicodechar{ʸ}{\ifmmode^{y}\else\textsuperscript{\ensuremath{y}}\fi}
\newunicodechar{ᵃ}{\ifmmode^{a}\else\textsuperscript{\ensuremath{a}}\fi}
\newunicodechar{ᵉ}{\ifmmode^{e}\else\textsuperscript{\ensuremath{e}}\fi}
\newunicodechar{ᵏ}{\ifmmode^{k}\else\textsuperscript{\ensuremath{k}}\fi}
\newunicodechar{ᵖ}{\ifmmode^{p}\else\textsuperscript{\ensuremath{p}}\fi}
\newunicodechar{ᴺ}{\ifmmode^{N}\else\textsuperscript{\ensuremath{N}}\fi}
\newunicodechar{ᶜ}{\ifmmode^{c}\else\textsuperscript{\ensuremath{c}}\fi}
\newunicodechar{ʰ}{\ifmmode^{h}\else\textsuperscript{\ensuremath{h}}\fi}
\newunicodechar{ᵠ}{\ifmmode^{q}\else\textsuperscript{\ensuremath{q}}\fi}
\newunicodechar{ₙ}{\ifmmode_{n}\else\textsubscript{\ensuremath{n}}\fi}
\newunicodechar{ₐ}{\ifmmode_{a}\else\textsubscript{\ensuremath{a}}\fi}
\newunicodechar{ᵢ}{\ifmmode_{i}\else\textsubscript{\ensuremath{i}}\fi}
\newunicodechar{ᵣ}{\ifmmode_{r}\else\textsubscript{\ensuremath{r}}\fi}
\newunicodechar{ₖ}{\ifmmode_{k}\else\textsubscript{\ensuremath{k}}\fi}
\newunicodechar{ᵦ}{\ifmmode_{\beta}\else\textsubscript{\ensuremath{\beta}}\fi}
\newunicodechar{ₓ}{\ifmmode_{x}\else\textsubscript{\ensuremath{x}}\fi}
\newfontfamily\popdisplay{ArchivoBlack-Regular.ttf}[Path=../../../fonts/]
\newfontfamily\poplogo{SpaceGrotesk-Bold.ttf}[Path=../../../fonts/]
\newfontfamily\popsemi{PlusJakartaSans-SemiBold.ttf}[Path=../../../fonts/]
\newfontfamily\popxbold{PlusJakartaSans-ExtraBold.ttf}[Path=../../../fonts/]
\newfontfamily\popxboldls{PlusJakartaSans-ExtraBold.ttf}[Path=../../../fonts/,LetterSpace=3.0]

\setlist[enumerate]{leftmargin=1.7em,itemsep=5pt,topsep=4pt}
\setlist[itemize]{leftmargin=1.7em,itemsep=4pt,topsep=4pt,label={\color{poporange}\textbullet}}
\setlength{\parindent}{0pt}
\setlength{\parskip}{5pt}
\renewcommand{\arraystretch}{1.25}

% pgfplots : style Pop par défaut
\pgfplotsset{
  every axis/.append style={axis line style={popink,line width=1pt},tick style={popink},
    grid style={popline,line width=0.5pt},label style={font=\small},tick label style={font=\footnotesize}},
  popcurve/.style={poporange,line width=1.8pt,no markers,smooth},
}
% Même style disponible dans un \draw TikZ simple (hors pgfplots)
\tikzset{popcurve/.style={poporange,line width=1.8pt}}
\tikzset{popgrid/.style={popline,very thin}}
\colorlet{popcurve}{poporange} % le nom du style est parfois utilisé comme couleur

% ============================================================
% Pill / squiggle
% ============================================================
\newcommand{\poppill}[3]{%
  \tikz[baseline=-0.55ex]\node[draw=popink,line width=1.2pt,rounded corners=2.6pt,%
    fill=#2,inner xsep=6.5pt,inner ysep=3pt]{%
    \popxboldls\fontsize{7}{7}\selectfont\color{#3}\MakeUppercase{#1}};%
}
\newcommand{\popsquiggle}[1]{%
  \tikz[baseline=-0.4ex]{
    \draw[#1,line width=2.6pt,line cap=round]
      plot [smooth] coordinates {(0,0.09) (0.34,0.01) (0.68,0.21) (1.02,0.01) (1.36,0.10)};
  }%
}
% Mot-clé surligné (marqueur orange sous le texte)
\newcommand{\cle}[1]{{\popxbold\color{popdark}#1}}

% ============================================================
% En-tête / pied
% ============================================================
\pagestyle{fancy}
\fancyhf{}
\renewcommand{\headrulewidth}{0pt}
\renewcommand{\footrulewidth}{0pt}
\lhead{\raisebox{-0.15ex}{\poplogo\fontsize{13}{13}\selectfont\color{popink}OnBuch\color{poporange}+}}
\rhead{\poppill{\DOCMATIERE\ \textbullet\ \DOCNIVEAU\ \textbullet\ Leçon \DOCLECON}{white}{popink}}
\lfoot{\popxbold\fontsize{7.6}{7.6}\selectfont\color{popmuted}\MakeUppercase{Cours signé OnBuch+}\ \textbullet\ {\popsemi\DOCTITRE}}
\rfoot{\tikz[baseline=-0.6ex]\node[rounded corners=8.5pt,fill=popink,minimum height=17pt,minimum width=17pt,inner xsep=5pt,inner sep=0]{\popxbold\fontsize{8}{8}\selectfont\color{popcream}\thepage\,/\,\pageref*{LastPage}};}

% ============================================================
% Titres
% ============================================================
\newcommand{\chaptitre}[2]{%
  \begin{tcolorbox}[enhanced,colback=poporange,colframe=popink,boxrule=2.6pt,arc=11pt,
      drop shadow=popink,left=15pt,right=15pt,top=12pt,bottom=11pt,before skip=3pt,after skip=18pt]
    \poppill{#2}{popink}{popcream}\par\vspace{9pt}
    {\raggedright\hyphenpenalty=10000\exhyphenpenalty=10000\popdisplay\fontsize{22}{24}\selectfont\color{popcream}\MakeUppercase{#1}\par}
  \end{tcolorbox}
}
% Section numérotée : pastille orange avec le numéro + titre Archivo
\newcounter{popsec}
\newcommand{\coursec}[1]{%
  \stepcounter{popsec}\setcounter{popsub}{0}%
  \par\vspace{16pt}\noindent
  \begin{minipage}{\linewidth}
  \tikz[baseline=-0.7ex]\node[draw=popink,line width=1.6pt,fill=poporange,rounded corners=4pt,minimum size=22pt,inner sep=0]{\popdisplay\fontsize{12}{12}\selectfont\color{popcream}\thepopsec};%
  \hspace{8pt}{\popdisplay\fontsize{15}{17}\selectfont\color{popink}\MakeUppercase{#1}}\par
  \vspace{4pt}\noindent{\color{poporange}\rule{36pt}{3.6pt}}
  \end{minipage}\par\nopagebreak\vspace{8pt}
}
\newcounter{popsub}
\newcommand{\courssub}[1]{%
  \stepcounter{popsub}%
  \par\vspace{9pt}\noindent{\popxbold\fontsize{11.5}{13}\selectfont\color{popink}{\color{poporange}\thepopsec.\thepopsub}\hspace{6pt}#1}\par\nopagebreak\vspace{4pt}
}

% ============================================================
% Boîtes pédagogiques — même recette : fond blanc, bordure épaisse colorée,
% ombre dure encre, badge plein. Titre optionnel [#1].
% ============================================================
\tcbset{popbase/.style={breakable,enhanced,colback=white,boxrule=2.3pt,arc=9pt,
  shadow={3pt}{-3pt}{0pt}{popink},left=14pt,right=14pt,top=11pt,bottom=11pt,before skip=11pt,after skip=15pt,
  fontupper=\popsemi\fontsize{10.6}{14.5}\selectfont\color{popink},
  fontlower=\popsemi\fontsize{10.6}{14.5}\selectfont\color{popink}}}
% Coche dessinée (le glyphe ✓ n'existe pas dans les polices du document)
\newcommand{\popcheck}[1]{\tikz[baseline=-0.5ex]\draw[#1,line width=1.6pt,line cap=round,line join=round](-0.13,0.0)--(-0.04,-0.09)--(0.13,0.1);}
\providecommand{\coche}{\popcheck{popgreen}}
\providecommand{\resultat}[1]{\par\smallskip\noindent\fcolorbox{popgreen}{popgreenL}{\parbox{\dimexpr\linewidth-2\fboxsep-2\fboxrule\relax}{\textbf{Résultat.} #1}}\par\smallskip}
\newcommand{\popboxhead}[3]{\poppill{#1}{#2}{white}\ifx\relax#3\relax\else\hspace{6pt}{\popxbold\fontsize{10.6}{12}\selectfont\color{popink}#3}\fi\par\vspace{7pt}}

\newtcolorbox{aretenir}[1][]{popbase,colframe=popgreen,before upper={\popboxhead{À retenir}{popgreen}{#1}}}
% Texte en chinois (document de lecture, dialogue, extrait) : cadre
% d'étude. Un titre optionnel (source, auteur) s'affiche dans l'en-tête.
\newtcolorbox{textebox}[1][]{popbase,colframe=popblue,before upper={\popboxhead{文本 · Texte}{popblue}{#1}},after upper={}}
% Marques ✓ / ✗ (forme correcte / fautive) souvent utilisées par les modèles
\providecommand{\cross}{\textcolor{poppink}{\ensuremath{\times}}}
\providecommand{\xmark}{\cross}\providecommand{\crossmark}{\cross}\providecommand{\cmark}{\ensuremath{\checkmark}}
% Ligne de réponse / justification à compléter (inventée par les modèles)
\providecommand{\justification}[1]{\par\noindent\textit{Justification~:} \dotfill\par}
% \textipa (package tipa non chargé) : la police ne couvre pas l'API -> simple passage
\providecommand{\textipa}[1]{#1}
% Soulignement (ulem non chargé) : \uline, \ul
\providecommand{\uline}[1]{\underline{#1}}\providecommand{\ul}[1]{\underline{#1}}
% \glossaire{\item ...} inventée par les modèles : liste de glossaire
\providecommand{\glossaire}[1]{\begin{itemize}#1\end{itemize}}
% \alert (beamer) inventée par les modèles : texte mis en évidence
\providecommand{\alert}[1]{\textbf{\textcolor{poppink}{#1}}}
% variantes ASCII d'ordinaux écrites en macro
\providecommand{\iere}{\textsuperscript{re}}\providecommand{\ieme}{\textsuperscript{e}}\providecommand{\ere}{\textsuperscript{re}}\providecommand{\eme}{\textsuperscript{e}}\providecommand{\er}{\textsuperscript{er}}\providecommand{\ie}{\textsuperscript{e}}
% Ordinaux français écrits en macro par les modèles : 1\ière, 3\ième
\newcommand{\ière}{\textsuperscript{re}}\newcommand{\ième}{\textsuperscript{e}}
% \exemple (inventée par les modèles) : étiquette « Exemple : »
\providecommand{\exemple}{\textbf{Exemple~:}\ }
% \square utilisable aussi hors mode math (cases à cocher)
\AtBeginDocument{\let\squareMath\square\renewcommand{\square}{\ensuremath{\squareMath}}}
% Mot ou phrase en chinois à l'intérieur d'un paragraphe français
\newcommand{\zh}[1]{\ifmmode\text{#1}\else{#1}\fi}
\let\eng\zh \let\esp\zh \let\all\zh % alias tolérés
% flèches d'intonation inventées par les modèles
\providecommand{\textsearrow}{}\renewcommand{\textsearrow}{\ensuremath{\searrow}}\providecommand{\textnearrow}{}\renewcommand{\textnearrow}{\ensuremath{\nearrow}}
% \fr{..} : mot français (inventé par les modèles)
\providecommand{\fr}[1]{\textit{#1}}
% zhbox : encadré de texte chinois inventé par les modèles
\newenvironment{zhbox}{\begin{quote}}{\end{quote}}
% \vs : « versus » inventé par les modèles
\providecommand{\vs}{\textit{vs}\ }
% \blank : case à compléter inventée par les modèles
\providecommand{\blank}[1][]{\underline{\hspace{1.6em}}}
\newtcolorbox{exemplebox}[1][]{popbase,colframe=poporange,before upper={\popboxhead{Exemple}{poporange}{#1}}}
\newtcolorbox{definition}[1][]{popbase,colframe=popblue,before upper={\popboxhead{Définition}{popblue}{#1}}}
\newtcolorbox{propriete}[1][]{popbase,colframe=poppurple,before upper={\popboxhead{Propriété}{poppurple}{#1}}}
\newtcolorbox{methode}[1][]{popbase,colframe=popgold,before upper={\popboxhead{Méthode}{popgold}{#1}}}
\newtcolorbox{attention}[1][]{popbase,colframe=poppink,before upper={\popboxhead{Attention}{poppink}{#1}}}
\newtcolorbox{experience}[1][]{popbase,colframe=popblue,colback=popblueL,before upper={\popboxhead{Expérience}{popblue}{#1}}}
% « Le savais-tu ? » : carte violette pleine (contexte, culture, Cameroun)
\newtcolorbox{savaistu}[1][]{popbase,colback=poppurple,colframe=popink,
  fontupper=\popsemi\fontsize{10.4}{14.2}\selectfont\color{white},
  before upper={\poppill{Le savais-tu\,?}{white}{poppurple}\ifx\relax#1\relax\else\hspace{6pt}{\popxbold\color{white}#1}\fi\par\vspace{7pt}}}
% Exercice résolu : énoncé en haut, \tcblower puis la solution sur fond crème
\newcounter{popexo}\newcounter{popexr}
\newtcolorbox{exoresolu}[1][]{popbase,colframe=poporange,colbacklower=popcream,
  segmentation style={popink,line width=1.2pt,dashed},
  before upper={\stepcounter{popexr}\popboxhead{Exercice résolu \thepopexr}{poporange}{#1}},
  before lower={\poppill{Solution}{popink}{popcream}\par\vspace{6pt}}}
% Exercice (non résolu en place) — la correction est dans la partie Corrigés
\newtcolorbox{exercice}[1][]{popbase,colframe=popink,
  before upper={\stepcounter{popexo}\popboxhead{Exercice \thepopexo}{popink}{#1}}}
% Corrigé
\newtcolorbox{corrige}[1][]{popbase,colframe=popgreen,colback=popgreenL,
  before upper={\popboxhead{Corrigé}{popgreen}{#1}}}
% Objectifs / prérequis / bilan
\newtcolorbox{objectifs}{popbase,colframe=popink,colback=poporangeL,
  before upper={\popboxhead{Objectifs de la leçon}{popink}{}}}
\newtcolorbox{prerequis}{popbase,colframe=popink,colback=popblueL,
  before upper={\popboxhead{Prérequis}{popblue}{}}}
\newtcolorbox{bilan}{popbase,colframe=popink,colback=white,boxrule=2.8pt,
  before upper={\popboxhead{Fiche bilan}{poporange}{L'essentiel à connaître pour l'examen}}}
% Figure encadrée avec légende
\newtcolorbox{popfigure}[1][]{enhanced,breakable=false,colback=white,colframe=popline,boxrule=1.4pt,arc=8pt,
  left=10pt,right=10pt,top=10pt,bottom=8pt,before skip=10pt,after skip=12pt,halign=center,
  lower separated=false,halign lower=center,
  fontlower=\popsemi\fontsize{9}{11.5}\selectfont\color{popmuted},
  #1}
\newcommand{\legende}[1]{\par\vspace{6pt}{\popsemi\fontsize{9}{11.5}\selectfont\color{popmuted}#1}}

% ============================================================
% Signature OnBuch+
% ============================================================
\newcommand{\poplogoinline}[1]{{\poplogo\fontsize{#1}{#1}\selectfont\color{popink}OnBuch\color{poporange}+}}
\newcommand{\signatureonbuch}{%
  \begin{tcolorbox}[enhanced,colback=popink,colframe=popink,boxrule=2.2pt,arc=10pt,
      drop shadow=poporange,left=14pt,right=14pt,top=10pt,bottom=10pt,before skip=0pt,after skip=0pt]
    \tikz[baseline=-0.7ex]\node[circle,fill=popgreen,draw=popcream,line width=1.3pt,minimum size=22pt,inner sep=0]{\popcheck{white}};%
    \hspace{9pt}\begin{minipage}[c]{0.62\linewidth}
      {\popxbold\fontsize{12}{13}\selectfont\color{popcream}Signé\ \textbullet\ {\poplogo OnBuch\color{poporange}+}}\par\vspace{2pt}
      {\raggedright\popsemi\fontsize{8}{10.5}\selectfont\color{popline}\MakeUppercase{Équipe pédagogique OnBuch+}\\\MakeUppercase{Conforme au programme officiel MINESEC}\par}
    \end{minipage}\hfill
    {\poplogo\fontsize{22}{22}\selectfont\color{popcream}OnBuch\color{poporange}+}
  \end{tcolorbox}%
}

% ============================================================
% Page de garde
% #1 titre · #2 matière · #3 niveau · #4 module · #5 n° leçon · #6 durée ·
% #7 description / objectifs (texte ou itemize)
% ============================================================
\newcommand{\pagegarde}[7]{%
  \thispagestyle{empty}
  \newgeometry{a4paper,margin=1.7cm,top=1.6cm,bottom=1.4cm}%
  \begin{tikzpicture}[remember picture,overlay]
    \fill[poporange!18] ([xshift=-3.2cm,yshift=-3.6cm]current page.north east) circle (4.6cm);
    \fill[poppurple!14] ([xshift=2.4cm,yshift=7.6cm]current page.south west) circle (3.8cm);
  \end{tikzpicture}%
  {\poplogo\fontsize{30}{30}\selectfont\color{popink}OnBuch\color{poporange}+}\hfill
  \raisebox{0.6ex}{\poppill{Cours signé \textbullet\ Premium}{popink}{popcream}}\par
  \vspace{1pt}\popsquiggle{poppurple}\par
  \vspace{8pt}
  \begin{tcolorbox}[enhanced,colback=poporange,colframe=popink,boxrule=2.8pt,arc=13pt,
      drop shadow=popink,left=18pt,right=18pt,top=12pt,bottom=13pt,after skip=10pt]
    \poppill{#2 \textbullet\ #3}{popink}{popcream}\hspace{5pt}\poppill{#4}{white}{popink}\par\vspace{8pt}
    {\popdisplay\fontsize{14}{14}\selectfont\color{popink}LEÇON #5}\par\vspace{4pt}
    {\raggedright\hyphenpenalty=10000\exhyphenpenalty=10000\popdisplay\fontsize{25}{27}\selectfont\color{popcream}\MakeUppercase{#1}\par}
    \vspace{8pt}
    {\popxbold\fontsize{9.5}{9.5}\selectfont\color{popink}\MakeUppercase{Durée conseillée : #6}}
  \end{tcolorbox}
  \begin{tcolorbox}[enhanced,colback=white,colframe=popink,boxrule=2.2pt,arc=10pt,
      drop shadow=popink,left=16pt,right=16pt,top=10pt,bottom=10pt,after skip=8pt]
    \poppill{Dans cette leçon}{poppurple}{white}\par\vspace{5pt}
    \popsemi\fontsize{9.3}{12.6}\selectfont\color{popink}#7
  \end{tcolorbox}
  \noindent\begin{minipage}[t]{0.487\linewidth}
    \begin{tcolorbox}[enhanced,colback=popgreen,colframe=popink,boxrule=2.2pt,arc=10pt,
        drop shadow=popink,left=11pt,right=11pt,top=10pt,bottom=10pt,height=3.1cm,valign=center]
      \tcbox[colback=white,colframe=popink,boxrule=1.3pt,arc=5pt,boxsep=0pt,nobeforeafter,
        left=3pt,right=3pt,top=3pt,bottom=3pt,valign=center]{\qrcode[height=1.9cm]{https://play.google.com/store/apps/details?id=com.luvvix.onbuch&hl=fr}}%
      \hspace{8pt}\begin{minipage}[c]{0.5\linewidth}
        {\popdisplay\fontsize{11}{12.5}\selectfont\color{white}\MakeUppercase{Télécharge OnBuch}}\par\vspace{4pt}
        {\raggedright\popsemi\fontsize{8.4}{11}\selectfont\color{white}Gratuit \textbullet\ Google Play\\Tuteur IA, annales, concours, résultats\par}
      \end{minipage}
    \end{tcolorbox}
  \end{minipage}\hfill
  \begin{minipage}[t]{0.487\linewidth}
    \begin{tcolorbox}[enhanced,colback=poppurple,colframe=popink,boxrule=2.2pt,arc=10pt,
        drop shadow=popink,left=11pt,right=11pt,top=10pt,bottom=10pt,height=3.1cm,valign=center]
      \tikz[baseline=-0.7ex]\node[circle,fill=white,draw=popink,line width=1.3pt,minimum size=1.6cm,inner sep=0]{\popdisplay\fontsize{24}{24}\selectfont\color{poppurple}+};%
      \hspace{8pt}\begin{minipage}[c]{0.6\linewidth}
        {\popdisplay\fontsize{11}{12.5}\selectfont\color{white}\MakeUppercase{Passe à OnBuch+}}\par\vspace{4pt}
        {\raggedright\popsemi\fontsize{8.4}{11}\selectfont\color{white}Tous les cours signés, Tuteur IA illimité, fiches PDF — onglet OnBuch+ de l'app\par}
      \end{minipage}
    \end{tcolorbox}
  \end{minipage}
  \vspace{8pt}
  \popcontenu
  \vfill
  \signatureonbuch
  \restoregeometry
  \newpage
}

% Rangée « Ce cours contient » (page de garde)
\newcommand{\poptile}[3]{% #1 couleur #2 gros texte #3 légende
  \begin{tcolorbox}[enhanced,colback=#1,colframe=popink,boxrule=2pt,arc=9pt,shadow={2.5pt}{-2.5pt}{0pt}{popink},
      left=6pt,right=6pt,top=9pt,bottom=9pt,halign=center,valign=center,height=1.75cm,width=\linewidth,nobeforeafter]
    {\popdisplay\fontsize{12.5}{13}\selectfont\color{white}\MakeUppercase{#2}}\par\vspace{3pt}
    {\centering\popsemi\fontsize{7.8}{9.5}\selectfont\color{white}#3\par}
  \end{tcolorbox}}
\newcommand{\popcontenu}{%
  \par\noindent{\popxboldls\fontsize{7.5}{7.5}\selectfont\color{popmuted}CE COURS SIGNÉ CONTIENT}\par\vspace{7pt}
  \noindent\begin{minipage}[t]{0.235\linewidth}\poptile{popblue}{Cours}{complet, illustré, pas à pas}\end{minipage}\hfill
  \begin{minipage}[t]{0.235\linewidth}\poptile{poporange}{Méthodes}{et exercices résolus}\end{minipage}\hfill
  \begin{minipage}[t]{0.235\linewidth}\poptile{poppink}{Exercices}{type examen + corrigés}\end{minipage}\hfill
  \begin{minipage}[t]{0.235\linewidth}\poptile{popgreen}{Fiche bilan}{l'essentiel à retenir}\end{minipage}\par}

% Sommaire
\newtcolorbox{sommaire}{enhanced,colback=white,colframe=popink,boxrule=2.2pt,arc=10pt,
  drop shadow=popink,left=18pt,right=18pt,top=13pt,bottom=13pt,before skip=6pt,after skip=20pt,
  before upper={\poppill{Sommaire}{poppurple}{white}\par\vspace{9pt}\popsemi\fontsize{10.6}{14.5}\selectfont\color{popink}}}

% Signature de fin de cours — poussée tout en bas de la dernière page
\newcommand{\findecours}{%
  \par\vfill\nopagebreak
  \begin{center}\popsquiggle{poporange}\end{center}\vspace{4pt}
  \signatureonbuch
}
\AtBeginDocument{\def\llbracket{\mathopen{[\mkern-3.5mu[}}\def\rrbracket{\mathclose{]\mkern-3.5mu]}}} % intervalles d'entiers (glyphe ⟦ absent de la police)
% Glyphes absents de Fira Math : redéfinis à partir de glyphes disponibles
\AtBeginDocument{%
  \def\setminus{\mathbin{\backslash}}\let\smallsetminus\setminus
  \def\vdots{\mathord{\vbox{\baselineskip3.2pt\lineskiplimit0pt\kern4pt\hbox{.}\hbox{.}\hbox{.}}}}%
  \def\ddots{\mathinner{\mkern1mu\raise6.5pt\hbox{.}\mkern2mu\raise3.6pt\hbox{.}\mkern2mu\raise0.7pt\hbox{.}\mkern1mu}}%
  \def\triangle{\mathord{\scalebox{0.9}{$\symup{\Delta}$}}}%
  \def\nabla{\mathord{\raisebox{\depth}{\scalebox{1}[-1]{$\symup{\Delta}$}}}}%
  \def\checkmark{\popcheck{popdark}}%
  \def\star{\mathbin{\ast}}\def\bigstar{\mathord{\ast}}%
  \def\diamond{\mathbin{\scalebox{0.6}{$\Diamond$}}}%
  \def\bigcup{\mathop{\mathchoice{\raisebox{-0.3em}{\scalebox{1.7}{$\cup$}}}{\scalebox{1.25}{$\cup$}}{\cup}{\cup}}\displaylimits}%
  \def\bigcap{\mathop{\mathchoice{\raisebox{-0.3em}{\scalebox{1.7}{$\cap$}}}{\scalebox{1.25}{$\cap$}}{\cap}{\cap}}\displaylimits}%
  \def\lesseqqgtr{\mathrel{\lessgtr}}\def\bigoplus{\mathop{\oplus}\displaylimits}%
}
\providecommand{\ssi}{si et seulement si} % abréviation fréquente
\AtBeginDocument{\providecommand{\wideparen}{\overparen}} % arc de cercle

%% FIGFIX-BEGIN v4 — correctifs globaux des figures (docs/onbuch-generation/tools/figfix)
\usepackage{adjustbox}\usepackage{etoolbox}
% toute figure TikZ/pgfplots plus large que la ligne est réduite à la largeur disponible
\BeforeBeginEnvironment{tikzpicture}{\begin{adjustbox}{max width=\linewidth}}
\AfterEndEnvironment{tikzpicture}{\end{adjustbox}}
% légendes pgfplots lisibles par-dessus les courbes
\pgfplotsset{every axis/.append style={legend style={fill=white,fill opacity=0.92,text opacity=1,draw=black!15,inner sep=3pt}}}
% étiquettes d'arêtes (syntaxe edge["..."]) avec fond blanc
\tikzset{every edge quotes/.append style={fill=white,inner sep=1.2pt}}
%% FIGFIX-END
\begin{document}

\pagegarde{Leçon 12 — Vocabulaire et compréhension de texte : paludisme}{Chinois}{Tle A}{Module 2 : Environnement et santé}{ZH12}{2 h}{Cette leçon de vocabulaire et de compréhension de texte porte sur le paludisme, thème sanitaire majeur au Cameroun. À partir d'un texte support original en chinois simplifié, les élèves acquièrent le lexique de la maladie, de sa transmission et de sa prévention, et s'entraînent à répondre à des questions de compréhension.
\vspace{4pt}
\begin{multicols}{2}\small
\begin{itemize}
\item À la fin de cette leçon, je sais lire à voix haute un court texte chinois sur le paludisme en respectant le pinyin et les tons.
\item À la fin de cette leçon, je sais reconnaître et traduire au moins quinze mots de vocabulaire liés à la santé et au paludisme.
\item À la fin de cette leçon, je sais identifier les clés (radicaux) et tracer l'ordre des traits de quelques caractères du thème.
\item À la fin de cette leçon, je sais comprendre globalement et en détail un texte documentaire simple en chinois.
\item À la fin de cette leçon, je sais répondre par écrit à des questions de compréhension en phrases complètes.
\item À la fin de cette leçon, je sais employer les collocations essentielles du thème (得疟疾, 发烧, 预防疾病) dans des phrases correctes.
\end{itemize}\end{multicols}}

\begin{sommaire}
\begin{multicols}{2}
\begin{enumerate}
\item Mise en route et texte support
\item Glossaire du thème
\item Clés (radicaux) et ordre des traits
\item Compréhension détaillée du texte
\item Collocations et production guidée
\item Activité d'intégration
\item Méthodes et exercices résolus
\item Exercices
\item Corrigés des exercices
\item Fiche bilan
\end{enumerate}
\end{multicols}
\end{sommaire}

\begin{objectifs}
\begin{itemize}
\item À la fin de cette leçon, je sais lire à voix haute un court texte chinois sur le paludisme en respectant le pinyin et les tons.
\item À la fin de cette leçon, je sais reconnaître et traduire au moins quinze mots de vocabulaire liés à la santé et au paludisme.
\item À la fin de cette leçon, je sais identifier les clés (radicaux) et tracer l'ordre des traits de quelques caractères du thème.
\item À la fin de cette leçon, je sais comprendre globalement et en détail un texte documentaire simple en chinois.
\item À la fin de cette leçon, je sais répondre par écrit à des questions de compréhension en phrases complètes.
\item À la fin de cette leçon, je sais employer les collocations essentielles du thème (得疟疾, 发烧, 预防疾病) dans des phrases correctes.
\item À la fin de cette leçon, je sais rédiger quelques phrases guidées pour décrire une maladie, ses symptômes et sa prévention.
\item À la fin de cette leçon, je sais comparer brièvement la lutte contre le paludisme au Cameroun et en Chine.
\end{itemize}
\end{objectifs}

\begin{prerequis}
\begin{itemize}
\item Vocabulaire de base du corps et de la santé vu en Première (头, 疼, 病, 医院, 医生)
\item Structure sujet-verbe-complément et particule 了 pour le procès accompli
\item Emploi de 很 / 非常 avec les adjectifs et de 会 pour la capacité ou la probabilité
\item Lecture du pinyin et maîtrise des quatre tons
\item Notions sur les radicaux et l'ordre général des traits (de haut en bas, de gauche à droite)
\end{itemize}
\end{prerequis}

\newpage

\coursec{Mise en route et texte support}

\courssub{Situation d'accroche et problématique}

À Yaoundé comme à Douala, chaque saison des pluies ramène son lot de moustiques et de cas de paludisme. Au centre de santé de quartier, l'infirmière accueille un malade et lui demande de décrire ses symptômes : fièvre, frissons, maux de tête, fatigue générale. Le diagnostic est presque toujours le même : le paludisme, cette maladie que tous les Camerounais connaissent bien.

Imaginez maintenant la scène suivante : une délégation de médecins chinois en mission au Cameroun visite votre lycée dans le cadre d'une campagne de sensibilisation. Ils vous posent des questions simples, mais en chinois : qu'est-ce que le paludisme ? Comment se transmet-il ? Comment s'en protège-t-on ? Seriez-vous capable de leur répondre ?

\textbf{Problématique de la leçon :} comment lire, comprendre et exploiter un court texte chinois sur le paludisme afin de pouvoir définir la maladie, expliquer sa transmission, décrire ses symptômes et présenter les moyens de prévention ?

\courssub{Brainstorming : que savez-vous du paludisme ?}

Avant de lire le texte chinois, activons nos connaissances en français. Le paludisme est un sujet que vous connaissez tous, et c'est justement cette familiarité qui va vous aider à deviner le sens des mots chinois inconnus.

\begin{experience}[Brainstorming en classe]
En groupes de quatre, répondez en français aux questions suivantes, puis notez vos idées au tableau :
\begin{enumerate}
\item Quels sont les \cle{symptômes} du paludisme ? (fièvre, frissons, maux de tête, fatigue, vomissements\ldots)
\item Comment le paludisme se \cle{transmet}-il ? (par la piqûre du moustique anophèle femelle)
\item Quelles sont les \cle{conséquences} du paludisme ? (absentéisme scolaire, hospitalisations, décès, surtout chez les enfants)
\item Comment peut-on \cle{prévenir} le paludisme ? (moustiquaires imprégnées, assainissement, destruction des eaux stagnantes, consultation rapide à l'hôpital)
\end{enumerate}
Chaque groupe présente ensuite une réponse à la classe. Gardez ces quatre catégories en tête : elles correspondent exactement à la structure du texte chinois que nous allons lire.
\end{experience}

\begin{attention}[Un mot de vocabulaire à manier avec prudence]
En français, on dit souvent que le moustique « transmet le virus du paludisme ». En réalité, le paludisme est causé par un \emph{parasite} (le \emph{Plasmodium}), pas par un virus. Notre texte chinois utilise le mot courant \zh{病毒} (bìngdú, « virus, agent pathogène ») dans son sens large et populaire de « ce qui rend malade ». Retenez simplement que dans la langue de tous les jours, en chinois comme en français, on emploie souvent des termes approximatifs.
\end{attention}

\courssub{Présentation du titre : \zh{疟疾} (nüèji)}

Le titre du texte est un seul mot : \zh{疟疾}, pinyin \emph{nüèji}, « le paludisme, la malaria ».

\begin{definition}[\zh{疟疾} nüèji]
Nom signifiant « paludisme ». C'est un mot à deux caractères :
\begin{itemize}
\item \zh{疟} (nüè) : caractère spécifique au nom de cette maladie ; il contient la clé de la maladie \zh{疒} (nè), que l'on retrouve dans \zh{病} (bìng, « maladie »), \zh{疼} (téng, « douloureux ») ou \zh{症} (zhèng, « symptôme »).
\item \zh{疾} (jí) : « maladie », que l'on retrouve dans \zh{疾病} (jíbìng, « maladie »).
\end{itemize}
Attention à la prononciation de \zh{疟} : la voyelle est \textbf{ü} (les lèvres arrondies comme pour dire « u » français, la langue placée comme pour « i »), avec le quatrième ton descendant : \emph{nüè}.
\end{definition}

\begin{attention}[La voyelle ü après n]
Après les consonnes \textbf{n} et \textbf{l}, le pinyin conserve les deux points sur le ü : \emph{nüè}, \emph{lǜ}. En revanche, après \textbf{j}, \textbf{q}, \textbf{x}, on écrit « u » mais on prononce toujours ü : \zh{去} s'écrit \emph{qù} mais se prononce avec un ü. Erreur fréquente des francophones : prononcer \emph{nüè} comme « noué ». Entraînez-vous : \emph{nǐ} (tu) → \emph{nǚ} (femme) → \emph{nüè} (palu- dans 疟疾).
\end{attention}

\courssub{Première lecture : saisir le sens global}

\begin{methode}[Comment comprendre un texte chinois en quatre étapes]
\begin{enumerate}
\item \textbf{Lire le titre.} Le titre annonce le thème : ici \zh{疟疾} (nüèji), le paludisme. Tout le texte tournera autour de ce mot.
\item \textbf{Repérer les mots connus.} Soulignez les mots déjà appris : \zh{是}, \zh{人}, \zh{孩子}, \zh{我们}, \zh{应该}, \zh{医院}\ldots{} Ils forment le squelette du texte.
\item \textbf{Deviner les mots inconnus par le contexte et les radicaux.} Un caractère avec la clé de la maladie \zh{疒} parle probablement de santé ; un caractère avec la clé de l'insecte \zh{虫} désigne probablement un insecte (\zh{蚊} dans \zh{蚊子}, « moustique »).
\item \textbf{Vérifier avec le glossaire.} Comparez vos hypothèses avec le tableau de vocabulaire et corrigez vos premières impressions.
\end{enumerate}
\end{methode}

\begin{experience}[Lecture silencieuse chronométrée]
Lisez le texte ci-dessous en silence pendant deux minutes, \textbf{sans regarder la traduction}. Appliquez la méthode en quatre étapes : soulignez les mots connus, entourez les mots inconnus, et essayez de répondre à cette question unique : \emph{De quoi parle ce texte, et quelles sont ses grandes parties ?} Puis comparez vos réponses avec votre voisin.
\end{experience}

\begin{textebox}[疟疾 — Le paludisme]
疟疾是一种很常见的疾病。\\
\emph{Nüèji shì yì zhǒng hěn chángjiàn de jíbìng.}\\
« Le paludisme est une maladie très courante. »\\[4pt]
蚊子把疟疾病毒传给人。\\
\emph{Wénzi bǎ nüèji bìngdú chuán gěi rén.}\\
« Le moustique transmet l'agent du paludisme à l'homme. »\\[4pt]
得了疟疾的人会发烧、头疼、全身没有力气。\\
\emph{Dé le nüèji de rén huì fāshāo, tóuténg, quánshēn méiyǒu lìqi.}\\
« La personne qui a attrapé le paludisme a de la fièvre, mal à la tête et n'a plus de forces dans tout le corps. »\\[4pt]
在非洲，很多孩子因为疟疾死亡。\\
\emph{Zài Fēizhōu, hěn duō háizi yīnwèi nüèji sǐwáng.}\\
« En Afrique, beaucoup d'enfants meurent à cause du paludisme. »\\[4pt]
为了预防疟疾，我们应该用蚊帐，保持环境干净，还要及时去医院看病。\\
\emph{Wèile yùfáng nüèji, wǒmen yīnggāi yòng wénzhàng, bǎochí huánjìng gānjìng, hái yào jíshí qù yīyuàn kànbìng.}\\
« Pour prévenir le paludisme, nous devons utiliser des moustiquaires, garder l'environnement propre, et aussi aller consulter à l'hôpital à temps. »
\end{textebox}

\courssub{Lecture à voix haute : corriger les tons}

La lecture à voix haute se fait \textbf{phrase par phrase}. Le professeur lit d'abord, les élèves répètent en chœur, puis individuellement. Voici les points de prononciation à surveiller :

\begin{center}
\begin{tabularx}{\linewidth}{|X|X|X|}
\hline
\rowcolor{popblueL} Mot & Pinyin & Piège de prononciation \\
\hline
疟疾 & nüèji & voyelle ü + 4\textsuperscript{e} ton descendant, puis 2\textsuperscript{e} ton montant \\
\hline
蚊子 & wénzi & le second caractère 子 se lit au \textbf{ton neutre} (zi), sans accent \\
\hline
发烧 & fāshāo & deux premiers tons plats et hauts, bien tenus \\
\hline
头疼 & tóuténg & deux 2\textsuperscript{es} tons montants successifs : ne pas les aplatir \\
\hline
力气 & lìqi & 气 au ton neutre dans ce mot \\
\hline
因为 & yīnwèi & 为 se lit ici \emph{wèi} (4\textsuperscript{e} ton), pas \emph{wéi} \\
\hline
蚊帐 & wénzhàng & 2\textsuperscript{e} ton montant puis 4\textsuperscript{e} ton descendant \\
\hline
及时 & jíshí & deux 2\textsuperscript{es} tons : « à temps » \\
\hline
\end{tabularx}
\end{center}

\begin{attention}[Les deux lectures du caractère 得]
Le caractère \zh{得} a plusieurs prononciations selon son emploi :
\begin{itemize}
\item \zh{得了疟疾} : 得 se prononce \textbf{dé} (2\textsuperscript{e} ton) ; \zh{得} + maladie signifie « attraper, contracter » une maladie. Exemple : \zh{他得了疟疾。} \emph{Tā dé le nüèji.} « Il a attrapé le paludisme. »
\item \zh{觉得} : 得 se prononce \textbf{de} (ton neutre) ; \zh{觉得} (juéde) signifie « penser, trouver, avoir le sentiment que ». Exemple : \zh{我觉得他很累。} \emph{Wǒ juéde tā hěn lèi.} « Je trouve qu'il est très fatigué. »
\item Il existe aussi la lecture \textbf{de} (ton neutre) comme particule structurale après un verbe : \zh{他说得很好。} \emph{Tā shuō de hěn hǎo.} « Il parle très bien. »
\end{itemize}
Ne confondez jamais \emph{dé le} (avoir attrapé) et \emph{juéde} (penser) : le sens de la phrase en dépend entièrement.
\end{attention}

\courssub{Traduction phrase par phrase}

\begin{center}
\begin{tabularx}{\linewidth}{|X|X|X|}
\hline
\rowcolor{popblueL} Chinois & Pinyin & Traduction française \\
\hline
疟疾是一种很常见的疾病。 & Nüèji shì yì zhǒng hěn chángjiàn de jíbìng. & Le paludisme est une maladie très courante. \\
\hline
蚊子把疟疾病毒传给人。 & Wénzi bǎ nüèji bìngdú chuán gěi rén. & Le moustique transmet l'agent du paludisme à l'homme. \\
\hline
得了疟疾的人会发烧、头疼、全身没有力气。 & Dé le nüèji de rén huì fāshāo, tóuténg, quánshēn méiyǒu lìqi. & Celui qui attrape le paludisme a de la fièvre, mal à la tête et n'a plus de forces dans tout le corps. \\
\hline
在非洲，很多孩子因为疟疾死亡。 & Zài Fēizhōu, hěn duō háizi yīnwèi nüèji sǐwáng. & En Afrique, beaucoup d'enfants meurent à cause du paludisme. \\
\hline
为了预防疟疾，我们应该用蚊帐，保持环境干净，还要及时去医院看病。 & Wèile yùfáng nüèji, wǒmen yīnggāi yòng wénzhàng, bǎochí huánjìng gānjìng, hái yào jíshí qù yīyuàn kànbìng. & Pour prévenir le paludisme, nous devons utiliser des moustiquaires, garder l'environnement propre, et aussi aller consulter à l'hôpital à temps. \\
\hline
\end{tabularx}
\end{center}

Quelques observations linguistiques utiles pour la traduction :

\begin{itemize}
\item \zh{一种} (yì zhǒng) : « une sorte de » ; \zh{种} est ici un classificateur de catégories. La structure \zh{是 + 一种 + adjectif + 的 + nom} correspond au français « est une sorte de \ldots{} » ou simplement « est un(e) \ldots{} ». Rappel : devant un 4\textsuperscript{e} ton, \zh{一} se prononce \emph{yí} ; devant un 1\textsuperscript{er}, 2\textsuperscript{e} ou 3\textsuperscript{e} ton, il se prononce \emph{yì} — d'où \emph{yì zhǒng}.
\item \zh{把} (bǎ) : particule qui fait passer le complément d'objet avant le verbe. \zh{蚊子把病毒传给人} suit le schéma \textbf{Sujet + 把 + objet + verbe + complément} : « le moustique, le virus, le transmet à l'homme ». Le français dit « transmet le virus à l'homme » : l'ordre des mots est donc différent.
\item \zh{会} (huì) dans \zh{会发烧} n'exprime pas la capacité (« savoir ») mais la \textbf{probabilité, la conséquence attendue} : « il se peut que / il arrive que la personne ait de la fièvre ».
\item \zh{因为} (yīnwèi) : « à cause de, parce que » ; \zh{为了} (wèile) : « pour, afin de ». Attention à ne pas les intervertir : l'un exprime la cause, l'autre le but.
\item \zh{看病} (kànbìng) : littéralement « regarder la maladie », c'est-à-dire « consulter (un médecin), se faire soigner ».
\end{itemize}

\courssub{Repérage des informations clés}

Le texte est court mais très bien organisé : chaque phrase apporte un type d'information précis. Repérons-les.

\begin{center}
\begin{tabularx}{\linewidth}{|X|X|X|}
\hline
\rowcolor{popblueL} Information clé & Phrase du texte (chinois + pinyin) & Traduction \\
\hline
Définition 定义 & 疟疾是一种很常见的疾病。 Nüèji shì yì zhǒng hěn chángjiàn de jíbìng. & Le paludisme est une maladie très courante. \\
\hline
Transmission 传播 & 蚊子把疟疾病毒传给人。 Wénzi bǎ nüèji bìngdú chuán gěi rén. & Le moustique transmet l'agent du paludisme à l'homme. \\
\hline
Symptômes 症状 & 发烧、头疼、全身没有力气 fāshāo, tóuténg, quánshēn méiyǒu lìqi & fièvre, maux de tête, grande fatigue \\
\hline
Conséquences 后果 & 很多孩子因为疟疾死亡。 hěn duō háizi yīnwèi nüèji sǐwáng & beaucoup d'enfants meurent du paludisme \\
\hline
Prévention 预防 & 用蚊帐，保持环境干净，及时去医院看病 yòng wénzhàng, bǎochí huánjìng gānjìng, jíshí qù yīyuàn kànbìng & moustiquaires, environnement propre, consultation rapide \\
\hline
\end{tabularx}
\end{center}

\begin{popfigure}
\begin{tikzpicture}[mindmap, grow cyclic,
  every node/.style={concept, align=center, font=\small},
  root concept/.append style={concept color=popblue, font=\large\bfseries},
  level 1/.append style={level distance=3.2cm, sibling angle=72},
  level 1 concept/.append style={concept color=poporange}]
\node[root concept, align=center]{疟疾\\ nüèji\\ \footnotesize paludisme}
  child { node[align=center]{定义\\ dìngyì\\ \footnotesize 常见的疾病} }
  child { node[align=center]{传播\\ chuánbō\\ \footnotesize 蚊子传给人} }
  child { node[align=center]{症状\\ zhèngzhuàng\\ \footnotesize 发烧、头疼} }
  child { node[align=center]{后果\\ hòuguǒ\\ \footnotesize 孩子死亡} }
  child { node[align=center]{预防\\ yùfáng\\ \footnotesize 蚊帐、医院} };
\end{tikzpicture}
\legende{Carte mentale du texte : les cinq informations clés autour du mot 疟疾 (nüèji).}
\end{popfigure}

\begin{popfigure}
\begin{tikzpicture}[node distance=0.9cm,
  box/.style={draw=popblue, fill=popblueL, rounded corners, thick, align=center, minimum width=2.2cm, minimum height=1cm, font=\small},
  fleche/.style={-{Stealth[length=3mm]}, very thick, poporange}]
\node[box, align=center] (mou){蚊子\\ wénzi\\ \footnotesize moustique};
\node[box, right=of mou, align=center] (vir){病毒\\ bìngdú\\ \footnotesize agent pathogène};
\node[box, right=of vir, align=center] (hom){人\\ rén\\ \footnotesize l'homme};
\node[box, right=of hom, align=center] (fie){发烧\\ fāshāo\\ \footnotesize fièvre};
\draw[fleche] (mou) -- node[above, font=\footnotesize]{传播 chuánbō} (vir);
\draw[fleche] (vir) -- node[above, font=\footnotesize]{传给 chuán gěi} (hom);
\draw[fleche] (hom) -- node[above, font=\footnotesize]{得了 dé le} (fie);
\end{tikzpicture}
\legende{Chaîne de transmission du paludisme : le moustique transmet l'agent pathogène à l'homme, qui développe alors la fièvre.}
\end{popfigure}

\begin{exoresolu}[Vérification de la compréhension globale]
Énoncé : sans regarder la traduction, répondez en français aux trois questions suivantes, en vous appuyant uniquement sur le texte chinois.
\begin{enumerate}
\item Qu'est-ce qui transmet le paludisme à l'homme ?
\item Citez deux symptômes du paludisme mentionnés dans le texte.
\item Citez deux moyens de prévention proposés par le texte.
\end{enumerate}
\tcblower
Solution détaillée :
\begin{enumerate}
\item C'est le \cle{moustique} (\zh{蚊子}, wénzi) qui transmet le paludisme : \zh{蚊子把疟疾病毒传给人。}
\item Deux symptômes possibles parmi trois : la \cle{fièvre} (\zh{发烧}, fāshāo), les \cle{maux de tête} (\zh{头疼}, tóuténg), la \cle{fatigue générale} (\zh{全身没有力气}, quánshēn méiyǒu lìqi).
\item Deux moyens de prévention possibles parmi trois : utiliser des \cle{moustiquaires} (\zh{用蚊帐}, yòng wénzhàng), garder l'\cle{environnement propre} (\zh{保持环境干净}, bǎochí huánjìng gānjìng), aller \cle{consulter à l'hôpital à temps} (\zh{及时去医院看病}, jíshí qù yīyuàn kànbìng).
\end{enumerate}
Stratégie : pour chaque question, repérez d'abord le mot-clé chinois correspondant (\zh{蚊子} pour la transmission, \zh{发烧} pour les symptômes, \zh{应该} pour les conseils), puis lisez la phrase entière qui l'entoure.
\end{exoresolu}

\begin{exercice}[À vous de jouer]
Replacez les cinq phrases du texte dans l'ordre logique (définition → transmission → symptômes → conséquences → prévention), puis recopiez chaque phrase avec son pinyin :
\begin{enumerate}
\item 为了预防疟疾，我们应该用蚊帐。
\item 疟疾是一种很常见的疾病。
\item 得了疟疾的人会发烧、头疼。
\item 蚊子把疟疾病毒传给人。
\item 在非洲，很多孩子因为疟疾死亡。
\end{enumerate}
\end{exercice}

\begin{corrige}[Exercice « À vous de jouer »]
L'ordre correct est : 2 → 4 → 3 → 5 → 1.
\begin{enumerate}
\item \zh{疟疾是一种很常见的疾病。} \emph{Nüèji shì yì zhǒng hěn chángjiàn de jíbìng.} (définition)
\item \zh{蚊子把疟疾病毒传给人。} \emph{Wénzi bǎ nüèji bìngdú chuán gěi rén.} (transmission)
\item \zh{得了疟疾的人会发烧、头疼。} \emph{Dé le nüèji de rén huì fāshāo, tóuténg.} (symptômes)
\item \zh{在非洲，很多孩子因为疟疾死亡。} \emph{Zài Fēizhōu, hěn duō háizi yīnwèi nüèji sǐwáng.} (conséquences)
\item \zh{为了预防疟疾，我们应该用蚊帐。} \emph{Wèile yùfáng nüèji, wǒmen yīnggāi yòng wénzhàng.} (prévention)
\end{enumerate}
\end{corrige}

\begin{savaistu}[La Chine, un pays désormais sans paludisme]
Le paludisme a longtemps sévi en Chine, comme au Cameroun. Mais en 2021, l'Organisation mondiale de la santé (OMS) a officiellement certifié la Chine \emph{exempte de paludisme}, après des décennies de lutte : utilisation massive des moustiquaires imprégnées, assainissement de l'environnement et traitements à base d'\cle{artémisinine}. Cette molécule a été découverte par la chercheuse chinoise \zh{屠呦呦} (Tú Yōuyōu), qui a reçu le prix Nobel de médecine en 2015 pour cette découverte qui a sauvé des millions de vies, notamment en Afrique. Les médicaments à base d'artémisinine que l'on utilise dans les hôpitaux camerounais sont directement issus de ses travaux : un beau pont scientifique entre la Chine et le Cameroun !
\end{savaistu}

\begin{aretenir}[Ce qu'il faut retenir de cette première étape]
\begin{itemize}
\item Le titre \zh{疟疾} (nüèji) annonce le thème ; la clé de la maladie \zh{疒} aide à deviner le sens des caractères inconnus.
\item Un texte chinois de santé se lit en quatre étapes : titre → mots connus → devinettes par le contexte et les radicaux → vérification au glossaire.
\item Le texte suit un plan en cinq points : \textbf{définition, transmission, symptômes, conséquences, prévention} — le même plan que votre brainstorming en français.
\item Attention aux tons (\emph{nüèji}, \emph{fāshāo}, \emph{yīnwèi}) et aux deux lectures de \zh{得} : \emph{dé} (attraper une maladie) et \emph{de} (dans \zh{觉得}, penser).
\item La construction en \zh{把} (Sujet + 把 + objet + verbe + complément) place l'objet avant le verbe : ne la traduisez jamais mot à mot.
\end{itemize}
\end{aretenir}

\coursec{Glossaire du thème}

Dans la section précédente, nous avons lu ensemble le texte support sur le paludisme au Cameroun et relevé les mots nouveaux qui y apparaissaient. Il est temps maintenant de nous arrêter sur chacun de ces mots : les observer, les prononcer correctement, comprendre leur construction et apprendre à les employer dans des phrases. Un bon glossaire ne se récite pas : il s'explore, se classe et se manipule. C'est exactement ce que nous allons faire.

\courssub{Présentation du glossaire}

Voici le tableau complet du vocabulaire de cette leçon. Pour chaque mot, tu trouveras les caractères, le pinyin avec les tons, la nature grammaticale et la traduction française. Lis chaque ligne à voix haute avant de poursuivre.

\begin{center}
\begin{tabularx}{\linewidth}{|X|X|X|X|}
\hline
\rowcolor{popblueL} Caractères & Pinyin & Nature & Traduction \\
\hline
疟疾 & nüèji & n. & paludisme \\
\hline
疾病 & jíbìng & n. & maladie \\
\hline
病毒 & bìngdú & n. & virus \\
\hline
蚊子 & wénzi & n. & moustique \\
\hline
蚊帐 & wénzhàng & n. & moustiquaire \\
\hline
医院 & yīyuàn & n. & hôpital \\
\hline
环境 & huánjìng & n. & environnement \\
\hline
传染 & chuánrǎn & v. & transmettre, contaminer \\
\hline
得 & dé & v. & attraper (une maladie) \\
\hline
发烧 & fāshāo & v. & avoir de la fièvre \\
\hline
预防 & yùfáng & v. & prévenir \\
\hline
保持 & bǎochí & v. & maintenir \\
\hline
死亡 & sǐwáng & v. & mourir \\
\hline
常见 & chángjiàn & adj. & courant, fréquent \\
\hline
干净 & gānjìng & adj. & propre \\
\hline
及时 & jíshí & adv. & à temps \\
\hline
\end{tabularx}
\end{center}

\begin{center}
\begin{tabularx}{\linewidth}{|X|X|X|X|}
\hline
\rowcolor{popblueL} Caractères & Pinyin & Nature & Traduction \\
\hline
全身 & quánshēn & expr. & tout le corps \\
\hline
没有力气 & méiyǒu lìqi & expr. & être sans force, être épuisé \\
\hline
因为 & yīnwèi & conj. & parce que \\
\hline
为了 & wèile & prép. & afin de, pour \\
\hline
应该 & yīnggāi & v. mod. & devoir (conseil) \\
\hline
\end{tabularx}
\end{center}

La figure ci-dessous reprend ces seize mots sous forme de tableau visuel à quatre colonnes, comme on le ferait sur une affiche de classe.

\begin{popfigure}
\begin{tikzpicture}[scale=0.62, every node/.style={transform shape}]
\fill[popblue] (0,0) rectangle (12,1);
\node[text=white, font=\bfseries] at (1.5,0.5) {汉字};
\node[text=white, font=\bfseries] at (4.5,0.5) {拼音};
\node[text=white, font=\bfseries] at (7.5,0.5) {词性};
\node[text=white, font=\bfseries] at (10.2,0.5) {法语};
\foreach \i/\zi/\py/\cx/\fr in {
1/疟疾/nüèji/n./paludisme,
2/疾病/jíbìng/n./maladie,
3/病毒/bìngdú/n./virus,
4/蚊子/wénzi/n./moustique,
5/蚊帐/wénzhàng/n./moustiquaire,
6/医院/yīyuàn/n./hôpital,
7/环境/huánjìng/n./environnement,
8/传染/chuánrǎn/v./transmettre,
9/得/dé/v./attraper,
10/发烧/fāshāo/v./avoir de la fièvre,
11/预防/yùfáng/v./prévenir,
12/保持/bǎochí/v./maintenir,
13/死亡/sǐwáng/v./mourir,
14/常见/chángjiàn/adj./courant,
15/干净/gānjìng/adj./propre,
16/及时/jíshí/adv./à temps}
{
\ifodd\i \fill[popcream] (0,-\i) rectangle (12,1-\i);
\else \fill[popblueL!50] (0,-\i) rectangle (12,1-\i); \fi
\node at (1.5,0.5-\i) {\zi};
\node at (4.5,0.5-\i) {\py};
\node at (7.5,0.5-\i) {\cx};
\node at (10.2,0.5-\i) {\fr};
}
\draw[popdark, thick] (0,0) rectangle (12,-16);
\foreach \x in {3,6,9} \draw[popline] (\x,0) -- (\x,-16);
\foreach \i in {1,...,15} \draw[popline] (0,-\i) -- (12,-\i);
\end{tikzpicture}
\legende{Tableau visuel du glossaire : seize mots, quatre colonnes (caractères, pinyin, nature, traduction).}
\end{popfigure}

\courssub{Les noms du thème}

\begin{definition}[疾病 jíbìng : la « maladie » en général]
\zh{疾病} \emph{jíbìng} est le terme général et soutenu pour dire « maladie ». Il est formé de deux caractères synonymes : \zh{疾} \emph{jí} (maladie) et \zh{病} \emph{bìng} (maladie). Dans la conversation quotidienne, on utilise souvent \zh{病} seul : \zh{我病了。} \emph{Wǒ bìng le.} « Je suis malade. » Mais dans un texte de santé publique, un article ou un exposé, on préfère \zh{疾病} : \zh{预防疾病} « prévenir les maladies ».
\end{definition}

Observons comment ces noms se construisent ; le chinois adore composer des mots à partir de deux caractères :

\begin{itemize}
\item \zh{疟疾} \emph{nüèji} « paludisme » : les deux caractères portent la clé de la maladie \zh{疒}. Ce radical est un indice précieux : un caractère avec \zh{疒} parle presque toujours de maladie (\zh{病}, \zh{疾}, \zh{疟}, \zh{疼}).
\item \zh{病毒} \emph{bìngdú} « virus » : littéralement « maladie + poison ». Logique et facile à retenir.
\item \zh{蚊子} \emph{wénzi} « moustique » : le suffixe \zh{子} (ton neutre) sert à former des noms concrets, comme dans \zh{孩子} « enfant » ou \zh{桌子} « table ».
\item \zh{蚊帐} \emph{wénzhàng} « moustiquaire » : littéralement « moustique + rideau ». Un rideau contre les moustiques : l'image est parlante.
\item \zh{医院} \emph{yīyuàn} « hôpital » : \zh{医} « médecine » + \zh{院} « institution, cour ». On retrouve \zh{院} dans \zh{学院} « institut ».
\item \zh{环境} \emph{huánjìng} « environnement » : \zh{环} « anneau, entourer » + \zh{境} « territoire, cadre » : ce qui nous entoure.
\end{itemize}

\begin{exemplebox}[Les noms en contexte]
\zh{疟疾是一种常见的疾病。} \emph{Nüèji shì yì zhǒng chángjiàn de jíbìng.} « Le paludisme est une maladie courante. »\\
\zh{蚊子传染这种疾病。} \emph{Wénzi chuánrǎn zhè zhǒng jíbìng.} « Le moustique transmet cette maladie. »\\
\zh{他去了医院。} \emph{Tā qù le yīyuàn.} « Il est allé à l'hôpital. »\\
\zh{我们要保持环境干净。} \emph{Wǒmen yào bǎochí huánjìng gānjìng.} « Nous devons maintenir l'environnement propre. »
\end{exemplebox}

\courssub{Les verbes et expressions verbales}

\begin{propriete}[得 dé + maladie : « attraper »]
Structure : \textbf{Sujet + 得 + (了) + nom de maladie}. Le verbe \zh{得} \emph{dé} suivi d'un nom de maladie signifie « attraper, contracter » : \zh{得疟疾} « attraper le paludisme », \zh{得病} « tomber malade ». Attention : ce \zh{得} (dé) n'a rien à voir avec le mot-outil \zh{得} (de) des compléments de degré ; c'est ici un vrai verbe.
\end{propriete}

\begin{exemplebox}[Verbes en action]
\zh{他得了疟疾。} \emph{Tā dé le nüèji.} « Il a attrapé le paludisme. »\\
\zh{我们应该预防疾病。} \emph{Wǒmen yīnggāi yùfáng jíbìng.} « Nous devons prévenir les maladies. »\\
\zh{孩子发烧了，全身没有力气。} \emph{Háizi fāshāo le, quánshēn méiyǒu lìqi.} « L'enfant a de la fièvre, tout son corps est sans force. »\\
\zh{很多病人死亡了。} \emph{Hěn duō bìngrén sǐwáng le.} « Beaucoup de malades sont morts. » (registre soutenu ; à l'oral on dirait plutôt \zh{死了})
\end{exemplebox}

\begin{attention}[发烧 : un verbe à complément incorporé]
\zh{发烧} \emph{fāshāo} est un verbe composé inséparable dans l'usage courant : \zh{发} « émettre » + \zh{烧} « brûlure, chaleur ». On dit \zh{我发烧了。} \emph{Wǒ fāshāo le.} « J'ai de la fièvre. » et JAMAIS \zh{我很发烧} : \zh{发烧} est un verbe, pas un adjectif, donc pas de \zh{很} devant. Pour insister sur le degré, on dit \zh{我发烧发得很厉害。} ou plus simplement \zh{我发高烧。} « J'ai une forte fièvre. »
\end{attention}

\courssub{Adjectifs, adverbe et connecteurs}

\begin{itemize}
\item \zh{常见} \emph{chángjiàn} (adj.) « courant » : \zh{常见病} « maladie courante ».
\item \zh{干净} \emph{gānjìng} (adj.) « propre » ; contraire : \zh{脏} \emph{zāng} « sale ».
\item \zh{及时} \emph{jíshí} (adv.) « à temps » : \zh{及时去医院} « aller à l'hôpital à temps ». Il se place avant le verbe.
\item \zh{因为} \emph{yīnwèi} « parce que » : \zh{因为蚊子传染疟疾，所以我们要小心。} « Parce que le moustique transmet le paludisme, nous devons être prudents. »
\item \zh{为了} \emph{wèile} « afin de » + verbe : \zh{为了预防疟疾，我们睡觉用蚊帐。} « Afin de prévenir le paludisme, nous dormons sous une moustiquaire. »
\item \zh{应该} \emph{yīnggāi} « devoir » (conseil) : \zh{你应该保持房间干净。} « Tu devrais garder ta chambre propre. »
\end{itemize}

\begin{methode}[Mémoriser par familles de mots]
\begin{enumerate}
\item Regroupe les mots autour d'un caractère commun : \zh{蚊} → \zh{蚊子} (moustique), \zh{蚊帐} (moustiquaire) ; \zh{病} → \zh{病毒} (virus), \zh{疾病} (maladie), \zh{看病} (consulter un médecin).
\item Pour chaque famille, fabrique une phrase simple : \zh{蚊子传染疾病。}
\item Récite les familles à voix haute, en frappant les tons dans tes mains.
\item Le lendemain, teste-toi : cache la colonne des caractères et retrouve-les à partir du français.
\end{enumerate}
\end{methode}

\courssub{Entraînement à la prononciation}

Lisons le glossaire en chœur, trois fois, en insistant sur les tons difficiles :

\begin{itemize}
\item \zh{疟} \emph{nüè} : quatrième ton, descendant et sec. La voyelle \emph{üe} se prononce les lèvres arrondies.
\item \zh{疾} \emph{jí} : deuxième ton, montant, comme une question en français.
\item \zh{传染} \emph{chuánrǎn} : deuxième ton puis troisième ton ; bien creuser le troisième ton.
\item \zh{死亡} \emph{sǐwáng} : troisième ton puis deuxième ton ; ne pas confondre avec \zh{希望} \emph{xīwàng} « espérer » !
\item \zh{及时} \emph{jíshí} : deux deuxièmes tons de suite, garder l'énergie montante.
\end{itemize}

\begin{exoresolu}[Mini-dictée de huit mots]
Énoncé : le professeur dicte lentement huit mots du glossaire (par exemple : 疟疾, 蚊子, 医院, 发烧, 预防, 干净, 及时, 环境). Écris chaque mot en caractères, avec son pinyin et sa traduction.
\tcblower
Solution détaillée :
\begin{enumerate}
\item \zh{疟疾} nüèji : paludisme (attention au quatrième ton de nüè).
\item \zh{蚊子} wénzi : moustique (zi au ton neutre).
\item \zh{医院} yīyuàn : hôpital.
\item \zh{发烧} fāshāo : avoir de la fièvre (verbe, pas d'adjectif).
\item \zh{预防} yùfáng : prévenir.
\item \zh{干净} gānjìng : propre.
\item \zh{及时} jíshí : à temps.
\item \zh{环境} huánjìng : environnement.
\end{enumerate}
Barème conseillé : un point par mot (0,5 pour les caractères, 0,25 pour le pinyin avec tons corrects, 0,25 pour la traduction).
\end{exoresolu}

\begin{exercice}[À toi de jouer]
Complète les phrases avec un mot du glossaire, puis traduis-les en français :
\begin{enumerate}
\item \zh{为了＿＿疟疾，我们睡觉用蚊帐。}
\item \zh{他＿＿了，全身没有力气，应该及时去医院。}
\item \zh{保持环境＿＿可以预防很多疾病。}
\end{enumerate}
\end{exercice}

\begin{corrige}[Exercice « À toi de jouer »]
1. \zh{预防} : « Afin de prévenir le paludisme, nous dormons sous une moustiquaire. »\\
2. \zh{发烧} : « Il a de la fièvre, tout son corps est sans force, il devrait aller à l'hôpital à temps. »\\
3. \zh{干净} : « Maintenir l'environnement propre permet de prévenir beaucoup de maladies. »
\end{corrige}

\begin{aretenir}[L'essentiel du glossaire]
Seize mots à connaître par cœur : sept noms, six verbes, deux adjectifs, un adverbe et cinq expressions. Trois réflexes : \zh{得 + maladie} pour « attraper » ; \zh{发烧} est un verbe (jamais \zh{很发烧}) ; \zh{为了 + verbe} pour le but, \zh{因为 + phrase} pour la cause. Classe les mots par familles (\zh{蚊}, \zh{病}) pour les mémoriser durablement.
\end{aretenir}

\coursec{Clés (radicaux) et ordre des traits}

Dans la section précédente, nous avons découvert le glossaire du paludisme : \zh{疟疾} (nüèji, paludisme), \zh{蚊子} (wénzi, moustique), \zh{蚊帐} (wénzhàng, moustiquaire), \zh{发烧} (fāshāo, avoir de la fièvre), \zh{疾病} (jíbìng, maladie). Observez attentivement ces caractères : ne remarquez-vous pas que certains morceaux reviennent sans cesse ? \zh{疟}, \zh{疾} et \zh{病} commencent tous par le même élément en haut à gauche ; \zh{蚊} contient un petit « insecte » à gauche. Ces éléments répétés ne sont pas un hasard : ce sont les \cle{clés} (ou \cle{radicaux}, \zh{部首} bùshǒu), et ils vont devenir vos meilleurs alliés pour lire, écrire et chercher les caractères chinois.

\courssub{Le rôle des radicaux : deviner le sens et utiliser le dictionnaire}

\begin{definition}[Qu'est-ce qu'une clé (radical) ?]
Une \cle{clé} (\zh{部首} bùshǒu) est un élément graphique qui se répète dans une famille de caractères et qui donne un indice sur le \textbf{sens général} du caractère. Le chinois possède environ 214 clés traditionnelles. Connaître les clés les plus fréquentes, c'est posséder une carte d'identité approximative de milliers de caractères.
\end{definition}

Prenons un exemple concret. Vous rencontrez le caractère \zh{疼} (téng, avoir mal) pour la première fois. Vous ne le connaissez pas, mais vous reconnaissez en haut à gauche la clé de la maladie \zh{疒}. Vous pouvez donc deviner : « ce caractère a un rapport avec la maladie ou la douleur ». C'est exact ! De même, si vous voyez un caractère inconnu contenant \zh{虫} (chóng, insecte), vous devinez qu'il s'agit probablement d'un insecte ou d'une petite bête.

\begin{aretenir}[Deux fonctions essentielles des clés]
\begin{enumerate}
\item \textbf{Deviner le sens} : la clé indique la famille sémantique (maladie, insecte, feu, eau, main, bouche...).
\item \textbf{Chercher dans un dictionnaire} : les dictionnaires chinois classent les caractères par clé, puis par nombre de traits restants.
\end{enumerate}
\end{aretenir}

\begin{methode}[Retrouver un mot inconnu dans le dictionnaire chinois]
\begin{enumerate}
\item \textbf{Repérer la clé} du caractère (souvent à gauche, en haut, ou en position enveloppante). Exemple : dans \zh{疾}, la clé est \zh{疒}.
\item \textbf{Compter les traits restants}, c'est-à-dire les traits qui ne font pas partie de la clé. Dans \zh{疾}, la clé \zh{疒} a 5 traits ; il reste 5 traits (\zh{矢}).
\item \textbf{Ouvrir le dictionnaire à la table des clés} et trouver la clé \zh{疒} (classée à 5 traits).
\item \textbf{Chercher dans la liste} des caractères de cette clé ceux qui ont 5 traits supplémentaires : vous trouvez \zh{疾} (jí).
\item \textbf{Lire la notice} : prononciation pinyin, sens, mots composés (\zh{疾病} jíbìng, maladie).
\end{enumerate}
Cette méthode fonctionne même sans connaître la prononciation du caractère : c'est sa grande force.
\end{methode}

\courssub{La clé 疒 (nè) : la famille de la maladie}

La clé \zh{疒} se prononce \emph{nè} et signifie « maladie ». Elle occupe toujours la partie haute et gauche du caractère, comme une petite couverture qui enveloppe le reste. Tous les caractères qui la contiennent appartiennent au domaine de la maladie, de la douleur ou des symptômes.

\begin{center}
\begin{tabularx}{\linewidth}{|X|X|X|X|}
\hline
\rowcolor{popblueL} Caractère & Pinyin & Clé & Traduction \\
\hline
疟 & nüè & 疒 & paludisme (dans 疟疾) \\
\hline
疾 & jí & 疒 & maladie (dans 疾病) \\
\hline
病 & bìng & 疒 & maladie ; être malade \\
\hline
疼 & téng & 疒 & douloureux, avoir mal \\
\hline
\end{tabularx}
\end{center}

Exemples en contexte :
\begin{itemize}
\item \zh{疟疾是喀麦隆常见的疾病。} \emph{Nüèji shì Kāmàilóng chángjiàn de jíbìng.} « Le paludisme est une maladie fréquente au Cameroun. »
\item \zh{他病了，头很疼。} \emph{Tā bìng le, tóu hěn téng.} « Il est tombé malade, il a très mal à la tête. »
\end{itemize}

Remarquez dans le deuxième exemple l'emploi de \zh{病} comme \textbf{verbe} : \zh{病了} (bìng le) signifie « être tombé malade », avec la particule \zh{了} qui marque le changement d'état. En chinois, un même caractère peut être nom ou verbe selon sa place dans la phrase.

\begin{propriete}[La clé 疒 se trace toujours en premier]
Dans tous les caractères contenant \zh{疒} (\zh{疟}, \zh{疾}, \zh{病}, \zh{疼}), la clé se trace \textbf{toujours en premier}, et elle \textbf{enveloppe} le reste du caractère. L'ordre des traits de la clé est fixe : \cle{点} (diǎn, le point en haut) → \cle{横} (héng, l'horizontale) → \cle{撇} (piě, la grande oblique descendante à gauche) → \cle{点} (diǎn, le petit point intérieur) → \cle{提} (tí, la petite montée oblique). On trace ensuite seulement l'intérieur du caractère. C'est la règle générale de l'écriture chinoise : \textbf{de l'extérieur vers l'intérieur} pour les caractères enveloppants.
\end{propriete}

\begin{savaistu}[Une clé qui dit l'histoire de la médecine]
Plus de 200 caractères chinois contiennent la clé \zh{疒} : \zh{疫} (yì, épidémie), \zh{症} (zhèng, symptôme), \zh{癌} (ái, cancer), \zh{疲} (pí, fatigue)... Cette abondance témoigne de l'importance très ancienne de la médecine dans la civilisation chinoise : la médecine traditionnelle chinoise (\zh{中医} zhōngyī) est pratiquée depuis plus de deux mille ans. Saviez-vous que l'artémisinine, un des traitements modernes du paludisme, a été découverte par la chercheuse chinoise Tu Youyou (\zh{屠呦呦} Tú Yōuyōu) en étudiant des textes de médecine traditionnelle ? Elle a reçu le prix Nobel de médecine en 2015 — une belle rencontre entre la Chine ancienne et la lutte contre une maladie qui touche aussi le Cameroun.
\end{savaistu}

\courssub{La clé 虫 (chóng) : la famille des insectes}

La clé \zh{虫} (chóng) signifie « insecte » ou « petite créature ». Elle se place généralement \textbf{à gauche} du caractère. Le mot \zh{蚊子} (wénzi, moustique) contient cette clé : logique, le moustique est un insecte !

\begin{center}
\begin{tabularx}{\linewidth}{|X|X|X|X|}
\hline
\rowcolor{popblueL} Caractère & Pinyin & Clé & Traduction \\
\hline
蚊 & wén & 虫 & moustique (dans 蚊子) \\
\hline
蚂 & mǎ & 虫 & fourmi (dans 蚂蚁) \\
\hline
蚁 & yǐ & 虫 & fourmi (dans 蚂蚁) \\
\hline
\end{tabularx}
\end{center}

Exemple : \zh{蚊子传播疟疾。} \emph{Wénzi chuánbō nüèji.} « Le moustique transmet le paludisme. »

Notez le verbe \zh{传播} (chuánbō, transmettre, propager), très utile pour parler des maladies : \zh{传播疾病} (chuánbō jíbìng, transmettre une maladie).

\courssub{La clé 巾 (jīn) dans 帐 et la clé 火 (huǒ) dans 烧}

\textbf{La clé 巾 (jīn, « tissu »).} Le caractère \zh{帐} (zhàng), dans \zh{蚊帐} (wénzhàng, moustiquaire), contient à gauche la clé \zh{巾}, qui représente un morceau de tissu suspendu. Le lien est évident : une moustiquaire est faite de tissu ! Autre caractère de la même famille : \zh{布} (bù, tissu, étoffe).

\textbf{La clé 火 (huǒ, « feu »).} Le caractère \zh{烧} (shāo, brûler), dans \zh{发烧} (fāshāo, avoir de la fièvre), contient à gauche la clé du feu. Quelle image parlante : quand on a de la fièvre, le corps « brûle » ! Le chinois dit littéralement « émettre de la brûlure » pour dire « avoir de la fièvre ». Autres caractères avec \zh{火} : \zh{热} (rè, chaud), \zh{灯} (dēng, lampe).

Exemple : \zh{孩子发烧了，妈妈带他去医院。} \emph{Háizi fāshāo le, māma dài tā qù yīyuàn.} « L'enfant a de la fièvre, sa mère l'emmène à l'hôpital. »

\begin{attention}[Ne pas confondre 帐 et 账 : même pinyin, sens différent !]
\zh{帐} (zhàng, moustiquaire, tente) et \zh{账} (zhàng, compte, facture) se prononcent \textbf{exactement pareil} : \emph{zhàng}. Seule la clé les distingue :
\begin{itemize}
\item \zh{帐} = clé \zh{巾} (tissu) → \zh{蚊帐} (wénzhàng, moustiquaire) ;
\item \zh{账} = clé \zh{贝} (bèi, coquillage, ancienne monnaie) → \zh{账户} (zhànghù, compte en banque).
\end{itemize}
Erreur fréquente des francophones : écrire \zh{蚊账} au lieu de \zh{蚊帐}. Retenez l'astuce : la moustiquaire est en \textbf{tissu} (\zh{巾}), l'argent se compte avec des \textbf{coquillages} (\zh{贝}). En chinois, le sens est dans les yeux, pas dans les oreilles !
\end{attention}

\courssub{Ordre des traits : règles générales et application}

\begin{propriete}[Les six règles d'or de l'ordre des traits]
\begin{enumerate}
\item De \textbf{haut en bas} : \zh{三} (sān, trois) se trace de la barre du haut vers celle du bas.
\item De \textbf{gauche à droite} : \zh{蚊} se trace avec \zh{虫} d'abord, puis \zh{文}.
\item L'\textbf{horizontale avant la verticale} : \zh{十} (shí, dix) : d'abord \cle{横}, puis \cle{竖}.
\item L'\textbf{extérieur avant l'intérieur} : \zh{疟} : d'abord \zh{疒}, puis l'intérieur.
\item On \textbf{ferme en dernier} : dans \zh{回}, on ferme le bas du cadre à la fin.
\item La \textbf{verticale centrale avant les côtés} dans les caractères symétriques : \zh{小} (xiǎo, petit).
\end{enumerate}
\end{propriete}

\textbf{Le caractère 疟 (nüè) : 8 traits.} On trace d'abord la clé \zh{疒} en 5 traits : 1) le point \cle{点} en haut ; 2) l'horizontale \cle{横} ; 3) la grande oblique \cle{撇} qui descend à gauche ; 4) le petit point \cle{点} à l'intérieur gauche ; 5) la montée \cle{提}. Puis l'intérieur : 6) l'horizontale supérieure ; 7) la verticale courbée qui descend ; 8) l'horizontale du bas qui ferme.

\begin{popfigure}
\begin{tikzpicture}[scale=1.1]
\draw[popline, dashed] (0,0) rectangle (4,4);
\draw[popline, dashed] (2,0) -- (2,4);
\draw[popline, dashed] (0,2) -- (4,2);
% trait 1 : point
\draw[line width=2.2pt, popink] (2.0,3.75) -- (2.15,3.45);
\node[poporange, font=\small\bfseries] at (1.55,3.75) {1};
% trait 2 : heng
\draw[line width=2.2pt, popink] (1.0,3.15) -- (3.6,3.15);
\node[poporange, font=\small\bfseries] at (3.75,3.35) {2};
% trait 3 : pie
\draw[line width=2.2pt, popink] (1.0,3.15) .. controls (0.85,2.3) .. (0.55,0.7);
\node[poporange, font=\small\bfseries] at (0.3,1.9) {3};
% trait 4 : point interieur
\draw[line width=2.2pt, popink] (1.25,2.35) -- (1.45,2.15);
\node[poporange, font=\small\bfseries] at (1.0,2.05) {4};
% trait 5 : ti
\draw[line width=2.2pt, popink] (1.2,1.55) -- (1.6,1.85);
\node[poporange, font=\small\bfseries] at (1.0,1.35) {5};
% trait 6 : heng interieur
\draw[line width=2.2pt, popblue] (1.75,2.5) -- (3.35,2.5);
\node[poporange, font=\small\bfseries] at (3.5,2.7) {6};
% trait 7 : verticale courbee
\draw[line width=2.2pt, popblue] (2.0,2.5) .. controls (1.95,1.7) .. (2.05,0.85);
\node[poporange, font=\small\bfseries] at (1.75,1.7) {7};
% trait 8 : heng bas
\draw[line width=2.2pt, popblue] (2.05,0.85) -- (3.35,0.85);
\node[poporange, font=\small\bfseries] at (3.5,0.65) {8};
\end{tikzpicture}
\legende{Le caractère 疟 (nüè) en 8 traits : la clé 疒 (traits 1 à 5, en noir) se trace d'abord et enveloppe l'intérieur (traits 6 à 8, en bleu).}
\end{popfigure}

\textbf{Le caractère 疾 (jí) : 10 traits.} Même logique : la clé \zh{疒} d'abord (traits 1 à 5, identiques à ceux de \zh{疟}), puis la partie \zh{矢} (shǐ, flèche) en 5 traits : 6) le petit \cle{撇} en haut ; 7) l'horizontale courte ; 8) l'horizontale longue ; 9) le \cle{撇} qui descend à gauche ; 10) le \cle{捺} (nà, l'appuyé) qui descend à droite.

\textbf{Le caractère 蚊 (wén) : 10 traits.} Ce caractère se trace de gauche à droite. D'abord la clé \zh{虫} en 6 traits : 1) la verticale ; 2) l'horizontale-pliée \cle{横折} (héngzhé) qui forme le haut et la droite du « corps » ; 3) l'horizontale du milieu ; 4) la verticale centrale ; 5) la montée \cle{提} ; 6) le point. Ensuite \zh{文} (wén) en 4 traits : 7) le point en haut ; 8) l'horizontale ; 9) le \cle{撇} ; 10) le \cle{捺} qui s'ouvre à droite.

\begin{popfigure}
\begin{tikzpicture}[
  grow=right,
  level 1/.style={sibling distance=1.6cm, level distance=3.2cm},
  every node/.style={draw, rounded corners, align=center, font=\small},
  edge from parent/.style={draw=popmuted, -{Stealth}, thick}
]
\node[fill=poporangeL, align=center]{疒 (nè)\\ maladie}
  child { node[fill=popblueL, align=center]{疟 nüè\\ paludisme} }
  child { node[fill=popblueL, align=center]{疾 jí\\ maladie} }
  child { node[fill=popblueL, align=center]{病 bìng\\ maladie} }
  child { node[fill=popblueL, align=center]{疼 téng\\ douleur} };
\begin{scope}[yshift=-4.6cm]
\node[fill=popgreenL, align=center]{虫 (chóng)\\ insecte}
  child { node[fill=popblueL, align=center]{蚊 wén\\ moustique} }
  child { node[fill=popblueL, align=center]{蚂 mǎ\\ fourmi} }
  child { node[fill=popblueL, align=center]{蚁 yǐ\\ fourmi} };
\end{scope}
\end{tikzpicture}
\legende{Arbre des familles de radicaux : la clé 疒 regroupe les caractères de la maladie, la clé 虫 ceux des insectes.}
\end{popfigure}

\begin{exoresolu}[Identifier la clé et compter les traits]
Énoncé. Pour chaque caractère, indique la clé, le nombre de traits de la clé et le nombre de traits restants : \zh{病}, \zh{烧}, \zh{帐}.
\tcblower
Solution détaillée.
\begin{itemize}
\item \zh{病} (bìng) : la clé est \zh{疒} (5 traits). La partie restante \zh{丙} compte 5 traits. Total : 10 traits. Famille : la maladie.
\item \zh{烧} (shāo) : la clé est \zh{火} (4 traits). La partie restante \zh{尧} compte 6 traits. Total : 10 traits. Famille : le feu.
\item \zh{帐} (zhàng) : la clé est \zh{巾} (3 traits). La partie restante \zh{长} compte 4 traits. Total : 7 traits. Famille : le tissu.
\end{itemize}
Vérification au dictionnaire : pour trouver \zh{病}, on ouvre la table des clés à \zh{疒}, puis on cherche parmi les caractères à 5 traits supplémentaires.
\end{exoresolu}

\begin{exercice}[Tracé sur cahier]
Sur ton cahier d'écriture, trace chaque caractère \textbf{trois fois}, en respectant l'ordre des traits appris en classe et en prononçant le pinyin à voix haute à chaque tracé : \zh{疟} (nüè, 8 traits), \zh{疾} (jí, 10 traits), \zh{病} (bìng, 10 traits), \zh{蚊} (wén, 10 traits), \zh{帐} (zhàng, 7 traits), \zh{烧} (shāo, 10 traits). Souligne ensuite dans chaque caractère la clé et écris à côté le nom de sa famille (maladie, insecte, tissu, feu).
\end{exercice}

\begin{corrige}[Exercice de tracé]
\zh{疟} : clé \zh{疒} → famille de la maladie. \zh{疾} : clé \zh{疒} → maladie. \zh{病} : clé \zh{疒} → maladie. \zh{蚊} : clé \zh{虫} → insecte. \zh{帐} : clé \zh{巾} → tissu. \zh{烧} : clé \zh{火} → feu. Vérifie que tu as bien tracé \zh{疒} en premier dans les trois premiers caractères, et \zh{虫} à gauche avant \zh{文} dans \zh{蚊}.
\end{corrige}

\begin{aretenir}[L'essentiel de la section]
\begin{itemize}
\item Les clés (\zh{部首}) donnent un indice de sens et permettent de chercher dans un dictionnaire.
\item \zh{疒} (maladie) : \zh{疟}, \zh{疾}, \zh{病}, \zh{疼} — elle se trace toujours en premier et enveloppe le caractère.
\item \zh{虫} (insecte) : \zh{蚊} ; \zh{巾} (tissu) : \zh{帐} ; \zh{火} (feu) : \zh{烧}.
\item Attention au piège : \zh{帐} (moustiquaire) et \zh{账} (compte) ont le même pinyin \emph{zhàng} mais des clés différentes.
\end{itemize}
\end{aretenir}

\coursec{Mise en route et texte support}

\begin{experience}[Situation de communication : au dispensaire]
Imaginez la scène : un élève de Terminale à Yaoundé accompagne son petit frère au dispensaire. L'infirmier dit : \zh{「可能是疟疾，要马上验血。」} (Kěnéng shì nüèjí, yào mǎshàng yànxiě. — « C'est peut-être le paludisme, il faut faire une prise de sang tout de suite. »). Comprendre le vocabulaire médical de base, les consignes de prévention et savoir décrire des symptômes en chinois est une compétence concrète, utile au Cameroun comme en Chine.
\end{experience}

\begin{textebox}[Texte support : 疟疾 — le paludisme]
疟疾是一种很常见的病。\\
Nüèjí shì yì zhǒng hěn chángjiàn de bìng.\\
« Le paludisme est une maladie très courante. »\\[4pt]
在非洲，很多人每年都会得疟疾。\\
Zài Fēizhōu, hěn duō rén měi nián dōu huì dé nüèjí.\\
« En Afrique, beaucoup de gens attrapent le paludisme chaque année. »\\[4pt]
因为蚊子传播疟疾，所以预防蚊子很重要。\\
Yīnwèi wénzi chuánbō nüèjí, suǒyǐ yùfáng wénzi hěn zhòngyào.\\
« Parce que les moustiques transmettent le paludisme, la prévention contre les moustiques est très importante. »\\[4pt]
得了疟疾的人会发烧、头疼，还会觉得很冷。\\
Dé le nüèjí de rén huì fāshāo, tóuténg, hái huì juéde hěn lěng.\\
« La personne qui a le paludisme a de la fièvre, mal à la tête, et se sent aussi très froid. »\\[4pt]
对孩子和老人来说，疟疾最危险。\\
Duì háizi hé lǎorén láishuō, nüèjí zuì wēixiǎn.\\
« Pour les enfants et les personnes âgées, le paludisme est le plus dangereux. »\\[4pt]
为了预防疟疾，我们应该用蚊帐，还应该保持环境干净。\\
Wèile yùfáng nüèjí, wǒmen yīnggāi yòng wénzhàng, hái yīnggāi bǎochí huánjìng gānjìng.\\
« Pour prévenir le paludisme, nous devons utiliser des moustiquaires et aussi garder l'environnement propre. »\\[4pt]
如果得了疟疾，要马上去医院看病。\\
Rúguǒ dé le nüèjí, yào mǎshàng qù yīyuàn kànbìng.\\
« Si l'on attrape le paludisme, il faut aller tout de suite à l'hôpital se faire soigner. »
\end{textebox}

\coursec{Glossaire du thème : le paludisme}

\begin{center}
\begin{tabularx}{\linewidth}{|X|X|X|X|}
\hline
\rowcolor{popblueL} \textbf{Caractères} & \textbf{Pinyin} & \textbf{Nature} & \textbf{Traduction} \\
\hline
疟疾 & nüèjí & nom & paludisme \\
\hline
病 & bìng & nom & maladie, malade \\
\hline
常见 & chángjiàn & adjectif & courant, fréquent \\
\hline
非洲 & Fēizhōu & nom propre & Afrique \\
\hline
每年 & měinián & adverbe & chaque année, tous les ans \\
\hline
蚊子 & wénzi & nom & moustique \\
\hline
传播 & chuánbō & verbe & transmettre, propager (maladie, info) \\
\hline
预防 & yùfáng & verbe & prévenir,防范 \\
\hline
蚊帐 & wénzhàng & nom & moustiquaire (lit : moustique + filet) \\
\hline
环境 & huánjìng & nom & environnement, cadre \\
\hline
干净 & gānjìng & adjectif & propre, net \\
\hline
症状 & zhèngzhuàng & nom & symptôme \\
\hline
发烧 & fāshāo & verbe & avoir de la fièvre \\
\hline
头疼 & tóuténg & verbe & avoir mal à la tête \\
\hline
冷 & lěng & adjectif & froid (ressenti) \\
\hline
孩子 & háizi & nom & enfant \\
\hline
老人 & lǎorén & nom & personne âgée, vieux \\
\hline
危险 & wēixiǎn & adjectif & dangereux \\
\hline
应该 & yīnggāi & verbe modal & devoir, il faut (conseil moral) \\
\hline
马上 & mǎshàng & adverbe & tout de suite, immédiatement \\
\hline
医院 & yīyuàn & nom & hôpital \\
\hline
看病 & kànbìng & verbe & consulter, se faire soigner \\
\hline
因为 & yīnwèi & conjonction & parce que (cause) \\
\hline
所以 & suǒyǐ & conjonction & donc, par conséquent (conséquence) \\
\hline
为了 & wèile & préposition & pour, afin de (but) \\
\hline
还 & hái & adverbe & aussi, en plus (addition) \\
\hline
\end{tabularx}
\end{center}

\begin{aretenir}[Points de vocabulaire clés]
\begin{itemize}
\item \cle{疟疾 (nüèjí)} : attention au \emph{ü} (u tréma) après n, l, j, q, x. Le ton est \textbf{4e ton} (nüè) + \textbf{2e ton} (jí).
\item \cle{传播 (chuánbō)} vs \cle{传给 (chuán gěi)} : \zh{传播} s'utilise avec une maladie ou une info comme objet direct (transmettre \emph{quelque chose}) ; \zh{传给} introduit le destinataire (transmettre \emph{à quelqu'un}). Le texte utilise \zh{传播疟疾}.
\item \cle{蚊帐 (wénzhàng)} : composé de \zh{蚊} (moustique) + \zh{帐} (tente, filet, compte). Classificateur : \zh{顶} (dǐng) pour les tentes/moustiquaires, ou \zh{个}.
\item \cle{发烧 (fāshāo)} : litt. « émettre de la chaleur ». \zh{发} (fā) = émettre, sortir ; \zh{烧} (shāo) = brûler, fièvre.
\item \cle{头疼 (tóuténg)} : structure \textbf{Corps + 疼} (mal à...). Autre : \zh{肚子疼} (mal au ventre), \zh{牙疼} (mal aux dents).
\item \cle{马上 (mǎshàng)} : litt. « sur le cheval » → immédiatement. Synonyme formel : \zh{立刻} (lìkè).
\end{itemize}
\end{aretenir}

\coursec{Clés (radicaux) et ordre des traits}

\begin{definition}[Rappel méthodologique : identifier la clé et compter les traits]
Pour chercher un caractère dans un dictionnaire (papier ou numérique par radical) : 1) Identifier la \cle{clé} (radical, 部首 bùshǒu) — souvent à gauche, en haut, en bas ou en bas à gauche. 2) Compter les \cle{traits restants} (hors clé). 3) Chercher dans l'index des radicaux par nombre de traits de la clé, puis dans le tableau des caractères par traits restants.
\end{definition}

\begin{center}
\begin{tabularx}{\linewidth}{|X|X|X|X|}
\hline
\rowcolor{popblueL} \textbf{Caractère} & \textbf{Clé (Radical)} & \textbf{Nb traits (total)} & \textbf{Ordre des traits principaux (description)} \\
\hline
\zh{疟} (nuè) & \zh{疒} (nè, « maladie », 5 traits) & 9 & Clé \zh{疒} (haut-gauche : point, horizontale, crochet, point, horizontale) + phonétique \zh{略} (lüè : haut \zh{田}, bas \zh{各}) \\
\hline
\zh{疾} (jí) & \zh{疒} (nè, « maladie ») & 9 & Clé \zh{疒} + phonétique \zh{矢} (shǐ, flèche : horizontale, verticale, horizontale, point, horizontale) \\
\hline
\zh{蚊} (wén) & \zh{虫} (chóng, « insecte », 6 traits) & 10 & Clé \zh{虫} (milieu-bas : verticale, horizontale pliée, horizontale, verticale, horizontale, point) + phonétique \zh{文} (wén, haut : point, horizontale, verticale, horizontale ; bas : point, horizontale) \\
\hline
\zh{帐} (zhàng) & \zh{巾} (jīn, « tissu », 3 traits) & 11 & Clé \zh{巾} (haut-gauche : horizontale, verticale crochet, horizontale) + phonétique \zh{长} (cháng/zhǎng, 8 traits : verticale, horizontale pliée, horizontale, verticale, horizontale, horizontale, horizontale, point) \\
\hline
\zh{预} (yù) & \zh{页} (yè, « tête/page », 6 traits) & 13 & Haut : \zh{予} (yǔ, 4 traits) + Bas : \zh{页} (6 traits : verticale, horizontale pliée, 3 horizontales, crochet final) \\
\hline
\zh{防} (fáng) & \zh{阝} (fù, « mound/oreille », 3 traits, à gauche) & 6 & Clé \zh{阝} (point, verticale, horizontale crochet) + phonétique \zh{方} (fāng, 4 traits : point, horizontale, verticale, horizontale) \\
\hline
\zh{症} (zhèng) & \zh{疒} (nè, « maladie ») & 12 & Clé \zh{疒} + phonétique \zh{正} (zhèng, 5 traits : verticale, horizontale, verticale, horizontale, horizontale) \\
\hline
\zh{状} (zhuàng) & \zh{犬} (quǎn, « chien », 4 traits, en bas) & 7 & Haut : \zh{爿} (qiáng, lit vertical, 4 traits) + Bas : \zh{犬} (point, courbe, point, point) \\
\hline
\end{tabularx}
\end{center}

\begin{experience}[Atelier écriture : « Course aux radicaux »]
En binôme. L'enseignant donne un radical (ex. \zh{疒}, \zh{虫}, \zh{巾}, \zh{阝}). Chaque binôme a 2 minutes pour lister au maximum de caractères du vocabulaire de la leçon (ou connus) contenant ce radical, les écrire au tableau avec l'ordre des traits, et donner le sens du radical. Ex. \zh{疒} → 疟, 疾, 症, 病, 痛, 疲… (maladie). \zh{虫} → 蚊, 蚂, 蚁, 蛋… (insecte/reptile).
\end{experience}

\coursec{Compréhension détaillée du texte}

Après avoir appris à écrire les caractères clés du thème (\zh{疟}, \zh{疾}, \zh{蚊}, \zh{帐}, \zh{预}, \zh{防}, \zh{症}, \zh{状}…) grâce aux radicaux et à l'ordre des traits, nous revenons maintenant sur le texte support pour le comprendre en profondeur. Lire un texte chinois, ce n'est pas seulement reconnaître des caractères : c'est repérer les informations, comprendre les liens logiques entre les phrases et savoir répondre à des questions en phrases complètes. C'est exactement la compétence exigée à l'examen du Baccalauréat.

\courssub{Relecture guidée du texte}

Relisons ensemble le texte support de la leçon. Cette fois, ne cherchez plus les mots nouveaux : écoutez les questions de repérage et montrez du doigt la phrase qui contient la réponse.

\textbf{Questions orales de repérage (à l'oral, réponse courte) :}
\begin{itemize}
\item \zh{疟疾常见吗？} \emph{Nüèjí chángjiàn ma ?} — Le paludisme est-il courant ? (Repérer : \zh{很常见}。)
\item \zh{谁传播疟疾？} \emph{Shéi chuánbō nüèjí ?} — Qui transmet le paludisme ? (Repérer : \zh{蚊子}。)
\item \zh{病人有什么症状？} \emph{Bìngrén yǒu shénme zhèngzhuàng ?} — Quels sont les symptômes du malade ? (Repérer : \zh{发烧、头疼，还会觉得很冷}。)
\item \zh{我们应该怎么预防？} \emph{Wǒmen yīnggāi zěnme yùfáng ?} — Comment devons-nous prévenir ? (Repérer : \zh{用蚊帐，保持环境干净}。)
\item \zh{得了疟疾怎么办？} \emph{Dé le nüèjí zěnme bàn ?} — Que faire si on a le paludisme ? (Repérer : \zh{马上去医院看病}。)
\end{itemize}

\courssub{Analyse des connecteurs logiques du texte}

Un texte bien construit relie ses phrases par des \cle{connecteurs}. En chinois, quatre mots-outils structurent notre texte. Observons-les dans leur phrase avant d'énoncer la règle.

\begin{center}
\begin{tabularx}{\linewidth}{|l|X|X|}
\hline
\rowcolor{popblueL} Connecteur & Fonction et structure & Exemple du texte + traduction \\
\hline
\zh{因为} yīnwèi & Cause. Répond à \zh{为什么} « pourquoi ». Souvent en paire : \zh{因为}…\zh{所以}… & \zh{因为蚊子传播疟疾。} « Parce que les moustiques transmettent le paludisme. » \\
\hline
\zh{为了} wèile & But. Se place en tête de phrase : \zh{为了} + but，sujet + verbe. & \zh{为了预防疟疾，我们应该用蚊帐。} « Pour prévenir le paludisme, nous devons utiliser des moustiquaires. » \\
\hline
\zh{还} hái & Addition : « aussi, en plus ». Se place avant le verbe / l'adjectif. & \zh{还会觉得很冷。} « et se sent aussi très froid. » \\
\hline
\zh{应该} yīnggāi & Conseil, devoir moral : « devoir, il faut ». Verbe modal, avant le verbe principal. & \zh{我们应该保持环境干净。} « Nous devons garder l'environnement propre. » \\
\hline
\end{tabularx}
\end{center}

\begin{propriete}[为了 + but : la phrase de but]
La structure du but est : \textbf{为了 + objectif，+ sujet + 应该/要/必须 + verbe}.\\
Contrairement au français (« pour + infinitif » peut aller partout), en chinois \zh{为了} se place presque toujours \textbf{en tête de phrase}, suivi d'une virgule.\\
\zh{为了预防疟疾，我们应该用蚊帐。} \emph{Wèile yùfáng nüèjí, wǒmen yīnggāi yòng wénzhàng.} « Pour prévenir le paludisme, nous devons utiliser des moustiquaires. »\\
Autre exemple : \zh{为了学好汉语，我每天练习写字。} \emph{Wèile xuéhǎo Hànyǔ, wǒ měi tiān liànxí xiězì.} « Pour bien apprendre le chinois, je m'entraîne chaque jour à écrire. »
\end{propriete}

\begin{attention}[因为 ou 所以 ? Ne confondez pas cause et conséquence]
Quand la question commence par \zh{为什么} « pourquoi », la réponse commence par \zh{因为}. Ne répondez JAMAIS par \zh{所以} à une question de cause : \zh{所以} introduit la \emph{conséquence}, pas la cause.\\
Question : \zh{为什么很多人得疟疾？} « Pourquoi beaucoup de gens attrapent-ils le paludisme ? »\\
Correct : \zh{因为蚊子传播疟疾。}\\
Incorrect : *\zh{所以蚊子传播疟疾。}\\
Astuce : \zh{因为} = « parce que » ; \zh{所以} = « donc ». La question « pourquoi » appelle « parce que », jamais « donc ».
\end{attention}

\begin{attention}[还 (hái) vs 也 (yě) : l'addition]
\zh{也} (aussi) lie deux sujets qui font la même action : \zh{我喜欢足球，他也喜欢足球。} (Moi j'aime le foot, lui \emph{aussi}).
\zh{还} (aussi, en plus) lie deux actions/états du \emph{même} sujet : \zh{我喜欢足球，还喜欢篮球。} (J'aime le foot, \emph{en plus} j'aime le basket).
Dans le texte : \zh{会发烧、头疼，还会觉得很冷} → même sujet (le malade), actions successives/ajoutées.
\end{attention}

\courssub{Méthode : répondre à une question de compréhension}

\begin{methode}[Répondre en phrase complète : la méthode en 4 étapes]
\begin{enumerate}
\item \textbf{Lire la question} et souligner le mot interrogatif : \zh{什么} (quoi), \zh{谁} (qui), \zh{怎么} (comment), \zh{为什么} (pourquoi), \zh{哪里} (où).
\item \textbf{Repérer le mot-clé} de la question dans le texte (ex. \zh{预防} pour une question sur la prévention ; \zh{症状} pour les symptômes).
\item \textbf{Recopier la phrase du texte} qui contient la réponse.
\item \textbf{Adapter} : reprendre les mots de la question dans la réponse pour former une phrase complète. Si la question est \zh{疟疾是一种什么病？}, la réponse commence obligatoirement par \zh{疟疾是一种…的病。}
\end{enumerate}
\end{methode}

\begin{popfigure}
\begin{tikzpicture}[
  node distance=0.45cm,
  box/.style={draw=popblue, fill=popblueL, rounded corners=3pt, align=center, text width=4.6cm, inner sep=5pt, font=\small},
  arr/.style={-{Stealth[length=2.5mm]}, thick, popdark}
]
\node[box, align=center] (q){\textbf{1. Question}\\ 我们应该怎么预防疟疾？};
\node[box, below=of q, align=center] (m){\textbf{2. Repérer le mot-clé}\\ 预防 → phrase 6 du texte};
\node[box, below=of m, align=center] (p){\textbf{3. Recopier la phrase}\\ 我们应该用蚊帐，还应该保持环境干净。};
\node[box, below=of p, align=center] (r){\textbf{4. Adapter à la question}\\ 我们应该用蚊帐，还应该保持环境干净。};
\draw[arr] (q) -- (m);
\draw[arr] (m) -- (p);
\draw[arr] (p) -- (r);
\end{tikzpicture}
\legende{Organigramme de compréhension : de la question à la réponse complète.}
\end{popfigure}

\courssub{Questions de compréhension écrites et correction collective}

Répondez par écrit, en phrases complètes, aux cinq questions suivantes. Puis nous corrigerons ensemble au tableau.

\begin{exercice}[Questions de compréhension]
\begin{enumerate}
\item \zh{疟疾是一种什么病？} \emph{Nüèjí shì yì zhǒng shénme bìng ?}
\item \zh{什么传播疟疾？} \emph{Shénme chuánbō nüèjí ?}
\item \zh{得了疟疾的人有什么症状？} \emph{Dé le nüèjí de rén yǒu shénme zhèngzhuàng ?}
\item \zh{在非洲，疟疾对谁最危险？} \emph{Zài Fēizhōu, nüèjí duì shéi zuì wēixiǎn ?}
\item \zh{我们应该怎么预防疟疾？} \emph{Wǒmen yīnggāi zěnme yùfáng nüèjí ?}
\end{enumerate}
\end{exercice}

\begin{exoresolu}[Correction collective : rédiger les réponses complètes]
Rédiger une réponse complète à la question 5. \tcblower
\textbf{Étape 1} : le mot interrogatif est \zh{怎么} « comment » ; le mot-clé est \zh{预防} « prévenir ».\\
\textbf{Étape 2} : dans le texte, la phrase contenant \zh{预防} est : \zh{为了预防疟疾，我们应该用蚊帐，还应该保持环境干净。}\\
\textbf{Étape 3} : on reprend les mots de la question (\zh{我们应该} + verbe) pour rédiger :\\
\zh{我们应该用蚊帐，还应该保持环境干净。}\\
\emph{Wǒmen yīnggāi yòng wénzhàng, hái yīnggāi bǎochí huánjìng gānjìng.}\\
« Nous devons utiliser des moustiquaires et aussi garder l'environnement propre. »
\end{exoresolu}

\begin{corrige}[Questions 1 à 4]
\begin{enumerate}
\item \zh{疟疾是一种很常见的病。} \emph{Nüèjí shì yì zhǒng hěn chángjiàn de bìng.} « Le paludisme est une maladie très courante. » (On reprend \zh{疟疾是一种…的病}.)
\item \zh{蚊子传播疟疾。} \emph{Wénzi chuánbō nüèjí.} « Les moustiques transmettent le paludisme. »
\item \zh{得了疟疾的人会发烧、头疼，还会觉得很冷。} \emph{Dé le nüèjí de rén huì fāshāo, tóuténg, hái huì juéde hěn lěng.} « Le malade a de la fièvre, mal à la tête et se sent aussi très froid. »
\item \zh{在非洲，疟疾对孩子和老人最危险。} \emph{Zài Fēizhōu, nüèjí duì háizi hé lǎorén zuì wēixiǎn.} « En Afrique, le paludisme est le plus dangereux pour les enfants et les personnes âgées. »
\end{enumerate}
\end{corrige}

\courssub{Vrai ou faux : vérifier sa compréhension globale}

\begin{exercice}[Vrai ou faux ? Justifiez avec une phrase du texte]
\begin{enumerate}
\item \zh{疟疾很少见。} \emph{Nüèjí hěn shǎojiàn.} « Le paludisme est rare. »
\item \zh{蚊帐可以预防疟疾。} \emph{Wénzhàng kěyǐ yùfáng nüèjí.} « La moustiquaire peut prévenir le paludisme. »
\item \zh{得了疟疾不用去医院。} \emph{Dé le nüèjí bú yòng qù yīyuàn.} « Quand on a le paludisme, pas besoin d'aller à l'hôpital. »
\end{enumerate}
\end{exercice}

\begin{corrige}[Vrai ou faux]
\begin{enumerate}
\item \zh{错。} Faux. Le texte dit : \zh{疟疾是一种很常见的病。} (très courante, pas rare).
\item \zh{对。} Vrai. Le texte dit : \zh{我们应该用蚊帐。} (L'usage de la moustiquaire est une mesure de prévention).
\item \zh{错。} Faux. Le texte dit : \zh{如果得了疟疾，要马上去医院看病。} (il faut aller tout de suite à l'hôpital).
\end{enumerate}
\end{corrige}

\courssub{Reformulation : résumer le texte avec ses propres mots}

Dernière étape : résumer le texte en \textbf{trois phrases}, sans recopier mot à mot. C'est l'exercice le plus difficile mais le plus utile : il prouve que vous avez vraiment compris.

\begin{exemplebox}[Modèle de résumé en trois phrases]
\zh{疟疾是非洲很常见的病，蚊子传播给人。}\\
\emph{Nüèjí shì Fēizhōu hěn chángjiàn de bìng, wénzi chuánbō gěi rén.}\\
« Le paludisme est une maladie très courante en Afrique, les moustiques le transmettent aux hommes. »\\[4pt]
\zh{病人会发烧、头疼，孩子和老人最危险。}\\
\emph{Bìngrén huì fāshāo, tóuténg, háizi hé lǎorén zuì wēixiǎn.}\\
« Le malade a de la fièvre et mal à la tête ; les enfants et les personnes âgées sont les plus menacés. »\\[4pt]
\zh{为了预防疟疾，我们要用蚊帐，保持环境干净。}\\
\emph{Wèile yùfáng nüèjí, wǒmen yào yòng wénzhàng, bǎochí huánjìng gānjìng.}\\
« Pour prévenir le paludisme, nous devons utiliser des moustiquaires et garder l'environnement propre. »
\end{exemplebox}

\coursec{Collocations et production guidée}

\courssub{Collocations fréquentes (搭配 dāipèi) du thème}

Maîtriser un mot, c'est savoir avec quels autres mots il s'associe naturellement. Complétez le tableau ci-dessous en associant le verbe de la colonne de gauche aux noms de la colonne de droite (plusieurs réponses possibles). Écrivez les collocations en caractères + pinyin.

\begin{center}
\begin{tabularx}{\linewidth}{|X|X|X|}
\hline
\rowcolor{popblueL} \textbf{Verbe} & \textbf{Nom (Objet / Complément)} & \textbf{Collocations correctes (Caractères + Pinyin)} \\
\hline
\zh{得} (dé, attraper) & 疟疾、病、感冒 & \zh{得疟疾 (dé nüèjí)}, \zh{得病 (dé bìng)}, \zh{得感冒 (dé gǎnmào)} \\
\hline
\zh{传播} (chuánbō, transmettre) & 疟疾、病毒、疾病、信息 & \zh{传播疟疾 (chuánbō nüèjí)}, \zh{传播病毒 (chuánbō bìngdú)}, \zh{传播疾病 (chuánbō jíbìng)} \\
\hline
\zh{预防} (yùfáng, prévenir) & 疟疾、疾病、感冒、传染病 & \zh{预防疟疾 (yùfáng nüèjí)}, \zh{预防疾病 (yùfáng jíbìng)}, \zh{预防传染病 (yùfáng chuánrǎnbìng)} \\
\hline
\zh{用} (yòng, utiliser) & 蚊帐、药、驱蚊液、网 & \zh{用蚊帐 (yòng wénzhàng)}, \zh{用药 (yòng yào)}, \zh{用驱蚊液 (yòng qūwényè)} \\
\hline
\zh{保持} (bǎochí, maintenir) & 环境干净、健康、安静、清醒 & \zh{保持环境干净 (bǎochí huánjìng gānjìng)}, \zh{保持健康 (bǎochí jiànkāng)} \\
\hline
\zh{去} (qù, aller) & 医院、诊所、药店 & \zh{去医院 (qù yīyuàn)}, \zh{去诊所 (qù zhěnsuǒ)}, \zh{去药店 (qù yàodiàn)} \\
\hline
\zh{看病} (kànbìng, se soigner) & (verbe intransitif ou + médecin) & \zh{看病 (kànbìng)}, \zh{找医生看病 (zhǎo yīshēng kànbìng)} \\
\hline
\end{tabularx}
\end{center}

\begin{exercice}[Production écrite guidée : Affiche de prévention]
Contexte : Le lycée organise une semaine « Santé et Environnement ». Vous devez rédiger une \textbf{affiche en chinois simplifié} (titre + 3 à 4 consignes) pour prévenir le paludisme au Cameroun. Utilisez le vocabulaire et les structures de la leçon (\zh{为了…应该…}, \zh{因为…所以…}, \zh{马上}, \zh{预防}).

\textbf{Contraintes :}
\begin{itemize}
\item Titre clair (ex. \zh{预防疟疾，从我做起}).
\item Au moins une phrase avec \zh{为了} (but).
\item Au moins une phrase avec \zh{应该} (conseil).
\item Vocabulaire : \zh{蚊帐}, \zh{环境干净}, \zh{蚊子}, \zh{医院}.
\item Écriture soignée (caractères corrects, ordre des traits respecté).
\end{itemize}
\end{exercice}

\begin{exemplebox}[Modèle d'affiche (corrigé type)]
\textbf{预防疟疾，从我做起}\\
\emph{Yùfáng nüèjí, cóng wǒ zuòqǐ}\\
« Prévenir le paludisme, ça commence par moi »\\[6pt]
\zh{因为蚊子传播疟疾，所以我们要防蚊。}\\
\emph{Yīnwèi wénzi chuánbō nüèjí, suǒyǐ wǒmen yào fángwén.}\\
« Parce que les moustiques transmettent le paludisme, nous devons nous protéger des moustiques. »\\[4pt]
\zh{为了预防疟疾，大家应该睡蚊帐。}\\
\emph{Wèile yùfáng nüèjí, dàjiā yīnggāi shuì wénzhàng.}\\
« Pour prévenir le paludisme, tout le monde doit dormir sous moustiquaire. »\\[4pt]
\zh{还应该保持环境干净，不留积水。}\\
\emph{Hái yīnggāi bǎochí huánjìng gānjìng, bù liú jīshuǐ.}\\
« Il faut aussi garder l'environnement propre, ne pas laisser d'eau stagnante. »\\[4pt]
\zh{如果发烧、头疼，马上去医院看病！}\\
\emph{Rúguǒ fāshāo, tóuténg, mǎshàng qù yīyuàn kànbìng!}\\
« Si fièvre et maux de tête, allez tout de suite à l'hôpital ! »
\end{exemplebox}

\courssub{Production orale : Exposé court (1 minute)}

\begin{experience}[Débat / Exposé : « Paludisme au Cameroun vs Chine »]
Préparez un exposé oral de 1 minute (sans lire, notes autorisées) sur l'un des sujets :
\begin{enumerate}
\item \textbf{Situation au Cameroun} : Zones à risque (Sud, Est, Littoral), saison des pluies, moustiquaires imprégnées (MILDA), traitement (ACT), rôle de l'école.
\item \textbf{Éradication en Chine} : La Chine certifiée « sans paludisme » par l'OMS en 2021. Stratégies : « 1-3-7 » (déclaration 1j, investigation 3j, riposte 7j), recherche sur l'artémisinine (屠呦呦 Tú Yōuyōu, prix Nobel 2015).
\item \textbf{Comparaison} : Points communs (moustique anophèle, prévention vectorielle) et différences (climat, système de santé, éradication vs contrôle).
\end{enumerate}
\textbf{Critères d'évaluation :} Prononciation (tons), structures (\zh{因为/所以}, \zh{为了/应该}, \zh{还}), vocabulaire précis, cohérence, contact visuel.
\end{experience}

\begin{savaistu}[Le saviez-vous ? L'artémisinine : un cadeau de la Chine au monde]
La découverte de l'\textbf{artémisinine} (青蒿素 qīnghāosù) par la pharmacologue chinoise \textbf{屠呦呦 (Tú Yōuyōu)} dans les années 1970, à partir de la médecine traditionnelle (\zh{青蒿} qīnghāo, armoise annuelle), a révolutionné le traitement du paludisme. Elle a reçu le \textbf{Prix Nobel de Physiologie ou Médecine 2015}. C'est le premier Nobel scientifique décerné à une chercheuse chinoise ayant travaillé en Chine. L'OMS recommande les thérapies combinées à base d'artémisinine (ACT) comme traitement de première ligne. La Chine a partagé ce savoir-faire avec l'Afrique (dont le Cameroun) via la coopération médicale et la construction d'usines de production locale.
\end{savaistu}

\begin{aretenir}[L'essentiel de la leçon ZH12 — Paludisme]
\begin{itemize}
\item \textbf{Vocabulaire} : 25 mots-clés (maladie, transmission, symptômes, prévention, soin) — Tableau glossaire à maîtriser (caractères, pinyin, sens).
\item \textbf{Écriture} : 8 caractères analysés par radical (\zh{疒, 虫, 巾, 页, 阝, 犬}) et ordre des traits.
\item \textbf{Grammaire / Connecteurs} :
      \begin{itemize}
      \item Cause : \zh{因为}…\zh{所以}… (question \zh{为什么} → réponse \zh{因为}).
      \item But : \zh{为了} + but, sujet + \zh{应该/要} + verbe (tête de phrase).
      \item Addition : \zh{还} (même sujet, action en plus) vs \zh{也} (sujets différents, même action).
      \item Conseil : \zh{应该} (devoir moral).
      \end{itemize}
\item \textbf{Compréhension} : Méthode en 4 étapes (Question → Mot-clé → Recopie → Adaptation). Réponse en phrase complète.
\item \textbf{Production} : Rédiger une affiche de prévention (structure \zh{为了…应该…}) ; Exposé oral comparatif Cameroun/Chine.
\item \textbf{Culture} : Artémisinine (屠呦呦), éradication en Chine (2021), coopération sino-africaine.
\end{itemize}
\end{aretenir}

\coursec{Collocations et production guidée}

Dans la section précédente, nous avons lu et analysé en détail le texte sur le paludisme : nous savons maintenant \emph{ce que dit} le texte. Il reste une étape décisive : savoir \emph{réutiliser} ce vocabulaire pour parler et écrire soi-même. En chinois comme en français, on ne retient pas les mots isolément, mais dans des \cle{collocations}, c'est-à-dire des associations stables de mots : on dit « attraper le paludisme », « avoir une forte fièvre », « utiliser une moustiquaire ». Le chinois a ses propres associations, qui ne coïncident pas toujours avec le français. Cette section vous donne les sept collocations essentielles du thème, puis vous guide pas à pas pour produire cinq phrases correctes sur le paludisme au Cameroun.

\courssub{Observer d'abord : les collocations dans des phrases}

\begin{exemplebox}[Les collocations en contexte]
\zh{蚊子传染疾病。} \emph{Wénzi chuánrǎn jíbìng.} « Les moustiques transmettent des maladies. »\\[2pt]
\zh{他发高烧，要马上去医院。} \emph{Tā fā gāoshāo, yào mǎshàng qù yīyuàn.} « Il a une forte fièvre, il doit aller immédiatement à l'hôpital. »\\[2pt]
\zh{去年我弟弟得了疟疾。} \emph{Qùnián wǒ dìdi dé le nüèjí.} « L'année dernière, mon petit frère a attrapé le paludisme. »\\[2pt]
\zh{我们应该用蚊帐睡觉。} \emph{Wǒmen yīnggāi yòng wénzhàng shuìjiào.} « Nous devrions dormir sous une moustiquaire. »\\[2pt]
\zh{保持环境干净可以预防疾病。} \emph{Bǎochí huánjìng gānjìng kěyǐ yùfáng jíbìng.} « Garder l'environnement propre permet de prévenir les maladies. »\\[2pt]
\zh{生病了要及时看病。} \emph{Shēngbìng le yào jíshí kànbìng.} « Quand on est malade, il faut consulter à temps. »
\end{exemplebox}

Observez : chaque collocation est un bloc \cle{verbe + complément}. En chinois, ce bloc se mémorise comme une unité insécable : on ne dit pas \zh{拿疟疾} ni \zh{做高烧}, exactement comme en français on ne dit pas « faire le paludisme ».

\courssub{Les sept collocations essentielles}

\begin{definition}[Collocation]
Une \cle{collocation} est une association habituelle et stable de deux mots (ici : verbe + nom). En chinois, le choix du verbe est fixe : \zh{得} avec \zh{疟疾}, \zh{发} avec \zh{高烧}, \zh{用} avec \zh{蚊帐}. Remplacer le verbe par un synonyme produit une phrase incorrecte ou artificielle.
\end{definition}

\begin{center}
\begin{tabularx}{\linewidth}{|X|X|X|}
\hline
\rowcolor{popblueL} Collocation & Pinyin & Traduction \\
\hline
\zh{得疟疾} & dé nüèjí & attraper le paludisme \\
\hline
\zh{传染疾病} & chuánrǎn jíbìng & transmettre une maladie \\
\hline
\zh{发高烧} & fā gāoshāo & avoir une forte fièvre \\
\hline
\zh{预防疾病} & yùfáng jíbìng & prévenir les maladies \\
\hline
\zh{用蚊帐} & yòng wénzhàng & utiliser une moustiquaire \\
\hline
\zh{保持环境干净} & bǎochí huánjìng gānjìng & garder l'environnement propre \\
\hline
\zh{及时看病} & jíshí kànbìng & consulter (un médecin) à temps \\
\hline
\end{tabularx}
\end{center}

\begin{popfigure}
\begin{tikzpicture}[
  verbe/.style={rectangle, rounded corners=3pt, draw=popblue, fill=popblueL, align=center, font=\small, minimum height=0.8cm},
  compl/.style={rectangle, rounded corners=3pt, draw=poporange, fill=poporangeL, align=center, font=\small, minimum height=0.8cm},
  trad/.style={font=\small\itshape, text=popdark, align=left},
  lien/.style={-{Stealth[length=2mm]}, thick, popmuted}
]
\node[verbe, align=center] (v1) at (0,0){得 dé\\attraper};
\node[compl, right=1.4cm of v1, align=center] (c1){疟疾 nüèjí\\le paludisme};
\node[trad, right=0.4cm of c1] {attraper le paludisme};
\draw[lien] (v1) -- (c1);

\node[verbe, align=center] (v2) at (0,-1.2){传染 chuánrǎn\\transmettre};
\node[compl, right=0.62cm of v2, align=center] (c2){疾病 jíbìng\\une maladie};
\node[trad, right=0.4cm of c2] {transmettre une maladie};
\draw[lien] (v2) -- (c2);

\node[verbe, align=center] (v3) at (0,-2.4){发 fā\\développer};
\node[compl, right=1.18cm of v3, align=center] (c3){高烧 gāoshāo\\forte fièvre};
\node[trad, right=0.4cm of c3] {avoir une forte fièvre};
\draw[lien] (v3) -- (c3);

\node[verbe, align=center] (v4) at (0,-3.6){预防 yùfáng\\prévenir};
\node[compl, right=0.85cm of v4, align=center] (c4){疾病 jíbìng\\les maladies};
\node[trad, right=0.4cm of c4] {prévenir les maladies};
\draw[lien] (v4) -- (c4);

\node[verbe, align=center] (v5) at (0,-4.8){用 yòng\\utiliser};
\node[compl, right=1.32cm of v5, align=center] (c5){蚊帐 wénzhàng\\moustiquaire};
\node[trad, right=0.4cm of c5] {utiliser une moustiquaire};
\draw[lien] (v5) -- (c5);

\node[verbe, align=center] (v6) at (0,-6.0){保持 bǎochí\\maintenir};
\node[compl, right=0.85cm of v6, align=center] (c6){环境干净\\huánjìng gānjìng};
\node[trad, right=0.4cm of c6] {garder l'environnement propre};
\draw[lien] (v6) -- (c6);

\node[verbe, align=center] (v7) at (0,-7.2){及时 jíshí\\à temps};
\node[compl, right=1.08cm of v7, align=center] (c7){看病 kànbìng\\consulter};
\node[trad, right=0.4cm of c7] {consulter à temps};
\draw[lien] (v7) -- (c7);
\end{tikzpicture}
\legende{Les sept collocations du thème : verbe (bleu) + complément (orange), avec la traduction en regard.}
\end{popfigure}

\begin{propriete}[Deux verbes-pivots à connaître : 得 et 发]
\begin{itemize}
\item \zh{得} (dé) + nom de maladie = « attraper, contracter » : \zh{得疟疾} « attraper le paludisme », \zh{得病} « tomber malade ». Accompli : \zh{得了疟疾} « a attrapé le paludisme ».
\item \zh{发} (fā) + nom de symptôme = « développer, avoir » : \zh{发高烧} « avoir une forte fièvre », \zh{发烧} « avoir de la fièvre ». Ici \zh{发} ne signifie pas « envoyer » : c'est un emploi lexicalisé.
\end{itemize}
Comparaison avec le français : le français utilise « avoir » (avoir de la fièvre) et « attraper » (attraper le paludisme) ; le chinois utilise deux verbes spécialisés, \zh{发} et \zh{得}, qu'il ne faut pas intervertir.
\end{propriete}

\begin{attention}[Pièges fréquents des francophones]
\begin{itemize}
\item On dit \zh{用蚊帐} (yòng wénzhàng) avec \zh{用} « utiliser », et non \zh{使蚊帐}. Le verbe \zh{使} (shǐ) signifie « faire, rendre » (\zh{使我很累} « cela me fatigue ») et ne s'emploie pas pour « se servir d'un objet ».
\item \zh{看病} (kànbìng) signifie « consulter un médecin », littéralement « regarder la maladie ». Ne traduisez pas mot à mot : \zh{我要去看病} = « je vais chez le médecin », et non « je vais regarder la maladie ».
\item Attention aux tons : \zh{疟疾} nüèjí (4e ton + 2e ton), \zh{高烧} gāoshāo (1er + 1er ton), \zh{蚊帐} wénzhàng (2e + 4e ton).
\item Ne confondez pas \zh{预防疾病} « prévenir les maladies » (action d'empêcher) et \zh{传染疾病} « transmettre les maladies » (action de propager) : sens presque opposés.
\end{itemize}
\end{attention}

\begin{exoresolu}[Complétion : insérer la bonne collocation]
Complétez avec : \zh{得疟疾 / 传染疾病 / 发高烧 / 预防疾病 / 用蚊帐 / 保持环境干净 / 及时看病}.\\[2pt]
1. 蚊子可以\_\_\_\_\_\_\_\_。\\
2. 他\_\_\_\_\_\_\_\_，体温三十九度。\\
3. 睡觉的时候，我们应该\_\_\_\_\_\_\_\_。\\
4. 打扫院子就是\_\_\_\_\_\_\_\_。\\
5. 感觉不舒服，要\_\_\_\_\_\_\_\_。
\tcblower
\textbf{Solution détaillée.}\\
1. \zh{蚊子可以传染疾病。} \emph{Wénzi kěyǐ chuánrǎn jíbìng.} « Les moustiques peuvent transmettre des maladies. » — sujet : les moustiques ; seule \zh{传染疾病} convient.\\
2. \zh{他发高烧，体温三十九度。} \emph{Tā fā gāoshāo, tǐwēn sānshíjiǔ dù.} « Il a une forte fièvre, sa température est de 39 degrés. » — l'indice « 39 degrés » impose \zh{发高烧}.\\
3. \zh{睡觉的时候，我们应该用蚊帐。} \emph{Shuìjiào de shíhou, wǒmen yīnggāi yòng wénzhàng.} « Quand nous dormons, nous devrions utiliser une moustiquaire. »\\
4. \zh{打扫院子就是保持环境干净。} \emph{Dǎsǎo yuànzi jiùshì bǎochí huánjìng gānjìng.} « Balayer la cour, c'est garder l'environnement propre. »\\
5. \zh{感觉不舒服，要及时看病。} \emph{Gǎnjué bù shūfu, yào jíshí kànbìng.} « Si on ne se sent pas bien, il faut consulter à temps. »
\end{exoresolu}

\begin{exercice}[À vous : complétion et traduction]
1. Complétez : 去年，我妹妹\_\_\_\_\_\_\_\_，在医院住了一个星期。(attraper le paludisme, accompli)\\
2. Complétez : 为了\_\_\_\_\_\_\_\_，我们要把脏水倒掉。(prévenir les maladies)\\
3. Traduisez en chinois : « Les enfants doivent dormir sous une moustiquaire. »
\end{exercice}

\begin{corrige}[Exercice « À vous »]
1. \zh{去年，我妹妹得了疟疾，在医院住了一个星期。} \emph{Qùnián, wǒ mèimei dé le nüèjí, zài yīyuàn zhù le yí ge xīngqī.} — accompli : \zh{得了疟疾}.\\
2. \zh{为了预防疾病，我们要把脏水倒掉。} \emph{Wèile yùfáng jíbìng, wǒmen yào bǎ zāng shuǐ dàodiào.}\\
3. \zh{孩子们应该用蚊帐睡觉。} \emph{Háizimen yīnggāi yòng wénzhàng shuìjiào.}
\end{corrige}

\courssub{Production guidée : cinq phrases sur le paludisme au Cameroun}

\begin{methode}[Rédiger cinq phrases en cinq étapes]
\begin{enumerate}
\item \textbf{Une phrase = une idée.} N'essayez pas de tout dire dans une seule phrase : définition, transmission, symptômes, prévention, conseil.
\item \textbf{Partez de la trame} : \zh{疟疾是…} / \zh{蚊子…} / \zh{病人会…} / \zh{我们应该…} / \zh{为了健康，…}
\item \textbf{Complétez avec les collocations} apprises : \zh{传染疾病}, \zh{发高烧}, \zh{用蚊帐}, \zh{及时看病}…
\item \textbf{Enrichissez} avec les connecteurs : \zh{因为} (parce que), \zh{为了} (afin de), \zh{还} (aussi, de plus).
\item \textbf{Relisez} : tons du pinyin, caractères, verbe correct dans chaque collocation.
\end{enumerate}
\end{methode}

\begin{textebox}[Modèle de production : 喀麦隆的疟疾]
\zh{疟疾是一种很危险的疾病。}\\
\emph{Nüèjí shì yì zhǒng hěn wēixiǎn de jíbìng.}\\
« Le paludisme est une maladie très dangereuse. »\\[3pt]
\zh{在喀麦隆，蚊子传染这种疾病。}\\
\emph{Zài Kāmàilóng, wénzi chuánrǎn zhè zhǒng jíbìng.}\\
« Au Cameroun, les moustiques transmettent cette maladie. »\\[3pt]
\zh{病人会发高烧，还会头疼。}\\
\emph{Bìngrén huì fā gāoshāo, hái huì tóuténg.}\\
« Le malade a une forte fièvre et aussi mal à la tête. »\\[3pt]
\zh{我们应该用蚊帐睡觉，保持环境干净。}\\
\emph{Wǒmen yīnggāi yòng wénzhàng shuìjiào, bǎochí huánjìng gānjìng.}\\
« Nous devrions dormir sous une moustiquaire et garder l'environnement propre. »\\[3pt]
\zh{为了健康，生病了要及时看病。健康很重要。}\\
\emph{Wèile jiànkāng, shēngbìng le yào jíshí kànbìng. Jiànkāng hěn zhòngyào.}\\
« Pour la santé, quand on tombe malade, il faut consulter à temps. La santé est très importante. »
\end{textebox}

Remarquez la progression : phrase 1 = définition (\zh{疟疾是…}) ; phrase 2 = transmission (\zh{蚊子传染…}) ; phrase 3 = symptômes (\zh{病人会…} enrichi avec \zh{还}) ; phrase 4 = prévention (\zh{我们应该…}) ; phrase 5 = conseil et conclusion (\zh{为了健康，…}).

\begin{aretenir}[Complément utile à l'examen : la formule de conclusion]
Pour terminer toute production écrite ou orale sur la santé, la phrase \zh{健康很重要。} \emph{Jiànkāng hěn zhòngyào.} « La santé est très importante. » est une conclusion simple, correcte et toujours pertinente. Variante enrichie : \zh{因为健康很重要，所以我们要预防疾病。} « Parce que la santé est très importante, nous devons prévenir les maladies. »
\end{aretenir}

\courssub{Relecture par les pairs}

\begin{experience}[Relecture croisée en binôme]
Échangez votre production de cinq phrases avec votre voisin. Vérifiez trois points, dans l'ordre :
\begin{enumerate}
\item \textbf{Les tons} : lisez chaque phrase à voix haute ; le pinyin porte-t-il les bonnes marques (\zh{nüèjí}, \zh{gāoshāo}, \zh{wénzhàng}) ?
\item \textbf{Les caractères} : chaque caractère est-il complet et correct (pas de confusion entre \zh{疾} et \zh{病}, entre \zh{帐} et \zh{账}) ?
\item \textbf{Les collocations} : le verbe est-il le bon ? (\zh{得疟疾}, pas \zh{有疟疾} ; \zh{用蚊帐}, pas \zh{使蚊帐} ; \zh{发高烧}, pas \zh{得高烧}.)
\end{enumerate}
Signalez chaque erreur au crayon, proposez une correction, puis rendez le texte à son auteur pour réécriture.
\end{experience}

\begin{savaistu}[Moustiquaires gratuites au Cameroun]
Au Cameroun, des campagnes nationales distribuent gratuitement des moustiquaires imprégnées d'insecticide aux familles, en particulier aux femmes enceintes et aux enfants. En chinois, on dit : \zh{免费发蚊帐} \emph{miǎnfèi fā wénzhàng} « distribuer gratuitement des moustiquaires » (\zh{免费} « gratuit », \zh{发} « distribuer » — encore un autre emploi du verbe \zh{发} !). La Chine, elle, a éliminé le paludisme sur son territoire après des décennies de programmes de prévention : un exemple souvent cité dans la coopération sanitaire sino-africaine.
\end{savaistu}

\begin{exercice}[Production finale]
Rédigez cinq phrases sur le paludisme au Cameroun en suivant la trame : \zh{疟疾是…} / \zh{蚊子…} / \zh{病人会…} / \zh{我们应该…} / \zh{为了健康，…}. Utilisez au moins quatre collocations de la leçon et un connecteur (\zh{因为}, \zh{为了} ou \zh{还}). Terminez par \zh{健康很重要。}
\end{exercice}

\begin{corrige}[Exercice « Production finale » — exemple de réponse attendue]
\zh{疟疾是一种很常见的疾病。} \emph{Nüèjí shì yì zhǒng hěn chángjiàn de jíbìng.} « Le paludisme est une maladie très courante. »\\
\zh{蚊子传染这种疾病。} \emph{Wénzi chuánrǎn zhè zhǒng jíbìng.} « Les moustiques transmettent cette maladie. »\\
\zh{病人会发高烧，还会觉得很累。} \emph{Bìngrén huì fā gāoshāo, hái huì juéde hěn lèi.} « Le malade a une forte fièvre et se sent aussi très fatigué. »\\
\zh{我们应该用蚊帐，还要保持环境干净。} \emph{Wǒmen yīnggāi yòng wénzhàng, hái yào bǎochí huánjìng gānjìng.} « Nous devrions utiliser une moustiquaire et aussi garder l'environnement propre. »\\
\zh{为了健康，生病了要及时看病。健康很重要。} \emph{Wèile jiànkāng, shēngbìng le yào jíshí kànbìng. Jiànkāng hěn zhòngyào.} « Pour la santé, quand on tombe malade, il faut consulter à temps. La santé est très importante. »\\
Barème indicatif : 1 point par phrase correcte (collocation exacte, ordre des mots, tons), 1 point pour le connecteur, 1 point pour la conclusion.
\end{corrige}

Vous savez désormais non seulement comprendre un texte sur le paludisme, mais aussi en parler et en écrire avec des collocations exactes : c'est précisément la compétence attendue à l'examen, en compréhension comme en expression écrite.

\newpage

\coursec{Activité d'intégration}

\coursec{Leçon 12 — Vocabulaire et compréhension de texte : le paludisme}

\courssub{Mise en route et texte support}

\courssub{Situation de communication}
Avril est le mois de la lutte contre le paludisme. À l'occasion de la Journée mondiale contre le paludisme (25 avril), le club de santé de ton lycée à Yaoundé organise une campagne de sensibilisation. Le proviseur a confié à la classe de chinois de Terminale la rédaction de l'affiche en chinois, langue de travail de nombreux partenaires sanitaires au Cameroun. Voici le document d'information distribué par le club, base de votre travail.

\begin{textebox}[Document d'information du club de santé — Avril 2025]
同学们，你们好！\\
Tóngxuémen, nǐmen hǎo!\\
« Chers camarades, bonjour ! »\\[4pt]
疟疾是喀麦隆很常见的疾病。蚊子传播疟疾寄生虫。得了疟疾的人常常发烧、头疼、全身没有力气。\\
Nüèjí shì Kāmàilóng hěn chángjiàn de jíbìng. Wénzi chuánbō nüèjí jìshēngchóng. Dé le nüèjí de rén chángcháng fāshāo, tóu téng, quánshēn méiyǒu lìqi.\\
« Le paludisme est une maladie très courante au Cameroun. Les moustiques transmettent le parasite du paludisme. Les personnes atteintes ont souvent de la fièvre, mal à la tête et manquent de forces dans tout le corps. »\\[4pt]
怎么预防疟疾呢？第一，晚上睡觉的时候要用蚊帐。第二，要保持环境干净，清除积水，防止蚊子滋生。第三，如果发烧了，要及时看病、吃药。\\
Zěnme yùfáng nüèjí ne? Dì-yī, wǎnshang shuìjiào de shíhou yào yòng wénzhàng. Dì-èr, yào bǎochí huánjìng gānjìng, qīngchú jīshuǐ, fángzhǐ wénzi zīshēng. Dì-sān, rúguǒ fāshāo le, yào jíshí kànbìng, chī yào.\\
« Comment prévenir le paludisme ? Premièrement, il faut utiliser une moustiquaire quand on dort le soir. Deuxièmement, il faut garder l'environnement propre, éliminer les eaux stagnantes pour empêcher la prolifération des moustiques. Troisièmement, en cas de fièvre, il faut consulter un médecin sans tarder et prendre des médicaments. »\\[4pt]
预防疟疾，保护健康！\\
Yùfáng nüèjí, bǎohù jiànkāng!\\
« Prévenons le paludisme, protégeons la santé ! »
\end{textebox}

\begin{attention}[Correction scientifique importante]
Dans la version originale, on lisait \zh{疟疾病毒} (nüèjí bìngdú, « virus du paludisme »). C'est une erreur : le paludisme est causé par un \textbf{parasite} (plasmodium), \zh{疟原虫} (nüèyuánchóng) ou \zh{寄生虫} (jìshēngchóng), et non par un virus. Le texte ci-dessus a été corrigé : \zh{蚊子传播疟疾寄生虫} (Wénzi chuánbō nüèjí jìshēngchóng). Retenez \zh{传播} (chuánbō, « transmettre, propager ») pour une maladie.
\end{attention}

\courssub{Glossaire du thème : vocabulaire du paludisme}

\begin{center}
\begin{tabularx}{\linewidth}{|X|X|X|X|}
\hline
\rowcolor{popblueL} \textbf{Caractères} & \textbf{Pinyin} & \textbf{Nature} & \textbf{Traduction} \\
\hline
疟疾 & nüèjí & n. & paludisme (malaria) \\
\hline
疾病 & jíbìng & n. & maladie, pathologie \\
\hline
症状 & zhèngzhuàng & n. & symptôme \\
\hline
发烧 & fāshāo & v. & avoir de la fièvre \\
\hline
头疼 & tóu téng & v. / n. & avoir mal à la tête / mal de tête \\
\hline
全身 & quánshēn & n. & tout le corps \\
\hline
力气 & lìqi & n. & force, énergie \\
\hline
蚊子 & wénzi & n. & moustique \\
\hline
蚊帐 & wénzhàng & n. & moustiquaire \\
\hline
寄生虫 & jìshēngchóng & n. & parasite \\
\hline
传播 & chuánbō & v. & transmettre, propager (maladie, info) \\
\hline
预防 & yùfáng & v. & prévenir \\
\hline
环境 & huánjìng & n. & environnement \\
\hline
干净 & gānjìng & adj. & propre \\
\hline
积水 & jīshuǐ & n. & eau stagnante \\
\hline
滋生 & zīshēng & v. & proliférer, se multiplier (bactéries, moustiques) \\
\hline
及时 & jíshí & adv. & à temps, sans tarder \\
\hline
看病 & kànbìng & v. & consulter un médecin \\
\hline
吃药 & chī yào & v. & prendre des médicaments \\
\hline
健康 & jiànkāng & n. / adj. & santé / en bonne santé \\
\hline
保护 & bǎohù & v. & protéger \\
\hline
\end{tabularx}
\end{center}

\courssub{Clés (radicaux) et ordre des traits}

\courssub{Principes généraux de l'ordre des traits (复习)}
\begin{enumerate}
\item \cle{Haut avant bas} (从上到下) : 三 、 言 。
\item \cle{Gauche avant droite} (从左到右) : 好 、 蚊 。
\item \cle{Horizontal avant vertical} (先横后竖) : 十 、 干 。
\item \cle{Extérieur avant intérieur} (先外后里) : 同 、 病 (pour le cadre) 。
\item \cle{Fermer en dernier} (最后封口) : 回 、 国 。
\item \cle{Trait central avant les ailes} (先中间后两边) : 小 、 水 。
\end{enumerate}

\begin{propriete}[Analyse de caractères clés du thème]
\textbf{1. 病  (bìng, maladie) — 10 traits}
\begin{itemize}
\item \cle{Clé (radical) :} \zh{疒} (nè, « maladie », 5 traits). Elle se place en haut à gauche (structure \cle{haut-gauche / bas-droite}).
\item \cle{Partie phonétique :} \zh{丙} (bǐng, 7 traits) à droite.
\item \cle{Ordre des traits détaillé :}
\begin{enumerate}
\item Points du haut de 疒 (丶 丶)
\item Trait horizontal du haut de 疒 (一)
\item Trait vertical gauche de 疒 (丨)
\item Crochet vertical du bas de 疒 (㇆)
\item Trait horizontal de 丙 (一)
\item Trait vertical de 丙 (丨)
\item Trait horizontal du milieu de 丙 (一)
\item Trait horizontal du bas de 丙 (一)
\item Crochet coudé du bas de 丙 (㇟)
\item Point final de 丙 (丶) — \emph{attention : le dernier trait de 丙 est un point, pas un horizontal}
\end{enumerate}
\end{itemize}

\textbf{2. 蚊  (wén, moustique) — 10 traits}
\begin{itemize}
\item \cle{Clé (radical) :} \zh{虫} (chóng, « insecte », 6 traits). Elle se place à gauche (structure \cle{gauche-droite}).
\item \cle{Partie phonétique :} \zh{文} (wén, 4 traits) à droite.
\item \cle{Ordre des traits détaillé :}
\begin{enumerate}
\item Trait horizontal de 虫 (一)
\item Crochet vertical de 虫 (㇆)
\item Trait horizontal du milieu de 虫 (一)
\item Trait horizontal du bas de 虫 (一)
\item Trait vertical central de 虫 (丨)
\item Point final de 虫 (丶) — \emph{fermeture du cadre}
\item Trait horizontal de 文 (一) — \emph{partie droite commence}
\item Trait vertical de 文 (丨)
\item Trait horizontal du bas de 文 (一)
\item Point final de 文 (丶)
\end{enumerate}
\end{itemize}

\textbf{3. 疟  (nüè, paludisme / fièvre intermittente) — 9 traits}
\begin{itemize}
\item \cle{Clé :} \zh{疒} (nè, maladie).
\item \cle{Partie phonétique :} \zh{略} (lüè).
\item \cle{Note :} Caractère complexe (HSK 5+), à reconnaître prioritairement dans \zh{疟疾}.
\end{itemize}

\textbf{4. 预  (yù, prévenir / à l'avance) — 16 traits}
\begin{itemize}
\item \cle{Clé :} \zh{页} (yè, tête/page), en bas (structure \cle{haut-bas}).
\item \cle{Partie haute :} \zh{予} (yǔ) + \zh{ㄗ} (simplifié de 矢).
\item \cle{Astuce :} \zh{预防} = « agir à l'avance pour empêcher ».
\end{itemize}
\end{propriete}

\begin{popfigure}
\begin{tikzpicture}[
  node distance=0.8cm and 1.2cm,
  box/.style={draw=popblue, thick, rounded corners=3pt, fill=popblueL, text width=3.2cm, align=center, font=\small},
  arrow/.style={-Stealth, thick, popdark}
]
\node[box, align=center] (b1){病 : Clé \zh{疒} (5 traits)\\Puis \zh{丙} (5 traits)\\Total : 10 traits};
\node[box, right=of b1, align=center] (b2){蚊 : Clé \zh{虫} (6 traits)\\Puis \zh{文} (4 traits)\\Total : 10 traits};
\draw[arrow] (b1.east) -- (b2.west) node[midway, above, font=\tiny, text=popmuted] {Structure gauche-droite};
\node[box, below=of b1, align=center] (b3){Règle : Gauche $\rightarrow$ Droite\\Haut $\rightarrow$ Bas\\Horizontal $\rightarrow$ Vertical};
\node[box, below=of b2, align=center] (b4){Règle : Extérieur $\rightarrow$ Intérieur\\ (pour \zh{疒})\\Puis Fermeture};
\draw[arrow] (b1.south) -- (b3.north);
\draw[arrow] (b2.south) -- (b4.north);
\end{tikzpicture}
\legende{Résumé visuel : ordre des traits et structure pour 病 et 蚊}
\end{popfigure}

\courssub{Compréhension détaillée du texte}

\begin{exercice}[Questions de compréhension — Répondez en français ou en chinois (caractères + pinyin)]
\begin{enumerate}
\item Selon le texte, quel agent pathogène cause le paludisme et comment est-il transmis ?
\item Citez les trois symptômes mentionnés pour une personne atteinte de paludisme.
\item Quels sont les trois conseils de prévention donnés dans le texte ? Résumez-les en français.
\item Traduisez en français : \zh{如果发烧了，要及时看病、吃药。}
\item Trouvez dans le texte l'équivalent chinois de : « éliminer les eaux stagnantes » et « empêcher la prolifération des moustiques ».
\item Le slogan final \zh{预防疟疾，保护健康！} suit une structure parallèle. Identifiez la structure grammaticale de chaque moitié (Verbe + Objet).
\end{enumerate}
\end{exercice}

\begin{corrige}[Exercice — Compréhension détaillée]
\begin{enumerate}
\item Le paludisme est causé par un \textbf{parasite} (\zh{寄生虫} jìshēngchóng / \zh{疟原虫} nüèyuánchóng), transmis par les \textbf{moustiques} (\zh{蚊子} wénzi) via la piqûre (\zh{传播} chuánbō).
\item 1. \zh{发烧} (fāshāo) — Fièvre. 2. \zh{头疼} (tóu téng) — Mal de tête. 3. \zh{全身没有力气} (quánshēn méiyǒu lìqi) — Faiblesse générale / pas de forces dans tout le corps.
\item 1. \textbf{Dormir sous moustiquaire} (\zh{晚上睡觉的时候要用蚊帐}). 2. \textbf{Assainir l'environnement} : le garder propre (\zh{保持环境干净}) et éliminer les eaux stagnantes (\zh{清除积水}) pour empêcher les moustiques de proliférer (\zh{防止蚊子滋生}). 3. \textbf{Consulter vite et se soigner} en cas de fièvre (\zh{如果发烧了，要及时看病、吃药}).
\item « Si tu as de la fièvre, il faut consulter un médecin sans tarder et prendre des médicaments. »
\item \zh{清除积水} (qīngchú jīshuǐ) — « éliminer l'eau stagnante » ; \zh{防止蚊子滋生} (fángzhǐ wénzi zīshēng) — « empêcher les moustiques de proliférer ».
\item \zh{预防} (v.) + \zh{疟疾} (n.) // \zh{保护} (v.) + \zh{健康} (n.). Structure \cle{Verbe + Complément d'objet} (动宾结构) des deux côtés, créant un parallélisme rythmique (对仗).
\end{enumerate}
\end{corrige}

\courssub{Collocations, structures et production guidée}

\begin{aretenir}[Boîte à outils : Collocations et structures utiles pour l'affiche]
\textbf{A. Collocations lexicales (固定搭配) — À mémoriser par cœur}
\begin{center}
\begin{tabularx}{\linewidth}{|X|X|X|}
\hline
\rowcolor{popgreenL} \textbf{Chinois (Caractères + Pinyin)} & \textbf{Structure interne} & \textbf{Traduction} \\
\hline
\zh{用蚊帐} (yòng wénzhàng) & Verbe + Objet & Utiliser une moustiquaire / Dormir sous moustiquaire \\
\hline
\zh{保持环境干净} (bǎochí huánjìng gānjìng) & V + O + Adj (Complément de résultat/état) & Garder l'environnement propre \\
\hline
\zh{清除积水} (qīngchú jīshuǐ) & Verbe + Objet & Éliminer l'eau stagnante \\
\hline
\zh{防止蚊子滋生} (fángzhǐ wénzi zīshēng) & Verbe + Objet + Verbe & Empêcher les moustiques de proliférer \\
\hline
\zh{及时看病} (jíshí kànbìng) & Adv (temps) + Verbe & Consulter un médecin à temps / sans tarder \\
\hline
\zh{吃药} (chī yào) & Verbe + Objet & Prendre des médicaments \\
\hline
\zh{发烧} (fāshāo) & Verbe intransitif & Avoir de la fièvre (\textbf{pas} \zh{有发烧}) \\
\hline
\zh{得了疟疾} (dé le nüèjí) & V + Particule + O & Attraper le paludisme / Être atteint de paludisme \\
\hline
\end{tabularx}
\end{center}

\textbf{B. Structures grammaticales (句法结构) pour l'affiche}
\begin{enumerate}
\item \cle{Obligation / Conseil fort :} \textbf{要  + V} \quad \zh{要用蚊帐。} (Il faut utiliser une moustiquaire.)
\item \cle{Interdiction / Conseil négatif :} \textbf{不要  + V} \quad \zh{不要留积水。} (Ne laissez pas d'eau stagnante.)
\item \cle{Condition - Conséquence :} \textbf{如果 …，就 要 / 要 …} \quad \zh{如果发烧了，就要及时看病。} (Si fièvre $\rightarrow$ consulter vite.) \emph{Note :} \zh{就} (jiù) marque la conséquence, souvent omis à l'oral mais recommandé à l'écrit formel.
\item \cle{Énumération ordonnée :} \textbf{第一，… 第二，… 第三，…} \quad (Premièrement…, Deuxièmement…, Troisièmement…)
\item \cle{Structure parallèle pour slogan :} \textbf{V + O ， V + O !} \quad \zh{预防疟疾，保护健康！} / \zh{人人用蚊帐，家家都健康！}
\end{enumerate}
\end{aretenir}

\begin{attention}[Pièges fréquents pour francophones — \cle{Check-list avant de rendre l'affiche}]
\begin{itemize}
\item \textbf{Tons :} \zh{疟疾} = \textbf{nüèjí} (4\textsuperscript{e} + 2\textsuperscript{e}), \textbf{pas} nuèjí ni ji. \zh{看病} = \textbf{kànbìng} (4\textsuperscript{e} + 4\textsuperscript{e}). \zh{蚊子} = \textbf{wénzi} (2\textsuperscript{e} + ton neutre). \zh{及时} = \textbf{jíshí} (2\textsuperscript{e} + 2\textsuperscript{e}).
\item \textbf{Ordre des mots — Adverbe de temps :} \zh{及时} (à temps) se place \textbf{AVANT} le verbe : \zh{及时看病} $\checkmark$ / \zh{看病及时} $\times$ (sauf structure \zh{看病很及时} « a consulté très vite »).
\item \textbf{« Avoir de la fièvre » :} Verbe \zh{发烧} seul. \zh{他发烧了。} $\checkmark$ — \zh{他有发烧。} $\times$ (calque de « avoir de la fièvre »).
\item \textbf{Classificateur pour moustiquaire :} \zh{一\cle{顶} 蚊帐} (yī \textbf{dǐng} wénzhàng) ou \zh{一\cle{床} 蚊帐} (yī \textbf{chuáng} wénzhàng). Pas \zh{个} (gè) seul.
\item \textbf{Parasite vs Virus :} \zh{寄生虫} (parasite) $\neq$ \zh{病毒} (virus). Le paludisme = parasite.
\item \textbf{Ponctuation liste :} Utilisez le signe \zh{、} (dùnhào) pour énumérer des noms/verbes courts dans une phrase : \zh{发烧、头疼、没力气}. Point final \zh{。} à la fin.
\end{itemize}
\end{attention}

\courssub{Tâche finale : Conception et présentation d'une affiche de sensibilisation (Travail de groupe)}

\begin{methode}[Déroulement de la séance (2 h / 2 périodes)]
\begin{enumerate}
\item \textbf{Période 1 — Préparation écrite (50 min)} \\
\underline{Étape 1 : Analyse du document (10 min).} En groupe de 4, relisez le \texttt{textebox} p. 1. Soulignez : 3 symptômes, 3 mesures de prévention, 5 mots-clés du glossaire. Vérifiez pinyin/ton.
\underline{Étape 2 : Rédaction brouillon (25 min).} Rédigez l'affiche selon le canevas ci-dessous. \textbf{Contrainte :} minimum 10 mots du glossaire (entourez-les). Utilisez les 3 collocations obligatoires : \zh{用蚊帐}, \zh{保持环境干净} (ou \zh{清除积水}), \zh{及时看病}.
\underline{Étape 3 : Calligraphie (15 min).} Choisissez 2 caractères (ex: \zh{病}, \zh{蚊}, \zh{预}, \zh{健}). Tracez-les en grand format (feuillet A4) en respectant l'ordre des traits. Notez la clé à côté.
\item \textbf{Période 2 — Finalisation et Oral (50 min)} \\
\underline{Étape 4 : Mise au net (15 min).} Recopiez l'affiche proprement sur papier affiche (A3 ou A2). Intégrez les 2 caractères calligraphiés. Soignez la mise en page (titre, parties distinctes, slogan).
\underline{Étape 5 : Répétition orale (10 min).} Chaque membre lit une section. Vérifiez les tons entre vous (correction par les pairs).
\underline{Étape 6 : Présentation et Vote (25 min).} Passage des groupes (2-3 min/groupe). Classe note sur grille : Exactitude chinois (5 pts) / Richesse vocabulaire (5 pts) / Prononciation/Tons (5 pts) / Esthétique/Calligraphie (5 pts).
\end{enumerate}
\end{methode}

\begin{experience}[Canevas de l'affiche — À reproduire sur votre brouillon]
\textbf{TITRE} : \underline{\hspace{8cm}} \quad (ex: \zh{预防疟疾，保护健康！} ou \zh{远离蚊子，拥抱健康！})

\vspace{0.5cm}
\textbf{PARTIE 1 : 症状 } (Zhèngzhuàng)
\begin{itemize}
\item 1. \underline{\hspace{7cm}} (ex: \zh{得了疟疾会发烧。})
\item 2. \underline{\hspace{7cm}} (ex: \zh{头疼很厉害。})
\item 3. \underline{\hspace{7cm}} (ex: \zh{全身没有力气。})
\end{itemize}

\vspace{0.5cm}
\textbf{PARTIE 2 : 预防办法 } (Yùfáng bànfǎ)
\begin{itemize}
\item \zh{第一，} \underline{\hspace{6.5cm}} (collocation : \zh{用蚊帐})
\item \zh{第二，} \underline{\hspace{6.5cm}} (collocation : \zh{保持环境干净} / \zh{清除积水})
\item \zh{第三，} \underline{\hspace{6.5cm}} (structure \zh{如果…就要…} + \zh{及时看病})
\end{itemize}

\vspace{0.5cm}
\textbf{SLOGAN FINAL} (2 phrases parallèles V+O) : \underline{\hspace{8cm}} \\
\underline{\hspace{8cm}}

\vspace{0.5cm}
\textbf{CALLIGRAPHIE} (2 caractères grands format + clé + nb traits) :
\begin{itemize}
\item Caractère 1 : \zh{\underline{\hspace{1cm}}} — Clé : \zh{\underline{\hspace{1cm}}} — Traits : \underline{\hspace{1cm}}
\item Caractère 2 : \zh{\underline{\hspace{1cm}}} — Clé : \zh{\underline{\hspace{1cm}}} — Traits : \underline{\hspace{1cm}}
\end{itemize}
\end{experience}

\begin{exoresolu}[Modèle d'affiche complet — Groupe « Santé pour tous » (Corrigé commenté)]
Énoncé : Produire une affiche complète (titre, symptômes, prévention, slogan) contenant $\ge$ 10 mots du glossaire, 3 collocations obligatoires, 2 caractères calligraphiés, structure parallèle pour le slogan.\par
\tcblower
\textbf{AFFICHE MODÈLE}\\[6pt]
\zh{\Large 预防疟疾，保护健康！}\\
Yùfáng nüèjí, bǎohù jiànkāng!\\
« Prévenons le paludisme, protégeons la santé ! »\\[8pt]
\zh{\textbf{症状 }}\\
\zh{得了疟疾的人常常发烧、头疼，全身没有力气。}\\
Dé le nüèjí de rén chángcháng fāshāo, tóu téng, quánshēn méiyǒu lìqi.\\
« Les malades du paludisme ont souvent de la fièvre, mal à la tête, et le corps sans force. »\\[8pt]
\zh{\textbf{预防办法 }}\\
\zh{第一，晚上睡觉要用蚊帐。}\\
Dì-yī, wǎnshang shuìjiào yào yòng wénzhàng.\\
« 1. Le soir, pour dormir, il faut utiliser une moustiquaire. »\\[4pt]
\zh{第二，要保持环境干净，清除积水。}\\
Dì-èr, yào bǎochí huánjìng gānjìng, qīngchú jīshuǐ.\\
« 2. Il faut garder l'environnement propre, éliminer l'eau stagnante. »\\[4pt]
\zh{第三，如果发烧了，就要及时看病、吃药。}\\
Dì-sān, rúguǒ fāshāo le, jiù yào jíshí kànbìng, chī yào.\\
« 3. Si on a de la fièvre, il faut consulter vite un médecin et prendre des médicaments. »\\[8pt]
\zh{\large 人人用蚊帐，家家都健康！}\\
Rénrén yòng wénzhàng, jiājiā dōu jiànkāng!\\
« Que chacun utilise une moustiquaire, et que chaque famille soit en bonne santé ! »\\[8pt]
\textbf{Caractères calligraphiés sur l'affiche :}
\begin{itemize}
\item \zh{\Large 病} — Clé \zh{疒} (nè, maladie) — 10 traits. Ordre : cadre 疒 (ext$\rightarrow$int) puis 丙.
\item \zh{\Large 蚊} — Clé \zh{虫} (chóng, insecte) — 10 traits. Ordre : clé 虫 (gauche, haut$\rightarrow$bas) puis 文 (droite).
\end{itemize}

\textbf{Commentaire linguistique (pour l'enseignant / auto-évaluation) :}
\begin{enumerate}
\item \textbf{Vocabulaire (12/10 min) :} \zh{疟疾、健康、症状、发烧、头疼、力气、预防、蚊帐、环境、干净、积水、及时、看病、吃药、保护}. \textit{Objectif dépassé.}
\item \textbf{Collocations imposées :} \zh{用蚊帐} (L. 9), \zh{保持环境干净} (L. 11), \zh{及时看病} (L. 13). \textit{Intégrées naturellement.}
\item \textbf{Grammaire :} \zh{要} (obligation) $\times$3 ; \zh{如果…就要…} (condition) ; \zh{第一/二/三} (énumération) ; \zh{、} (liste symptômes). \textit{Varié et correct.}
\item \textbf{Slogan :} Parallélisme \zh{人人} (sujet général) + \zh{用蚊帐} (V-O) // \zh{家家} (sujet général) + \zh{都健康} (Adv-V/Adj). Rythmique 4+3 / 4+3 syllabes. \textit{Efficace.}
\item \textbf{Calligraphie :} Choix pertinents (clés transparentes 疒/虫), ordre des traits respecté.
\end{enumerate}
\end{exoresolu}

\courssub{Bilan de la leçon et auto-évaluation}

\begin{aretenir}[Ce que je dois savoir faire après cette leçon (Compétences visées)]
\begin{itemize}
\item \textbf{Lexique :} Reconnaître, prononcer (tons exacts), écrire (pinyin + caractères) et employer les 20 mots du glossaire dans un contexte de santé.
\item \textbf{Lecture :} Comprendre un texte court (100-150 caractères) sur le paludisme au Cameroun, en extraire symptômes, causes, prévention.
\item \textbf{Grammaire :} Construire des conseils de santé avec \zh{要/不要}, \zh{如果…就…}, \zh{第一…} ; placer correctement \zh{及时}.
\item \textbf{Écriture :} Tracer correctement \zh{病} (clé 疒) et \zh{蚊} (clé 虫) + 2 autres caractères du thème ; respecter l'ordre des traits.
\item \textbf{Expression orale :} Lire l'affiche à voix haute avec tons corrects, débit fluide, en respectant les groupes rythmiques (ex: \zh{预防疟疾 / 保护健康}).
\item \textbf{Culture :} Savoir expliquer brièvement la coopération Cameroun-Chine dans la lutte anti-paludique (centres de traitement, dons de moustiquaires, artémisinine \zh{青蒿素} découverte par Tu Youyou \zh{屠呦呦}, prix Nobel 2015).
\end{itemize}
\end{aretenir}

\begin{savaistu}[Le saviez-vous ? — Chine / Cameroun : Lutte contre le paludisme]
\begin{itemize}
\item L'\textbf{artémisinine} (\zh{青蒿素} qīnghāosù), molécule clé des traitements antipaludiques actuels (ACT), a été isolée en 1972 par la pharmacologue chinoise \textbf{Tu Youyou} (\zh{屠呦呦}), prix Nobel de physiologie/médecine 2015. Elle s'est inspirée de la pharmacopée traditionnelle (\zh{本草纲目}).
\item La Chine a été certifiée \textbf{« exempte de paludisme »} par l'OMS en \textbf{juin 2021} (zéro cas autochtone depuis 2017). Elle partage son expérience (stratégie « 1-3-7 » : déclaration sous 1 jour, investigation sous 3 jours, riposte sous 7 jours) avec les pays africains, dont le Cameroun.
\item Au Cameroun, le \textbf{Programme National de Lutte contre le Paludisme (PNLP)} distribue gratuitement des moustiquaires imprégnées (MILDA) lors de campagnes de masse (souvent avec appui du Fonds mondial, de la Chine, des USA). Le vaccin RTS,S/AS01 (Mosquirix) a commencé à être déployé en 2024 chez les jeunes enfants.
\end{itemize}
\end{savaistu}

\begin{exercice}[Entraînement individuel — À faire au cahier pour la prochaine évaluation]
\begin{enumerate}
\item \textbf{Traduction (Version) :} Traduisez en français.
\begin{enumerate}
\item 晚上睡觉一定要用蚊帐，防止蚊子叮咬。
\item 如果发现发烧、头疼，要及时去医院看病，不要自己随便吃药。
\item 保持环境干净，清除积水，是预防疟疾最重要的办法。
\end{enumerate}
\item \textbf{Traduction (Thème) :} Traduisez en chinois (caractères + pinyin).
\begin{enumerate}
\item Le paludisme est une maladie courante transmise par les moustiques.
\item Les symptômes sont : fièvre, mal de tête, pas de force.
\item Il faut garder la cour propre et ne pas laisser d'eau stagnante.
\item En cas de fièvre, allez vite à l'hôpital voir le médecin.
\end{enumerate}
\item \textbf{Ordre des traits :} Écrivez l'ordre des traits (numérotés 1 à 10) pour \zh{疟} (clé 疒 + 略) et \zh{预} (haut + 页). Indiquez la clé.
\item \textbf{Complétez avec la bonne collocation :} \zh{用蚊帐 / 保持环境干净 / 及时看病 / 清除积水 / 吃药}
\begin{enumerate}
\item 晚上睡觉前，请 \underline{\hspace{3cm}}。
\item 雨季来了，院子里有很多水，要赶紧 \underline{\hspace{3cm}}。
\item 生病了不舒服，必须 \underline{\hspace{3cm}}，听医生的话。
\item 为了健康，我们要 \underline{\hspace{3cm}}，不乱扔垃圾。
\item 医生开了三天的药，按时 \underline{\hspace{3cm}}。
\end{enumerate}
\end{enumerate}
\end{exercice}

\begin{corrige}[Exercice — Entraînement individuel]
\begin{enumerate}
\item \textbf{Version :}
\begin{enumerate}
\item « Le soir en dormant, il faut absolument utiliser une moustiquaire pour empêcher les moustiques de piquer. »
\item « Si l'on découvre fièvre et mal de tête, il faut aller à l'hôpital consulter sans tarder, ne pas prendre de médicaments n'importe comment soi-même. »
\item « Garder l'environnement propre, éliminer l'eau stagnante, c'est le moyen le plus important de prévenir le paludisme. »
\end{enumerate}
\item \textbf{Thème :}
\begin{enumerate}
\item \zh{疟疾是蚊子传播的常见疾病。} (Nüèjí shì wénzi chuánbō de chángjiàn jíbìng.)
\item \zh{症状是：发烧、头疼、没有力气。} (Zhèngzhuàng shì: fāshāo, tóu téng, méiyǒu lìqi.)
\item \zh{要保持院子干净，不要留积水。} (Yào bǎochí yuànzi gānjìng, bú yào liú jīshuǐ.)
\item \zh{如果发烧了，要赶紧去医院看病。} (Rúguǒ fāshāo le, yào gǎnjǐn qù yīyuàn kànbìng.)
\end{enumerate}
\item \textbf{Ordre des traits :}
\begin{itemize}
\item \zh{疟} (9 traits) : Clé \zh{疒} (traits 1-5 : 丶丶一丨㇆) + \zh{略} (traits 6-9 : 丨一丨一 — \emph{simplifié : 略 s'écrit 纟+各, ici partie phonétique simplifiée}). \textit{Note : 疟 est HSK 5+, reconnaissance prioritaire.}
\item \zh{预} (16 traits) : Partie haute \zh{予} (traits 1-4 : 丶一丨㇟) + \zh{ㄗ} (traits 5-6 : 丨一) + Clé \zh{页} (traits 7-16 : 一丨一一一丨㇆一一一). Règle : Haut avant bas.
\end{itemize}
\item \textbf{Collocations :}
\begin{enumerate}
\item \zh{用蚊帐}
\item \zh{清除积水}
\item \zh{及时看病}
\item \zh{保持环境干净}
\item \zh{吃药}
\end{enumerate}
\end{enumerate}
\end{corrige}

\newpage

\coursec{Méthodes et exercices résolus}

\coursec{Leçon 12 — Vocabulaire et compréhension de texte : le paludisme}

\courssub{Mise en route et texte support}

\begin{methode}[Lire un texte de santé publique en chinois]
Un texte sur le paludisme (疟疾 nǜèjí) suit presque toujours la même logique : présentation de la maladie, causes (transmission), symptômes, prévention et prise en charge. Voici la démarche à adopter pour une lecture efficace.
\begin{enumerate}
\item \textbf{Lecture globale sans dictionnaire.} Lis le texte une première fois en repérant les mots-clés du thème : \zh{病} (maladie), \zh{发烧} (fièvre), \zh{蚊子} (moustique), \zh{医院} (hôpital), \zh{预防} (prévenir). Ne t'arrête pas sur un caractère inconnu.
\item \textbf{Repérer la structure logique et les connecteurs.} Souligne les marqueurs de relation : \zh{因为……所以……} (parce que... donc...), \zh{如果……就……} (si... alors...), \zh{为了……应该……} (pour... il faut...), \zh{被} (marque du passif/agent). Ils portent l'argumentation.
\item \textbf{Deviner par le contexte et les radicaux (clés).} Un caractère avec la clé \zh{疒} (nè, « maladie ») désigne presque toujours une pathologie ou un symptôme : \zh{疟} (paludisme), \zh{疾} (maladie grave), \zh{疼} (douleur), \zh{痛} (mal). La clé \zh{虫} (chóng, « insecte ») signale un insecte ou reptile : \zh{蚊} (moustique), \zh{蝇} (mouche).
\item \textbf{Identifier les phrases de conseil et d'obligation.} Les recommandations sanitaires utilisent \zh{应该} / \zh{要} / \zh{得 (děi)} + verbe : \zh{睡觉的时候要挂蚊帐。} (Shuìjiào de shíhou yào guà wénzhàng.) « Au moment de dormir, il faut suspendre une moustiquaire. »
\item \textbf{Reformuler en français.} Traduis phrase par phrase en respectant l'ordre chinois : Sujet + Temps/Lieu + Verbe + Complément.
\end{enumerate}
\end{methode}

\begin{textebox}[Le paludisme au Cameroun — Texte support original]
\zh{在喀麦隆，疟疾是一种很常见的病。人被传播疟疾的蚊子咬了以后，会发烧、头疼，有时候还会呕吐。为了预防疟疾，我们应该挂蚊帐，保持环境干净，生病了要马上去医院看病。}\\
\emph{Zài Kāmàilóng, nǜèjí shì yì zhǒng hěn chángjiàn de bìng. Rén bèi chuánbō nǜèjí de wénzi yǎo le yǐhòu, huì fāshāo, tóu téng, yǒu shíhou hái huì ǒutù. Wèile yùfáng nǜèjí, wǒmen yīnggāi guà wénzhàng, bǎochí huánjìng gānjìng, shēngbìng le yào mǎshàng qù yīyuàn kànbìng.}\\
« Au Cameroun, le paludisme est une maladie très courante. Après avoir été piqué par un moustique transmettant le paludisme, on a de la fièvre, mal à la tête, et parfois on vomit. Pour prévenir le paludisme, nous devons utiliser des moustiquaires, garder l'environnement propre, et si l'on tombe malade, aller immédiatement consulter à l'hôpital. »
\end{textebox}

\begin{exoresolu}[Compréhension globale du texte]
\textbf{Énoncé.} Lis le texte ci-dessus et réponds en français :
\begin{enumerate}
\item Comment attrape-t-on le paludisme ? (Transmission)
\item Quels sont les trois symptômes cités ? (Symptômes)
\item Que faut-il faire pour se protéger et que faire si l'on est malade ? (Prévention et soin)
\end{enumerate}
\tcblower
\textbf{Solution détaillée.}
\begin{enumerate}
\item \textbf{Transmission.} La phrase clé est \zh{人被传播疟疾的蚊子咬了以后}. Structure passive : \zh{被} + agent (\zh{传播疟疾的蚊子}) + verbe (\zh{咬}) + \zh{了} (accompli) + \zh{以后} (après). Réponse : On attrape le paludisme quand on est piqué par un moustique qui transmet le paludisme (porteur du parasite).
\item \textbf{Symptômes.} Le verbe modal \zh{会} (risquer de / avoir comme conséquence) introduit la liste : \zh{发烧} (fièvre), \zh{头疼} (mal de tête), et avec \zh{有时候还会} (parfois aussi) : \zh{呕吐} (vomir). Réponse : La fièvre, les maux de tête et parfois les vomissements.
\item \textbf{Prévention et soin.} \zh{为了} + but (\zh{预防疟疾}) introduit les conseils avec \zh{应该} et \zh{要} : 1. \zh{挂蚊帐} (installer une moustiquaire), 2. \zh{保持环境干净} (garder l'environnement propre), 3. \zh{生病了要马上去医院看病} (si malade, aller tout de suite à l'hôpital consulter). Note : \zh{看病} = « regarder la maladie » $\rightarrow$ « consulter un médecin ».
\end{enumerate}
\textbf{Conclusion.} La méthode gagnante : repérer les marqueurs logiques (\zh{被}, \zh{会}, \zh{为了}, \zh{应该}) avant de traduire mot à mot.
\end{exoresolu}

\courssub{Glossaire du thème : vocabulaire clé}

\begin{aretenir}[Tableau de vocabulaire — Module « Environnement et santé : Paludisme »]
\begin{center}
\begin{tabularx}{\linewidth}{|X|X|X|X|}
\hline
\rowcolor{popblueL} \textbf{Caractères} & \textbf{Pinyin} & \textbf{Nature} & \textbf{Traduction} \\ \hline
疟疾 & nǜèjí & n. & paludisme (malaria) \\ \hline
病毒 & bìngdú & n. & virus (attention : le paludisme est causé par un parasite 疟原虫 nǜèyuánchóng, mais \zh{病毒} reste fréquent dans le langage courant pour « agent pathogène ») \\ \hline
蚊子 & wénzi & n. & moustique \\ \hline
叮咬 & dīngyǎo & v. & piquer (insecte) — plus précis que \zh{咬} pour les moustiques \\ \hline
传播 & chuánbō & v. & transmettre, propager (maladie) \\ \hline
发烧 & fāshāo & v. & avoir de la fièvre (collocation : \zh{发高烧} = forte fièvre) \\ \hline
头疼 / 头痛 & tóu téng / tóu tòng & v./n. & mal de tête / céphalée \\ \hline
呕吐 & ǒutù & v. & vomir \\ \hline
症状 & zhèngzhuàng & n. & symptôme \\ \hline
预防 & yùfáng & v. & prévenir, prévention \\ \hline
蚊帐 & wénzhàng & n. & moustiquaire (clé \zh{巾} jīn, tissu) \\ \hline
挂 & guà & v. & suspendre, accrocher (collocation : \zh{挂蚊帐}) \\ \hline
保持 & bǎochí & v. & garder, maintenir (état) — collocation : \zh{保持干净 / 保持卫生} \\ \hline
环境 & huánjìng & n. & environnement, cadre de vie \\ \hline
干净 & gānjìng & adj. & propre, clean \\ \hline
卫生 & wèishēng & n./adj. & hygiène / hygiénique \\ \hline
医院 & yīyuàn & n. & hôpital \\ \hline
医生 & yīshēng & n. & médecin \\ \hline
看病 / 看医生 & kànbìng / kàn yīshēng & v. & consulter (un médecin) \\ \hline
吃药 & chī yào & v. & prendre un médicament (litt. « manger le médicament ») \\ \hline
药 & yào & n. & médicament, remède (ton 4) — ne pas confondre avec \zh{要} yào (vouloir/falloir) \\ \hline
马上 & mǎshàng & adv. & immédiatement, tout de suite \\ \hline
得 & dé / děi & v. & auxiliaire : \zh{得} dé (devoir, obligation forte, oral) / \zh{děi} (falloir, obligation) \\ \hline
\end{tabularx}
\end{center}
\end{aretenir}

\begin{methode}[Mémoriser et employer le lexique de la santé]
\begin{enumerate}
\item \textbf{Apprendre par familles sémantiques (champs lexicaux).} Regroupe : \textbf{Maladie/agent} : \zh{疟疾, 疟原虫, 病毒, 细菌, 传播} ; \textbf{Symptômes} : \zh{发烧, 头疼, 呕吐, 咳嗽, 流鼻涕, 腹泻} ; \textbf{Prévention} : \zh{预防, 蚊帐, 挂, 干净, 卫生, 环境, 驱蚊液} ; \textbf{Soins} : \zh{医院, 医生, 看病, 吃药, 打针, 康复}.
\item \textbf{Apprendre le mot avec son verbe (collocation).} En chinois, le nom seul ne suffit pas : \zh{发烧} (avoir de la fièvre), \zh{挂蚊帐} (installer une moustiquaire), \zh{吃药} (prendre un médicament), \zh{打针} (faire une piqûre/injection), \zh{看医生} (voir le médecin).
\item \textbf{Soigner les tons dès le départ.} Distingue \zh{药} yào (médicament, 4\up{e} ton) de \zh{要} yào (vouloir/falloir, 4\up{e} ton aussi mais sens différent, le contexte tranche) ; \zh{疟} nǜè (4\up{e} ton) vs \zh{虐} nǜè (maltraiter) ; \zh{疾} jí (2\up{e} ton, maladie grave) vs \zh{及} jí (atteindre). Attention à \zh{病} bìng (4\up{e} ton) et \zh{疼} téng (2\up{e} ton).
\item \textbf{Vérifier par une phrase complète contextualisée.} Pour chaque mot, fabrique une phrase simple : \zh{我弟弟发高烧了，妈妈马上带他去医院看医生。} (Wǒ dìdi fā gāoshāo le, māma mǎshàng dài tā qù yīyuàn kàn yīshēng.) « Mon petit frère a une forte fièvre, maman l'emmène tout de suite à l'hôpital voir le médecin. »
\end{enumerate}
\end{methode}

\begin{exoresolu}[Vocabulaire : phrases à compléter]
\textbf{Énoncé.} Complète avec le mot qui convient parmi : \zh{蚊帐 / 发烧 / 预防 / 呕吐 / 医院 / 挂 / 保持 / 吃药}.
\begin{enumerate}
\item \zh{他昨天晚上吃了不干净的东西，今天一直＿＿。}
\item \zh{睡觉以前，妈妈把＿＿＿＿好了。}
\item \zh{孩子＿＿到三十九度，得马上去＿＿。}
\item \zh{打疫苗可以＿＿很多传染病。}
\item \zh{医生说：＿＿环境卫生，少生病。}
\item \zh{感冒了要多喝水，按时＿＿。}
\end{enumerate}
\tcblower
\textbf{Solution détaillée.}
\begin{enumerate}
\item \textbf{呕吐} (ǒutù). Indice : « manger quelque chose de pas propre » $\rightarrow$ vomissements ; \zh{一直} (sans cesse) + verbe d'action. \emph{Tā zuótiān wǎnshàng chī le bù gānjìng de dōngxi, jīntiān yīzhí ǒutù.}
\item \textbf{蚊帐 / 挂} (wénzhàng / guà). Collocation obligatoire : \zh{挂蚊帐} ; \zh{挂好} + \zh{了} = résultat accompli. \emph{Shuìjiào yǐqián, māma bǎ wénzhàng guà hǎo le.}
\item \textbf{发烧 ; 医院} (fāshāo ; yīyuàn). \zh{发烧到} + température ; \zh{得} (děi, il faut) + \zh{马上去医院}. \emph{Háizi fāshāo dào sānshíjiǔ dù, děi mǎshàng qù yīyuàn.}
\item \textbf{预防} (yùfáng). Structure : \zh{可以} + V + CO ; \zh{打疫苗} (se faire vacciner) \zh{预防传染病}. \emph{Dǎ yìmiáo kěyǐ yùfáng hěnduō chuánrǎnbìng.}
\item \textbf{保持} (bǎochí). Collocation : \zh{保持环境卫生} / \zh{保持干净}. \emph{Yīshēng shuō: bǎochí huánjìng wèishēng, shǎo shēngbìng.}
\item \textbf{吃药} (chī yào). Collocation standard pour « prendre un médicament ». \emph{Gǎnmào le yào duō hē shuǐ, ànshí chī yào.}
\end{enumerate}
\textbf{Conclusion.} Le bon réflexe : chercher d'abord la \cle{collocation} (quel verbe va avec quel nom), puis vérifier le sens global de la phrase.
\end{exoresolu}

\courssub{Clés (radicaux) et ordre des traits}

\begin{propriete}[Analyse des radicaux (clés) et ordre des traits — Caractères du thème]
La connaissance des clés (部首 bùshǒu) permet de deviner le sens et de retrouver le caractère dans un dictionnaire. L'ordre des traits (笔顺 bǐshùn) est indispensable pour écrire correctement et vite (règle générale : haut $\rightarrow$ bas, gauche $\rightarrow$ droite, horizontal avant vertical, trait central avant les ailes, cadre avant contenu, fermeture en dernier).
\end{propriete}

\begin{exemplebox}[Radicaux clés du thème « Paludisme »]
\begin{center}
\begin{tabularx}{\linewidth}{|X|X|X|X|}
\hline
\rowcolor{popgreenL} \textbf{Caractère} & \textbf{Clé (Radical)} & \textbf{Sens de la clé} & \textbf{Ordre des traits (résumé)} \\ \hline
\zh{疟} (nǜè) & \zh{疒} (nè, bìngzìtóu) & Maladie, symptôme & 1. Points haut (丶丶) 2. Horizontal (一) 3. Vertical (丨) 4. Enveloppement (冂) 5. Intérieur \zh{勺} \\ \hline
\zh{疾} (jí) & \zh{疒} (nè) & Maladie grave & 1. Clé \zh{疒} (5 traits) 2. \zh{矢} (shǐ, flèche : 横、竖、横、撇、捺) \\ \hline
\zh{蚊} (wén) & \zh{虫} (chóng, Chongzi di) & Insecte, reptile & 1. Haut \zh{文} (wén : 点、横、撇、捺、横) 2. Bas \zh{虫} (竖、横折、横、竖、横) \\ \hline
\zh{帐} (zhàng) & \zh{巾} (jīn, Jinzi pang) & Tissu, étoffe & 1. Clé \zh{巾} (3 traits : 竖、横折钩、横) 2. Droite \zh{长} (cháng : 撇、横、竖、横) \\ \hline
\zh{预} (yù) & \zh{页} (yè, Yezipang) & Tête, page & 1. Haut \zh{ㄨ} (横折、横) 2. \zh{页} (横、竖、横折、横、横、横) \\ \hline
\zh{防} (fáng) & \zh{阝} (fù, Zuoerdao) & Monticule, rempart (à gauche) & 1. Haut \zh{方} (fāng : 点、横、撇、捺、横) 2. Clé \zh{阝} (横折弯钩、竖) \\ \hline
\zh{呕} (ǒu) & \zh{口} (kǒu, Kouzipang) & Bouche, action buccale & 1. Clé \zh{口} (竖、横折、横) 2. Droite \zh{区} (qū : ㄈ、横、竖、横) \\ \hline
\zh{吐} (tù) & \zh{口} (kǒu) & Bouche, action buccale & 1. Clé \zh{口} 2. Droite \zh{土} (tǔ : 横、竖、横) \\ \hline
\end{tabularx}
\end{center}
\end{exemplebox}

\begin{popfigure}
\begin{tikzpicture}[
  mindmap,
  grow cyclic,
  every node/.style={concept, execute at begin node=\small},
  concept color=popblue!60,
  level 1/.append style={level distance=3.5cm, sibling angle=90, concept color=popblue!40},
  level 2/.append style={level distance=2.8cm, sibling angle=45, concept color=popgreen!40},
  level 3/.append style={level distance=2.2cm, sibling angle=30, concept color=poporange!40},
  text=popdark
]
\node {Clés du thème}
  child { node {疒 (Maladie)}
    child { node {疟 nǜè} }
    child { node {疾 jí} }
    child { node {疼 téng} }
    child { node {痛 tòng} }
  }
  child { node {虫 (Insecte)}
    child { node {蚊 wén} }
    child { node {蝇 yíng} }
  }
  child { node {巾 (Tissu)}
    child { node {帐 zhàng} }
    child { node {帽 mào} }
  }
  child { node {口 (Bouche)}
    child { node {呕 ǒu} }
    child { node {吐 tù} }
    child { node {咳 ké} }
  };
\end{tikzpicture}
\legende{Carte mentale des radicaux (clés) fréquents dans le vocabulaire de la santé.}
\end{popfigure}

\begin{exercice}[Entraînement à l'écriture et reconnaissance des clés]
\begin{enumerate}
\item Écris trois fois chaque caractère ci-dessous en respectant l'ordre des traits : \zh{疟}、\zh{疾}、\zh{蚊}、\zh{帐}、\zh{预}、\zh{防}、\zh{呕}、\zh{吐}.
\item Pour chaque caractère, entoure la clé (radical) et donne son nom en français.
\item Classifie les mots suivants selon leur clé : \zh{疼、痛、蝇、帽、咳、唇、痒、疗}. (Utilise un tableau : Caractère | Clé | Sens de la clé).
\end{enumerate}
\end{exercice}

\begin{corrige}[Exercice : Clés et ordre des traits]
\begin{enumerate}
\item \textit{Pratique individuelle sur cahier d'écriture (田字格). Vérifier :} \zh{疟} (clé \zh{疒} 5 traits + \zh{勺}) ; \zh{疾} (clé \zh{疒} + \zh{矢}) ; \zh{蚊} (haut \zh{文} + bas \zh{虫}) ; \zh{帐} (clé \zh{巾} à gauche + \zh{长} à droite) ; \zh{预} (haut \zh{ㄨ} + bas \zh{页}) ; \zh{防} (haut \zh{方} + clé \zh{阝} gauche) ; \zh{呕/吐} (clé \zh{口} gauche + phonétique droite).
\item Classification :
\begin{center}
\begin{tabularx}{\linewidth}{|X|X|X|}
\hline
\rowcolor{popblueL} Caractère & Clé & Sens de la clé \\ \hline
\zh{疼、痛、痒、疗} & \zh{疒} (bìngzìtóu) & Maladie, symptôme, soin \\ \hline
\zh{蝇} & \zh{虫} (chóng) & Insecte \\ \hline
\zh{帽} & \zh{巾} (jīn) & Tissu, vêtement \\ \hline
\zh{咳、唇} & \zh{口} (kǒu) & Bouche, organe buccal \\ \hline
\end{tabularx}
\end{center}
\end{enumerate}
\end{corrige}

\courssub{Compréhension détaillée du texte}

\begin{methode}[Rédiger une réponse de compréhension en chinois (niveau Bac)]
\begin{enumerate}
\item \textbf{Repérer la phrase source.} Les questions avec \zh{为什么} (pourquoi) appellent une réponse avec \zh{因为……} ; les questions avec \zh{怎么 / 怎样} (comment) appellent \zh{应该 / 要 / 得} + verbe ; les questions \zh{什么} (quoi) portent sur le COD.
\item \textbf{Reprendre le vocabulaire du texte, adapter la personne.} Si la question porte sur « il/elle » (\zh{他/她}), la réponse garde la 3\up{e} personne. Ne recopie jamais la question entière.
\item \textbf{Construire une phrase complète, ponctuée.} Sujet + (Temps/Lieu) + Verbe + Complément + \zh{。}. Jamais de réponse par un seul mot au Bac.
\item \textbf{Vérifier la place des compléments (Règle d'or).} Complément de TEMPS et de LIEU se placent \textbf{AVANT} le verbe : \zh{他昨天去医院看病了。} (Lui hier aller hôpital consulter) $\neq$ *\zh{他去医院看病昨天}.
\end{enumerate}
\end{methode}

\begin{exoresolu}[Questions de compréhension — Réponses en chinois]
\textbf{Énoncé.} Réponds en chinois, par des phrases complètes, d'après le texte « Le paludisme au Cameroun ».
\begin{enumerate}
\item \zh{人怎么会得疟疾？} (Comment attrape-t-on le paludisme ?)
\item \zh{得了疟疾有什么症状？} (Quels symptômes le paludisme provoque-t-il ?)
\item \zh{为了预防疟疾，我们应该做什么？} (Pour prévenir le paludisme, que devons-nous faire ?)
\item \zh{生病了以后应该怎么办？} (Une fois malade, que faut-il faire ?)
\end{enumerate}
\tcblower
\textbf{Solution détaillée.}
\begin{enumerate}
\item Question \zh{怎么} (comment). Réponse reprenant la structure passive \zh{被} : \zh{人被传播疟疾的蚊子咬了以后，就会得疟疾。} \emph{(Rén bèi chuánbō nǜèjí de wénzi yǎo le yǐhòu, jiù huì dé nǜèjí.)} « Après avoir été piqué par un moustique transmettant le paludisme, on attrape le paludisme. » Justification : \zh{被} introduit l'agent, \zh{了以后} marque l'antériorité, \zh{就} relie cause et conséquence.
\item Question \zh{什么症状}. Réponse énumérée avec \zh{、} et \zh{还会} : \zh{得了疟疾会发烧、头疼，有时候还会呕吐。} \emph{(Dé le nǜèjí huì fāshāo, tóu téng, yǒu shíhou hái huì ǒutù.)} « Quand on a le paludisme, on a de la fièvre, mal à la tête, et parfois on vomit. »
\item Question \zh{应该做什么}. Réponse listant les 3 mesures avec \zh{应该} et \zh{要} : \zh{我们应该挂蚊帐，保持环境干净，生病了要马上去医院看病。} \emph{(Wǒmen yīnggāi guà wénzhàng, bǎochí huánjìng gānjìng, shēngbìng le yào mǎshàng qù yīyuàn kànbìng.)} La question utilise \zh{我们}, la réponse garde \zh{我们}.
\item Question \zh{怎么办}. Réponse centrée sur l'action immédiate : \zh{生病了要马上去医院看病，听医生的话，按时吃药。} \emph{(Shēngbìng le yào mǎshàng qù yīyuàn kànbìng, tīng yīshēng de huà, ànshí chī yào.)} « Une fois malade, il faut aller tout de suite à l'hôpital, écouter le médecin, prendre les médicaments à l'heure. »
\end{enumerate}
\textbf{Conclusion.} Une bonne réponse = mots du texte + structure complète (SVO + compléments bien placés) + personne grammaticale adaptée + ponctuation chinoise。
\end{exoresolu}

\begin{exercice}[Compréhension écrite — Nouveaux items]
\textbf{Énoncé.} D'après le texte, réponds en chinois par des phrases complètes.
\begin{enumerate}
\item \zh{疟疾在喀麦隆常见吗？}
\item \zh{蚊子咬人以后，人会马上发烧吗？(注意文中“以后”的用法)}
\item \zh{为什么要挂蚊帐？}
\item \zh{如果你发烧、头疼，你应该怎么做？}
\end{enumerate}
\end{exercice}

\begin{corrige}[Exercice : Compréhension écrite]
\begin{enumerate}
\item \zh{是的，疟疾在喀麦隆是一种很常见的病。} (Shì de, nǜèjí zài Kāmàilóng shì yì zhǒng hěn chángjiàn de bìng.)
\item \zh{文中没说“马上”。说“被蚊子咬了以后，会发烧……”，意思是过一段时间后出现症状。} (Wénzhōng méi shuō "mǎshàng". Shuō "bèi wénzi yǎo le yǐhòu, huì fāshāo...", yìsi shì guò yì duàn shíjiān hòu chūxiàn zhèngzhuàng.) \textit{Note : "以后" indique l'antériorité, pas l'immédiateté.}
\item \zh{为了预防疟疾，应该挂蚊帐，防止蚊子咬人。} (Wèile yùfáng nǜèjí, yīnggāi guà wénzhàng, fángzhǐ wénzi yǎo rén.)
\item \zh{我应该马上去医院看病，听医生的话，按时吃药。} (Wǒ yīnggāi mǎshàng qù yīyuàn kànbìng, tīng yīshēng de huà, ànshí chī yào.)
\end{enumerate}
\end{corrige}

\courssub{Collocations et production guidée}

\begin{methode}[Traduire du français vers le chinois (Thème) sans calquer]
\begin{enumerate}
\item \textbf{Identifier le noyau de la phrase :} Sujet + Verbe principal.
\item \textbf{Placer les compléments :} TEMPS (hier, maintenant) et LIEU (à l'hôpital, à la maison) \textbf{AVANT} le verbe. \zh{我昨天在医院看病。}
\item \textbf{Traduire « il faut / il est nécessaire » :} \zh{应该} (conseil moral/général), \zh{要} (nécessité objective/forte), \zh{得 (děi)} (obligation contrainte, oral).
\item \textbf{Rendre le passif français par \zh{被} (bèi) :} « Être piqué par un moustique » $\rightarrow$ \zh{被蚊子咬} / \zh{被蚊子叮咬}. Sujet = victime.
\item \textbf{Utiliser \zh{了} (le) pour :} 1. Action accomplie (\zh{咬了}) 2. Changement d'état / nouvelle situation (\zh{发烧了} « il a maintenant de la fièvre »).
\item \textbf{Utiliser les collocations chinoises (pas le calque) :} « Prendre un médicament » = \zh{吃药} ; « Avoir de la fièvre » = \zh{发烧} ; « Consulter » = \zh{看病 / 看医生} ; « Garder propre » = \zh{保持……干净/卫生} ; « Installer moustiquaire » = \zh{挂蚊帐}.
\item \textbf{Relire : tons, caractères proches, ordre des mots.}
\end{enumerate}
\end{methode}

\begin{exoresolu}[Thème : traduction de phrases (Français $\rightarrow$ Chinois)]
\textbf{Énoncé.} Traduis en chinois (caractères + pinyin).
\begin{enumerate}
\item Mon petit frère a été piqué par un moustique, maintenant il a de la fièvre.
\item Pour prévenir le paludisme, il faut garder la maison propre.
\item Hier, il est allé à l'hôpital voir le médecin parce qu'il vomissait.
\item Le médecin lui a dit : « Il faut prendre le médicament à l'heure et boire plus d'eau. »
\end{enumerate}
\tcblower
\textbf{Solution détaillée.}
\begin{enumerate}
\item Sujet \zh{我弟弟} ; Passif \zh{被蚊子咬了} (ou \zh{叮咬了}) ; Conséquence \zh{现在} + \zh{发烧了} (changement d'état). \zh{我弟弟被蚊子咬了，现在发烧了。} \emph{(Wǒ dìdi bèi wénzi yǎo le, xiànzài fāshāo le.)}
\item « Pour + infinitif » = \zh{为了} + V ; « il faut » = \zh{要/应该} ; « garder propre » = \zh{保持……干净}. \zh{为了预防疟疾，我们要保持家里干净。} \emph{(Wèile yùfáng nǜèjí, wǒmen yào bǎochí jiā lǐ gānjìng.)} (Pas *\zh{让家干净}).
\item Temps \zh{昨天} ; Lieu \zh{在医院} ; Cause \zh{因为……所以……} ; Verbe \zh{呕吐}. \zh{昨天他因为呕吐，去医院看病了。} \emph{(Zuótiān tā yīnwèi ǒutù, qù yīyuàn kànbìng le.)} Ou : \zh{昨天他呕吐了，所以去医院看病。}
\item Discours rapporté direct : \zh{医生对他说：“你要按时吃药，多喝水。”} \emph{(Yīshēng duì tā shuō: "Nǐ yào ànshí chī yào, duō hē shuǐ.")} \textit{Note :} \zh{按时} (à l'heure/régulièrement), \zh{多喝水} (boire plus d'eau / beaucoup d'eau).
\end{enumerate}
\textbf{Conclusion.} La traduction réussie repose sur trois contrôles : 1. Ordre des mots (Temps/Lieu avant V), 2. Collocations (\zh{保持干净}, \zh{挂蚊帐}, \zh{吃药}), 3. Marqueurs aspectuels (\zh{了}, \zh{被}, \zh{过}).
\end{exoresolu}

\begin{experience}[Production orale guidée : Jeu de rôle « À la consultation »]
\textbf{Contexte :} Un élève joue le \textbf{patient} (symptômes : fièvre 39°C, mal de tête, vomissements), l'autre le \textbf{médecin} (diagnostic : paludisme probable, prescription : prise de sang, médicaments, repos, moustiquaire).
\textbf{Consignes :}
\begin{enumerate}
\item Patient : Décrire les symptômes avec \zh{我……了} (changement d'état) et \zh{很/特别} (intensité).
\item Médecin : Poser des questions avec \zh{有没有……} / \zh{多久了} ; Donner conseils avec \zh{你应该……} / \zh{别……} / \zh{要……}.
\item Utiliser le vocabulaire de la leçon : \zh{发烧、头疼、呕吐、验血、吃药、挂蚊帐、保持环境卫生}.
\end{enumerate}
\textbf{Modèle de dialogue :}
\zh{病人：医生，我发烧了，头疼，还想吐。}\\
\zh{医生：你发烧多久了？有没有呕吐？}\\
\zh{病人：两天了，吐了两次。}\\
\zh{医生：可能是疟疾。我给你开单子验血。你要按时吃药，多休息，回家挂蚊帐。}\\
\zh{病人：好的，谢谢医生。}
\end{experience}

\begin{exercice}[Expression écrite : Message de prévention]
\textbf{Consigne.} Tu es délégué santé de ton lycée. Rédige une \textbf{affiche} (通知 tōngzhī) en chinois (80-100 caractères) pour sensibiliser les élèves au paludisme. Inclus : 1. Danger (paludisme fréquent), 2. Cause (moustique), 3. Symptômes (3), 4. 3 gestes prévention, 5. Conseil si malade. Utilise des phrases impératives (\zh{要、别、应该}).
\end{exercice}

\begin{corrige}[Exercice : Expression écrite — Affiche de prévention]
\textbf{Modèle de réponse (environ 90 caractères) :}
\begin{textebox}[通知：预防疟疾，人人有责]
\zh{同学们：}\\
\zh{疟疾在喀麦隆很常见，是蚊子传播的。得病会发烧、头疼、呕吐，很危险。}\\
\zh{请大家注意：}\\
\zh{1. 睡觉要挂蚊帐。}\\
\zh{2. 保持宿舍干净，别留积水。}\\
\zh{3. 如果发烧、头疼，马上去医院看病，别自己买药。}\\
\zh{健康最重要！}\\
\zh{卫生委员 敬上}
\end{textebox}
\textbf{Traduction :} « Camarades : Le paludisme est très courant au Cameroun, transmis par les moustiques. La maladie donne fièvre, maux de tête, vomissements, c'est dangereux. Attention : 1. Dormir sous moustiquaire. 2. Garder le dortoir propre, pas d'eau stagnante. 3. Si fièvre/maux de tête, aller à l'hôpital tout de suite, ne pas s'automédicamenter. La santé avant tout ! Le délégué santé. »
\end{corrige}

\courssub{Éléments socioculturels : Cameroun et Chine}

\begin{savaistu}[Paludisme : regards croisés Cameroun — Chine]
\begin{itemize}
\item \textbf{Cameroun :} Zone à transmission stable et élevée. Le paludisme (\zh{疟疾 nǜèjí}) est la première cause de consultation (env. 30-40\% des consultations) et de mortalité infantile. Les espèces dominantes : \zh{恶性疟原虫} (Plasmodium falciparum). Stratégies nationales : Moustiquaires imprégnées (MILD) gratuites pour femmes enceintes/enfants <5 ans, Traitement préventif intermittent (TPI) grossesse, Pulvérisation intra-domiciliaire (PID) dans le Nord. Vocabulaire local : « Palu » (familier), « Fièvre typhoïde » souvent confondu.
\item \textbf{Chine :} \textbf{Certifiée « sans paludisme » par l'OMS en juin 2021} (après 70 ans de lutte). Dernier cas autochtone signalé en 2016 (Yunnan). Stratégie historique : « 1-3-7 » (Déclaration 1j, Enquête 3j, Réponse 7j), traitement de masse, assainissement, médicament \zh{青蒿素} (qīnghāosù, artémisine) découverte par \zh{屠呦呦} (Tú Yōuyōu, Prix Nobel Médecine 2015). Aujourd'hui : surveillance des cas importés (retour d'Afrique/Asie du Sud-Est) aux aéroports/frontières.
\item \textbf{Comparaison :} Le Cameroun est en phase de \zh{控制} (contrôle) / \zh{消除} (élimination visée 2030) ; la Chine est en phase de \zh{防止再传入} (empêcher la réintroduction). Coopération : Équipes médicales chinoises au Cameroun, don de médicaments (artémisine), formation de personnel, construction d'hôpitaux (ex: Hôpital de l'Amitié à Yaoundé).
\end{itemize}
\end{savaistu}

\begin{attention}[Les 5 erreurs qui coûtent des points au Bac — Spéciale Leçon 12]
\begin{enumerate}
\item \textbf{Confondre les caractères proches (homophones / graphies voisines).}
\begin{itemize}
\item \zh{帐} (wénzhàng, moustiquaire — clé \zh{巾} jīn, tissu) $\neq$ \zh{账} (zhàng, compte, facture — clé \zh{贝} bèi, argent/coquillage).
\item \zh{疟} (nǜè, paludisme — clé \zh{疒}) $\neq$ \zh{虐} (nǜè, maltraiter — clé \zh{虎} hǔ/tigre en bas).
\item \zh{疾} (jí, maladie — clé \zh{疒} + \zh{矢}) $\neq$ \zh{及} (jí, atteindre — pas de clé maladie).
\item \zh{呕} (ǒu, nausée — clé \zh{口} + \zh{区}) $\neq$ \zh{欧} (ōu, Europe — clé \zh{欠}).
\end{itemize}
\item \textbf{Oublier la structure passive \zh{被} (bèi) pour la victime.}
« Il a été piqué par un moustique » $\rightarrow$ \zh{他被蚊子咬了。} (Sujet = victime).
\textit{Erreur fréquente :} *\zh{蚊子咬了他} (Le moustique l'a piqué — actif, change le focus). Au Bac, si le sujet imposé est la personne, \zh{被} est obligatoire.
\item \textbf{Mal placer les compléments de TEMPS et de LIEU.}
Français : « Il va à l'hôpital demain. » $\rightarrow$ Chinois : \zh{他明天去医院。} (Temps \zh{明天} AVANT verbe \zh{去}).
Français : « Il a pris le médicament à la maison. » $\rightarrow$ \zh{他在家吃了药。} (Lieu \zh{在家} AVANT verbe \zh{吃}).
\textit{Règle :} Sujet + Temps + Lieu + Verbe + Complément.
\item \textbf{Calquer les collocations françaises (Chinglish).}
\begin{tabularx}{\linewidth}{|X|X|X|}
\hline
\rowcolor{poporangeL} \textbf{Français} & \textbf{Faux chinois (Calque)} & \textbf{Vrai chinois (Collocation)} \\ \hline
Prendre un médicament & *\zh{拿药} / *\zh{取药} & \zh{吃药} (chī yào) \\ \hline
Avoir de la fièvre & *\zh{有发烧} / *\zh{得发烧} & \zh{发烧} (fāshāo) / \zh{发高烧} \\ \hline
Consulter un médecin & *\zh{问医生} / *\zh{找医生} & \zh{看病} / \zh{看医生} (kàn bìng / kàn yīshēng) \\ \hline
Garder la maison propre & *\zh{让家干净} / *\zh{使家干净} & \zh{保持家里干净} (bǎochí jiālǐ gānjìng) \\ \hline
Installer une moustiquaire & *\zh{安装蚊帐} / *\zh{放蚊帐} & \zh{挂蚊帐} (guà wénzhàng) \\ \hline
\end{tabularx}
\item \textbf{Négliger les tons qui changent le sens (Pinyin sans tons = 0 point).}
\begin{itemize}
\item \zh{药} yào (4\up{e}, médicament) vs \zh{要} yào (4\up{e}, vouloir/falloir) — même ton, sens radicalement différents, seul le contexte/caractère distingue.
\item \zh{疟} nǜè (4\up{e}) vs \zh{虐} nǜè (4\up{e}) — même son, clés différentes.
\item \zh{疾} jí (2\up{e}, maladie) vs \zh{鸡} jī (1\up{e}, poulet) — tons différents !
\item \zh{吐} tǔ (3\up{e}, vomir/cracher) vs \zh{土} tǔ (3\up{e}, terre) — même son/ton, clés différentes (\zh{口} vs \zh{土}).
\end{itemize}
\textbf{Astuce :} Toujours écrire le pinyin \textbf{AVEC les marques de tons} (ā á ǎ à) sur les voyelles, jamais de chiffres (ni3 hao3).
\end{enumerate}
\end{attention}

\newpage

\coursec{Exercices}

\courssub{Je vérifie mes connaissances}

\begin{exercice}[QCM : le vocabulaire du paludisme]
Entoure la bonne réponse.
\begin{enumerate}
\item « Paludisme » se dit en chinois : \textbf{a)} \zh{感冒} \quad \textbf{b)} \zh{疟疾} \quad \textbf{c)} \zh{头疼}
\item \zh{蚊子} (wénzi) signifie : \textbf{a)} la mouche \quad \textbf{b)} le moustique \quad \textbf{c)} l'abeille
\item \zh{发烧} (fāshāo) signifie : \textbf{a)} avoir froid \quad \textbf{b)} avoir de la fièvre \quad \textbf{c)} avoir mal à la tête
\item \zh{蚊帐} (wénzhàng) signifie : \textbf{a)} la moustiquaire \quad \textbf{b)} le lit \quad \textbf{c)} la fenêtre
\item \zh{预防} (yùfáng) signifie : \textbf{a)} soigner \quad \textbf{b)} prévenir \quad \textbf{c)} transmettre
\end{enumerate}
\end{exercice}

\begin{exercice}[Vrai ou faux]
Réponds par « vrai » ou « faux » et justifie ta réponse par une phrase en français.
\begin{enumerate}
\item \zh{疟疾} (nüèji) désigne une maladie transmise par les moustiques.
\item \zh{及时} (jíshí) signifie « trop tard ».
\item \zh{医院} (yīyuàn) est le lieu où l'on va pour se faire soigner.
\item \zh{传染} (chuánrǎn) signifie « se reposer ».
\end{enumerate}
\end{exercice}

\begin{exercice}[Associe le mot à sa traduction]
Associe chaque mot chinois à sa traduction française.

\begin{center}
\begin{tabularx}{\linewidth}{|X|X|}\hline
\rowcolor{popblueL} \textbf{Mot chinois} & \textbf{Traduction} \\ \hline
1. \zh{疾病} & a. le médicament \\ \hline
2. \zh{病人} & b. la maladie \\ \hline
3. \zh{药} & c. le malade \\ \hline
4. \zh{医生} & d. la santé \\ \hline
5. \zh{健康} & e. le médecin \\ \hline
\end{tabularx}
\end{center}
\end{exercice}

\begin{exercice}[Associe le pinyin aux caractères]
Associe chaque caractère à son pinyin correct, puis recopie chaque paire sur ton cahier.

\begin{center}
\begin{tabularx}{\linewidth}{|X|X|}\hline
\rowcolor{popblueL} \textbf{Caractère} & \textbf{Pinyin} \\ \hline
1. \zh{疟} & a. shāo \\ \hline
2. \zh{疾} & b. zhàng \\ \hline
3. \zh{蚊} & c. nüè \\ \hline
4. \zh{帐} & d. jí \\ \hline
5. \zh{烧} & e. wén \\ \hline
\end{tabularx}
\end{center}
\end{exercice}

\courssub{Je m'entraîne}

\begin{exercice}[Traduction : du français vers le chinois]
Traduis les phrases suivantes en chinois (caractères et pinyin).
\begin{enumerate}
\item Le moustique transmet la maladie.
\item Mon petit frère a de la fièvre, il doit aller à l'hôpital.
\item Nous utilisons des moustiquaires pour prévenir le paludisme.
\end{enumerate}
\end{exercice}

\begin{exercice}[Texte à trous]
Complète le texte suivant avec les mots de la liste : \zh{疟疾、蚊子、发烧、蚊帐、医院}。

\begin{textebox}[Texte à compléter]
昨天晚上我弟弟\zh{＿＿}了，烧\zh{得}很高。妈妈说，可能是\zh{＿＿}，因为最近\zh{＿＿}很多。今天早上我们去了\zh{＿＿}。医生说要好好休息，睡觉的时候一定要用\zh{＿＿}。\\
\emph{Zuótiān wǎnshang wǒ dìdi \_\_ le, shāo de hěn gāo. Māma shuō, kěnéng shì \_\_, yīnwèi zuìjìn \_\_ hěn duō. Jīntiān zǎoshang wǒmen qù le \_\_. Yīshēng shuō yào hǎohao xiūxi, shuìjiào de shíhou yídìng yào yòng \_\_.}\\
« Hier soir, mon petit frère \_\_ ; sa fièvre était très élevée. Maman a dit que c'était peut-être \_\_, car en ce moment il y a beaucoup de \_\_. Ce matin, nous sommes allés à \_\_. Le médecin a dit qu'il fallait bien se reposer et toujours utiliser \_\_ pour dormir. »
\end{textebox}
\end{exercice}

\begin{exercice}[Remets les mots dans l'ordre]
Remets les mots dans l'ordre pour former des phrases correctes, puis traduis-les en français.
\begin{enumerate}
\item \zh{疟疾 / 是 / 一种 / 常见的 / 疾病}
\item \zh{蚊子 / 把 / 疾病 / 传染 / 给 / 人们}
\item \zh{我们 / 应该 / 及时 / 去 / 医院 / 看病}
\end{enumerate}
\end{exercice}

\begin{exercice}[Classe par radical]
Recopie les caractères suivants et classe-les selon leur clé (radical) : \zh{疟、疾、病、蚊、帐、烧、疼、疗}。Pour chaque caractère, indique la clé (\zh{疒} « maladie », \zh{虫} « insecte », \zh{巾} « tissu », \zh{火} « feu ») et le nombre total de traits.
\end{exercice}

\courssub{Je me prépare au Bac}

\begin{exercice}[Compréhension écrite type Bac]
Lis le texte suivant, puis réponds aux questions \textbf{en chinois}, par des phrases complètes.

\begin{textebox}[Le paludisme au Cameroun]
在喀麦隆，疟疾是一种很常见的疾病。每年雨季的时候，蚊子特别多，很多孩子会发烧、头疼。医生说，疟疾是蚊子传染给人的，所以我们要好好预防。第一，睡觉的时候要用蚊帐；第二，要常常打扫房子，不让蚊子进来；第三，生病的时候要马上去医院，及时吃药。雅温得的一所中学还组织学生打扫校园，教他们怎么保护自己和家人。现在，越来越多的人知道怎么预防疟疾了。\\
\emph{Zài Kāmàilóng, nüèji shì yì zhǒng hěn chángjiàn de jíbìng. Měi nián yǔjì de shíhou, wénzi tèbié duō, hěn duō háizi huì fāshāo, tóuténg. Yīshēng shuō, nüèji shì wénzi chuánrǎn gěi rén de, suǒyǐ wǒmen yào hǎohao yùfáng. Dì-yī, shuìjiào de shíhou yào yòng wénzhàng ; dì-èr, yào chángcháng dǎsǎo fángzi, bú ràng wénzi jìnlai ; dì-sān, shēngbìng de shíhou yào mǎshàng qù yīyuàn, jíshí chī yào. Yǎwēndé de yì suǒ zhōngxué hái zǔzhī xuésheng dǎsǎo xiàoyuán, jiāo tāmen zěnme bǎohù zìjǐ hé jiārén. Xiànzài, yuè lái yuè duō de rén zhīdào zěnme yùfáng nüèji le.}\\
« Au Cameroun, le paludisme est une maladie très courante. Chaque année, pendant la saison des pluies, les moustiques sont particulièrement nombreux et beaucoup d'enfants ont de la fièvre et mal à la tête. Les médecins disent que le paludisme est transmis aux hommes par les moustiques ; nous devons donc bien le prévenir. Premièrement, il faut utiliser une moustiquaire pour dormir ; deuxièmement, il faut nettoyer souvent la maison pour empêcher les moustiques d'entrer ; troisièmement, quand on est malade, il faut aller tout de suite à l'hôpital et prendre des médicaments à temps. Un collège de Yaoundé organise même les élèves pour nettoyer la cour de l'école et leur apprend à se protéger, eux et leur famille. Aujourd'hui, de plus en plus de gens savent comment prévenir le paludisme. »
\end{textebox}

Questions :
\begin{enumerate}
\item \zh{疟疾是怎么传染给人的？}
\item \zh{雨季的时候，很多孩子怎么样？}
\item \zh{医生提出了哪三个预防疟疾的办法？}
\item \zh{雅温得的中学组织学生做什么？}
\item \zh{« 及时吃药 » 是什么意思？用自己的话解释。}
\end{enumerate}
\end{exercice}

\begin{exercice}[Thème et version]
\textbf{A. Thème.} Traduis en chinois (caractères et pinyin) :

« Le paludisme est une maladie courante. Les moustiques transmettent la maladie aux hommes. Pour prévenir le paludisme, nous devons utiliser des moustiquaires. »

\textbf{B. Version.} Traduis en français le paragraphe suivant :

\begin{textebox}[À traduire]
我妹妹昨天发烧了，烧得很高。妈妈马上带她去医院。医生说，她得的是疟疾，要按时吃药，好好休息。医生还说，睡觉的时候要用蚊帐，这样才能预防疟疾。\\
\emph{Wǒ mèimei zuótiān fāshāo le, shāo de hěn gāo. Māma mǎshàng dài tā qù yīyuàn. Yīshēng shuō, tā dé de shì nüèji, yào ànshí chī yào, hǎohao xiūxi. Yīshēng hái shuō, shuìjiào de shíhou yào yòng wénzhàng, zhèyàng cái néng yùfáng nüèji.}
\end{textebox}
\end{exercice}

\begin{exercice}[Expression écrite guidée et tracé]
\textbf{A. Expression écrite.} Ton ami chinois \zh{小明} (Xiǎo Míng) est malade. Rédige-lui un message en chinois de \textbf{5 à 8 phrases} (environ 60 à 90 caractères) pour :
\begin{itemize}
\item lui demander comment il va ;
\item lui conseiller d'aller à l'hôpital et de prendre ses médicaments à temps ;
\item lui rappeler d'utiliser une moustiquaire pour se protéger des moustiques ;
\item lui souhaiter un bon rétablissement.
\end{itemize}
Utilise au moins trois mots du glossaire de la leçon (\zh{发烧、及时、预防、蚊帐、医院}…).

\textbf{B. Tracé des caractères.} Écris les caractères \zh{疟、疾、蚊、帐、烧} sur ton cahier en respectant l'ordre des traits (de haut en bas, de gauche à droite, l'extérieur avant l'intérieur). Pour chaque caractère, indique sa clé et son pinyin.
\end{exercice}

\newpage

\coursec{Corrigés des exercices}

\begin{corrige}[Exercice 1 — QCM : le vocabulaire du paludisme]
\begin{enumerate}
\item \textbf{b)} \zh{疟疾} (nüèjí) = le paludisme. \zh{感冒} (gǎnmào) = le rhume ; \zh{头疼} (tóuténg) = avoir mal à la tête.
\item \textbf{b)} \zh{蚊子} (wénzi) = le moustique. La clé \zh{虫} (chóng, « insecte ») présente dans les deux caractères est un indice.
\item \textbf{b)} \zh{发烧} (fāshāo) = avoir de la fièvre. Littéralement « émettre de la chaleur » (\zh{发} + \zh{烧}, clé du feu \zh{火}).
\item \textbf{a)} \zh{蚊帐} (wénzhàng) = la moustiquaire. Littéralement « rideau (\zh{帐}) contre les moustiques (\zh{蚊}) ».
\item \textbf{b)} \zh{预防} (yùfáng) = prévenir. « Soigner » se dit \zh{治疗} (zhìliáo) ; « transmettre » se dit \zh{传染} (chuánrǎn).
\end{enumerate}
\textbf{Point de vigilance :} ne pas confondre \zh{疟疾} (paludisme) et \zh{感冒} (rhume) : les deux contiennent la clé de la maladie \zh{疒}, mais ce sont des maladies différentes.
\end{corrige}

\begin{corrige}[Exercice 2 — Vrai ou faux]
\begin{enumerate}
\item \textbf{Vrai.} \zh{疟疾} (nüèjí) désigne le paludisme, une maladie transmise à l'homme par la piqûre du moustique (anophèle).
\item \textbf{Faux.} \zh{及时} (jíshí) signifie « à temps, sans retard ». « Trop tard » se dit \zh{太晚} (tài wǎn) ou \zh{来不及} (láibují).
\item \textbf{Vrai.} \zh{医院} (yīyuàn) = l'hôpital, le lieu où l'on va se faire soigner : \zh{去医院看病} (qù yīyuàn kànbìng) = « aller à l'hôpital consulter ».
\item \textbf{Faux.} \zh{传染} (chuánrǎn) signifie « transmettre, contaminer ». « Se reposer » se dit \zh{休息} (xiūxi).
\end{enumerate}
\end{corrige}

\begin{corrige}[Exercice 3 — Associe le mot à sa traduction]
\begin{center}
\begin{tabularx}{\linewidth}{|X|X|X|}\hline
\rowcolor{popblueL} \textbf{Mot chinois} & \textbf{Pinyin} & \textbf{Traduction} \\ \hline
1. \zh{疾病} & jíbìng & \textbf{b.} la maladie \\ \hline
2. \zh{病人} & bìngrén & \textbf{c.} le malade \\ \hline
3. \zh{药} & yào & \textbf{a.} le médicament \\ \hline
4. \zh{医生} & yīshēng & \textbf{e.} le médecin \\ \hline
5. \zh{健康} & jiànkāng & \textbf{d.} la santé \\ \hline
\end{tabularx}
\end{center}
\textbf{Astuce :} \zh{病人} se décompose en \zh{病} (maladie) + \zh{人} (personne) = « personne malade » ; \zh{医生} en \zh{医} (médecine) + \zh{生} = « celui qui pratique la médecine ».
\end{corrige}

\begin{corrige}[Exercice 4 — Associe le pinyin aux caractères]
\begin{center}
\begin{tabularx}{\linewidth}{|X|X|X|}\hline
\rowcolor{popblueL} \textbf{Caractère} & \textbf{Pinyin} & \textbf{Mot de la leçon} \\ \hline
1. \zh{疟} & \textbf{c.} nüè & \zh{疟疾} nüèjí \\ \hline
2. \zh{疾} & \textbf{d.} jí & \zh{疾病} jíbìng \\ \hline
3. \zh{蚊} & \textbf{e.} wén & \zh{蚊子} wénzi \\ \hline
4. \zh{帐} & \textbf{b.} zhàng & \zh{蚊帐} wénzhàng \\ \hline
5. \zh{烧} & \textbf{a.} shāo & \zh{发烧} fāshāo \\ \hline
\end{tabularx}
\end{center}
\textbf{Point de vigilance :} \zh{疟} se prononce nüè (4\textsuperscript{e} ton, ü après n) ; ne pas confondre \zh{帐} (zhàng, moustiquaire, clé \zh{巾} « tissu ») avec \zh{账} (zhàng, compte, clé \zh{贝} « argent ») — mêmes sons, sens différents.
\end{corrige}

\begin{corrige}[Exercice 5 — Traduction : du français vers le chinois]
\begin{enumerate}
\item \zh{蚊子传染疾病。}\\
\emph{Wénzi chuánrǎn jíbìng.}\\
« Le moustique transmet la maladie. »\\
Structure : Sujet + verbe + complément. On peut aussi dire \zh{蚊子把疾病传染给人。} (construction avec \cle{把}).
\item \zh{我弟弟发烧了，他要去医院。}\\
\emph{Wǒ dìdi fāshāo le, tā yào qù yīyuàn.}\\
« Mon petit frère a de la fièvre, il doit aller à l'hôpital. »\\
\zh{了} marque ici le changement d'état (il est tombé malade) ; \zh{要} exprime l'obligation (« devoir »).
\item \zh{我们用蚊帐预防疟疾。}\\
\emph{Wǒmen yòng wénzhàng yùfáng nüèjí.}\\
« Nous utilisons des moustiquaires pour prévenir le paludisme. »\\
\zh{用} (yòng) = « utiliser, au moyen de » ; en chinois, « pour prévenir » se rend simplement par l'enchaînement des verbes, sans mot-outil.
\end{enumerate}
\textbf{Point de vigilance :} pas d'article en chinois (ni « le », ni « des ») ; le classificateur n'est pas obligatoire ici car on parle en général.
\end{corrige}

\begin{corrige}[Exercice 6 — Texte à trous]
Texte complété :

\zh{昨天晚上我弟弟发烧了，烧得很高。妈妈说，可能是疟疾，因为最近蚊子很多。今天早上我们去了医院。医生说要好好休息，睡觉的时候一定要用蚊帐。}

\emph{Zuótiān wǎnshang wǒ dìdi fāshāo le, shāo de hěn gāo. Māma shuō, kěnéng shì nüèjí, yīnwèi zuìjìn wénzi hěn duō. Jīntiān zǎoshang wǒmen qù le yīyuàn. Yīshēng shuō yào hǎohao xiūxi, shuìjiào de shíhou yídìng yào yòng wénzhàng.}

Ordre des réponses : \textbf{1.} \zh{发烧} \quad \textbf{2.} \zh{疟疾} \quad \textbf{3.} \zh{蚊子} \quad \textbf{4.} \zh{医院} \quad \textbf{5.} \zh{蚊帐}.

\textbf{Justification :} « 烧得很高 » (la fièvre est très élevée) exige \zh{发烧} au premier trou ; la maman formule une hypothèse (\zh{可能是}) : \zh{疟疾} ; la cause (\zh{因为}) : beaucoup de \zh{蚊子} ; on va se faire soigner : \zh{医院} ; pour dormir, on utilise la \zh{蚊帐}.
\end{corrige}

\begin{corrige}[Exercice 7 — Remets les mots dans l'ordre]
\begin{enumerate}
\item \zh{疟疾是一种常见的疾病。}\\
\emph{Nüèjí shì yì zhǒng chángjiàn de jíbìng.}\\
« Le paludisme est une maladie courante. »\\
Structure : Sujet + \zh{是} + classificateur \zh{种} + adjectif + \zh{的} + nom.
\item \zh{蚊子把疾病传染给人们。}\\
\emph{Wénzi bǎ jíbìng chuánrǎn gěi rénmen.}\\
« Les moustiques transmettent la maladie aux hommes. »\\
Construction en \cle{把} : Sujet + \zh{把} + objet + verbe + \zh{给} + bénéficiaire.
\item \zh{我们应该及时去医院看病。}\\
\emph{Wǒmen yīnggāi jíshí qù yīyuàn kànbìng.}\\
« Nous devons aller à l'hôpital consulter à temps. »\\
L'adverbe \zh{及时} se place avant le verbe ; \zh{应该} (devoir) précède tout le groupe verbal.
\end{enumerate}
\textbf{Point de vigilance :} en chinois, l'adverbe se place toujours \emph{avant} le verbe, jamais en fin de phrase comme en français.
\end{corrige}

\begin{corrige}[Exercice 8 — Classe par radical]
\begin{center}
\begin{tabularx}{\linewidth}{|X|X|X|X|}\hline
\rowcolor{popblueL} \textbf{Caractère} & \textbf{Clé} & \textbf{Traits} & \textbf{Pinyin} \\ \hline
\zh{疟} & \zh{疒} (maladie) & 9 & nüè \\ \hline
\zh{疾} & \zh{疒} (maladie) & 10 & jí \\ \hline
\zh{病} & \zh{疒} (maladie) & 10 & bìng \\ \hline
\zh{疼} & \zh{疒} (maladie) & 10 & téng \\ \hline
\zh{疗} & \zh{疒} (maladie) & 7 & liáo \\ \hline
\zh{蚊} & \zh{虫} (insecte) & 10 & wén \\ \hline
\zh{帐} & \zh{巾} (tissu) & 7 & zhàng \\ \hline
\zh{烧} & \zh{火} (feu) & 10 & shāo \\ \hline
\end{tabularx}
\end{center}
\textbf{Regroupement :} clé \zh{疒} : \zh{疟、疾、病、疼、疗} (cinq caractères liés à la maladie) ; clé \zh{虫} : \zh{蚊} ; clé \zh{巾} : \zh{帐} ; clé \zh{火} : \zh{烧}.

\textbf{Ordre des traits (rappel) :} pour les caractères à clé \zh{疒}, on trace d'abord l'enveloppe \zh{疒} (de gauche à droite, de haut en bas), puis l'intérieur ; pour \zh{蚊}, la clé \zh{虫} à gauche d'abord, puis \zh{文} à droite ; pour \zh{烧}, la clé \zh{火} à gauche d'abord, puis \zh{尧} à droite.
\end{corrige}

\begin{corrige}[Exercice 9 — Compréhension écrite type Bac]
\textbf{Réponses modèles (en chinois, phrases complètes) :}
\begin{enumerate}
\item \zh{疟疾是蚊子传染给人的。}\\
\emph{Nüèjí shì wénzi chuánrǎn gěi rén de.}\\
« Le paludisme est transmis aux hommes par les moustiques. »
\item \zh{雨季的时候，很多孩子会发烧、头疼。}\\
\emph{Yǔjì de shíhou, hěn duō háizi huì fāshāo, tóuténg.}\\
« Pendant la saison des pluies, beaucoup d'enfants ont de la fièvre et mal à la tête. »
\item \zh{第一，睡觉的时候要用蚊帐；第二，要常常打扫房子，不让蚊子进来；第三，生病的时候要马上去医院，及时吃药。}\\
\emph{Dì-yī, shuìjiào de shíhou yào yòng wénzhàng ; dì-èr, yào chángcháng dǎsǎo fángzi, bú ràng wénzi jìnlái ; dì-sān, shēngbìng de shíhou yào mǎshàng qù yīyuàn, jíshí chī yào.}\\
« Premièrement, utiliser une moustiquaire pour dormir ; deuxièmement, nettoyer souvent la maison pour empêcher les moustiques d'entrer ; troisièmement, aller tout de suite à l'hôpital quand on est malade et prendre des médicaments à temps. »
\item \zh{雅温得的中学组织学生打扫校园，教他们怎么保护自己和家人。}\\
\emph{Yǎwēndé de zhōngxué zǔzhī xuésheng dǎsǎo xiàoyuán, jiāo tāmen zěnme bǎohù zìjǐ hé jiārén.}\\
« Le collège de Yaoundé organise les élèves pour nettoyer la cour et leur apprend à se protéger, eux et leur famille. »
\item \zh{« 及时吃药 » 的意思是：生病以后要马上吃药，不能等。}\\
\emph{« Jíshí chī yào » de yìsi shì : shēngbìng yǐhòu yào mǎshàng chī yào, bù néng děng.}\\
« "Prendre ses médicaments à temps" signifie : quand on est malade, il faut prendre ses médicaments tout de suite, sans attendre. »
\end{enumerate}
\textbf{Barème indicatif (sur 10 points) :} 2 points par question (1 point pour l'exactitude du sens, 1 point pour la correction de la phrase : ordre des mots, tons, caractères).

\textbf{Points de vigilance :} répondre par une phrase complète (sujet + verbe), pas par un mot isolé ; reprendre les mots du texte est autorisé et conseillé ; attention à \zh{传染给} (transmettre à) suivi directement du bénéficiaire ; pour la question 5, « expliquer avec ses propres mots » exige une reformulation, pas une recopie.
\end{corrige}

\begin{corrige}[Exercice 10 — Thème et version]
\textbf{A. Thème — proposition de traduction :}

\zh{疟疾是一种常见的疾病。蚊子把疾病传染给人。为了预防疟疾，我们要用蚊帐。}

\emph{Nüèjí shì yì zhǒng chángjiàn de jíbìng. Wénzi bǎ jíbìng chuánrǎn gěi rén. Wèile yùfáng nüèjí, wǒmen yào yòng wénzhàng.}

\textbf{Justifications :} « une maladie courante » = \zh{一种常见的疾病} (classificateur \zh{种} pour les sortes) ; « transmettent aux hommes » = construction \cle{把} + \zh{传染给人} ; « pour prévenir » = \zh{为了} + verbe, placé en tête de phrase ; « devons » = \zh{要} ou \zh{应该}.

\textbf{B. Version — traduction :}

« Ma petite sœur a eu de la fièvre hier, une fièvre très élevée. Maman l'a tout de suite emmenée à l'hôpital. Le médecin a dit qu'elle avait le paludisme, qu'elle devait prendre ses médicaments aux heures prévues et bien se reposer. Le médecin a dit aussi qu'il fallait utiliser une moustiquaire pour dormir : c'est seulement ainsi qu'on peut prévenir le paludisme. »

\textbf{Points de vigilance :} \zh{烧得很高} = « la fièvre est très élevée » (complément de degré avec \zh{得}) ; \zh{按时} (ànshí) = « aux heures prévues, ponctuellement » (à ne pas confondre avec \zh{及时} « à temps ») ; \zh{这样……才……} = « c'est seulement ainsi que… » ; \zh{得的是疟疾} = « ce qu'elle a attrapé, c'est le paludisme » (construction \cle{是……的}).

\textbf{Barème indicatif (sur 10) :} thème 5 points (1 point par segment correct : \zh{疟疾是一种常见的疾病} / \zh{蚊子把疾病传染给人} / \zh{为了预防疟疾} / \zh{我们要用蚊帐} + 1 point pour l'ensemble : tons, caractères) ; version 5 points (fidélité du sens 3 points, qualité du français 2 points).
\end{corrige}

\begin{corrige}[Exercice 11 — Expression écrite guidée et tracé]
\textbf{A. Barème indicatif (sur 10) :} respect des quatre consignes (4 points) ; emploi d'au moins trois mots du glossaire (2 points) ; correction grammaticale et ordre des mots (2 points) ; caractères et tons corrects (2 points).

\textbf{Points de vigilance :} \zh{别忘了} = « n'oublie pas de » (impératif négatif avec \zh{别}) ; \zh{祝你早日康复} est la formule idiomatique de souhait de rétablissement ; éviter de traduire mot à mot « bon rétablissement » par \zh{好恢复}.

\textbf{B. Tracé des caractères :}
\begin{center}
\begin{tabularx}{\linewidth}{|X|X|X|X|}\hline
\rowcolor{popblueL} \textbf{Caractère} & \textbf{Clé} & \textbf{Pinyin} & \textbf{Ordre des traits} \\ \hline
\zh{疟} & \zh{疒} & nüè & enveloppe \zh{疒} d'abord, puis l'intérieur \zh{丄} + \zh{彐} \\ \hline
\zh{疾} & \zh{疒} & jí & enveloppe \zh{疒}, puis \zh{矢} à l'intérieur \\ \hline
\zh{蚊} & \zh{虫} & wén & \zh{虫} à gauche d'abord, puis \zh{文} à droite \\ \hline
\zh{帐} & \zh{巾} & zhàng & \zh{巾} à gauche d'abord, puis \zh{长} à droite \\ \hline
\zh{烧} & \zh{火} & shāo & \zh{火} à gauche d'abord, puis \zh{尧} à droite \\ \hline
\end{tabularx}
\end{center}
\textbf{Règles générales appliquées :} de haut en bas, de gauche à droite ; l'extérieur (enveloppe \zh{疒}) avant l'intérieur ; la partie gauche avant la partie droite pour les caractères composés horizontalement.
\end{corrige}

\begin{textebox}[Message modèle — Exercice 11]
小明，你好吗？听说你生病了，我很担心。你发烧吗？你应该马上去医院看病，按时吃药，好好休息。睡觉的时候别忘了用蚊帐，这样可以预防疟疾，不让蚊子传染你。祝你早日康复！\\
\emph{Xiǎo Míng, nǐ hǎo ma ? Tīngshuō nǐ shēngbìng le, wǒ hěn dānxīn. Nǐ fāshāo ma ? Nǐ yīnggāi mǎshàng qù yīyuàn kànbìng, ànshí chī yào, hǎohao xiūxi. Shuìjiào de shíhou bié wàng le yòng wénzhàng, zhèyàng kěyǐ yùfáng nüèjí, bú ràng wénzi chuánrǎn nǐ. Zhù nǐ zǎorì kāngfù !}\\
« Xiǎo Míng, comment vas-tu ? J'ai entendu dire que tu étais malade, je suis très inquiet. As-tu de la fièvre ? Tu devrais aller tout de suite à l'hôpital consulter, prendre tes médicaments aux heures prévues et bien te reposer. Quand tu dors, n'oublie pas d'utiliser une moustiquaire : ainsi tu peux prévenir le paludisme et empêcher les moustiques de te contaminer. Je te souhaite un prompt rétablissement ! »
\end{textebox}

\newpage

\coursec{Fiche bilan}

\begin{popfigure}
\begin{tikzpicture}[mindmap, grow cyclic, every node/.style={concept, font=\small},
  concept color=popblue, text=white,
  level 1/.append style={level distance=3.2cm, sibling angle=60},
  level 2/.append style={level distance=2.0cm, sibling angle=45, font=\scriptsize}]
\node[concept, font=\small\bfseries, align=center]{Leçon 12\\ Le paludisme\\ \zh{疟疾}}
  [counterclockwise from=90]
  child[concept color=poporange] { node[align=center]{Lexique\\ maladie} 
    child { node {\zh{发烧} fāshāo} }
    child { node {\zh{头疼} tóuténg} }
  }
  child[concept color=popgreen] { node[align=center]{Lexique\\ prévention}
    child { node {\zh{蚊帐} wénzhàng} }
    child { node {\zh{蚊子} wénzi} }
  }
  child[concept color=poppurple] { node[align=center]{Lexique\\ soins}
    child { node {\zh{医院} yīyuàn} }
    child { node {\zh{吃药} chī yào} }
  }
  child[concept color=popgold] { node {Radicaux}
    child { node {疒 maladie} }
    child { node {虫 insecte} }
  }
  child[concept color=poppink] { node {Structures}
    child { node {\zh{因为……所以……}} }
    child { node {\zh{应该} + verbe} }
  }
  child[concept color=popblue!70] { node[align=center]{Compréhension\\ du texte}
    child { node {repérer qui / quoi / où} }
  };
\end{tikzpicture}
\legende{Carte mentale de la leçon 12 : le paludisme (\zh{疟疾}).}
\end{popfigure}

\begin{bilan}
\textbf{1. Lexique clé à connaître par cœur}
\begin{itemize}
  \item \zh{疟疾} nüèji : le paludisme ; \zh{生病} shēngbìng : tomber malade ; \zh{病} bìng : maladie.
  \item \zh{发烧} fāshāo : avoir de la fièvre ; \zh{头疼} tóuténg : avoir mal à la tête ; \zh{咳嗽} késou : tousser.
  \item \zh{蚊子} wénzi : moustique ; \zh{蚊帐} wénzhàng : moustiquaire ; \zh{打蚊子} dǎ wénzi : tuer les moustiques.
  \item \zh{医院} yīyuàn : hôpital ; \zh{医生} yīshēng : médecin ; \zh{药} yào : médicament ; \zh{看病} kànbìng : consulter (un médecin).
  \item \zh{预防} yùfáng : prévenir ; \zh{健康} jiànkāng : santé, en bonne santé ; \zh{注意} zhùyì : faire attention à.
\end{itemize}

\textbf{2. Structures et règles de grammaire}
\begin{center}
\begin{tabularx}{\linewidth}{|X|X|X|}
\hline
\rowcolor{popblueL} Structure & Exemple & Traduction \\
\hline
因为 + cause，所以 + résultat & \zh{因为蚊子多，所以容易得疟疾。} Yīnwèi wénzi duō, suǒyǐ róngyì dé nüèji. & Comme il y a beaucoup de moustiques, on attrape facilement le paludisme. \\
\hline
应该 + verbe (conseil) & \zh{我们应该用蚊帐。} Wǒmen yīnggāi yòng wénzhàng. & Nous devrions utiliser des moustiquaires. \\
\hline
Sujet + 要 + verbe (nécessité) & \zh{发烧了要去看医生。} Fāshāo le yào qù kàn yīshēng. & Quand on a de la fièvre, il faut aller voir le médecin. \\
\hline
Verb + 得 + adjectif (manière) & \zh{他病得很重。} Tā bìng de hěn zhòng. & Il est gravement malade. \\
\hline
\end{tabularx}
\end{center}

\textbf{3. Radicaux et écriture}
\begin{itemize}
  \item Clé 疒 (nè, « maladie ») : \zh{病}, \zh{疼}, \zh{疟} — on écrit d'abord le point et l'horizontale, puis le trait tombant à gauche, enfin l'intérieur.
  \item Clé 虫 (chóng, « insecte ») : \zh{蚊}, \zh{疟} (sens lié) — ordre : gauche avant droite, haut avant bas.
  \item Clé 艹 (cǎo, « herbe ») : \zh{药} — le radical « herbe » s'écrit en premier, en haut.
\end{itemize}

\textbf{4. Méthode express : comprendre un texte de santé}
\begin{enumerate}
  \item Première lecture : repérer le thème (ici : une maladie) et les mots connus.
  \item Souligner le vocabulaire du corps et des symptômes (\zh{发烧}, \zh{头疼}).
  \item Repérer les connecteurs \zh{因为……所以……} pour identifier cause et conséquence.
  \item Répondre aux questions en reprenant les mots du texte, phrases courtes.
\end{enumerate}
\end{bilan}

\begin{aretenir}[Checklist avant le Bac]
\begin{itemize}
  \item $\square$ Je sais lire et écrire \zh{疟疾}, \zh{发烧}, \zh{头疼}, \zh{蚊子}, \zh{蚊帐}, \zh{医院}, \zh{药}, \zh{预防} avec le bon pinyin et les bons tons.
  \item $\square$ Je connais les radicaux 疒 (maladie), 虫 (insecte) et 艹 (herbe) et je sais les reconnaître dans un caractère inconnu.
  \item $\square$ Je respecte l'ordre des traits : haut avant bas, gauche avant droite, horizontal avant vertical.
  \item $\square$ J'emploie correctement \zh{因为……所以……} sans oublier la virgule entre les deux propositions.
  \item $\square$ Je sais donner un conseil avec \zh{应该} + verbe et une nécessité avec \zh{要} + verbe.
  \item $\square$ Je distingue \zh{看病} (consulter) et \zh{看医生} (voir le médecin).
  \item $\square$ Je place correctement le complément de manière avec \zh{得} : verbe + 得 + adjectif.
  \item $\square$ Je sais décrire des symptômes : \zh{我发烧、头疼。}
  \item $\square$ Je sais répondre à une question de compréhension avec une phrase complète et courte.
  \item $\square$ Je peux expliquer en trois phrases comment prévenir le paludisme au Cameroun (moustiquaire, consultation, propreté).
\end{itemize}
\end{aretenir}

\findecours

\end{document}
